%% GENERATED by generate_papers.py -- edit the bank_*.py files, not this file.
\documentclass[11pt,a4paper]{article}
\usepackage[a4paper, top=20mm, bottom=20mm, left=18mm, right=18mm, headheight=14pt]{geometry}
\usepackage[T1]{fontenc}
\usepackage{lmodern}
\usepackage{amsmath, amssymb, mathtools}
\usepackage{braket}
\usepackage{array, tabularx, booktabs, multicol}
\usepackage{enumitem}
\usepackage[dvipsnames]{xcolor}
\usepackage[most]{tcolorbox}
\usepackage{needspace}
\usepackage{fancyhdr, lastpage}
\usepackage{pifont}
\usepackage{tikz}
\usepackage{graphicx}
\usepackage{eso-pic}
% SILICO FELLER watermark: official logo, centred, ~8 % opacity, behind all content, on every page
\AddToShipoutPictureBG{\AtPageCenter{\makebox(0,0){\tikz\node[opacity=0.08, inner sep=0pt]{\includegraphics[width=0.72\paperwidth]{../../../assets/branding/silico_feller_logo.png}};}}}
\usepackage[hidelinks]{hyperref}
\newcolumntype{L}{>{\raggedright\arraybackslash}X}
\definecolor{brand}{HTML}{1F4E79}
\definecolor{okgreen}{HTML}{1B7F3B}
\setlength{\parindent}{0pt}
\setlength{\parskip}{3pt}
\newcommand{\optlabel}[1]{\textbf{(#1)}}
\newcommand{\correctlabel}[1]{\textcolor{okgreen}{\textbf{(#1)}\,\ding{51}}}
\newenvironment{qitem}[1]{\begin{list}{}{\setlength{\leftmargin}{3.3em}\setlength{\labelwidth}{3.1em}%
  \setlength{\labelsep}{0.2em}\setlength{\topsep}{3pt}\setlength{\parsep}{3pt}}\item[\textbf{Q.#1}\hfill]}{\end{list}}
\newcommand{\qsep}{\par\vspace{2pt}{\color{black!15}\hrule}\vspace{4pt}}
\newcommand{\qmeta}[2]{{\footnotesize\color{black!60}\textsc{#1} \;$\cdot$\; #2}\par}
\newtcolorbox{passagebox}[1]{enhanced, breakable, colback=brand!4, colframe=brand!70, fonttitle=\bfseries,
  title={#1}, left=6pt, right=6pt, top=4pt, bottom=4pt}
\newtcolorbox{sectionbanner}{enhanced, colback=brand, colframe=brand, coltext=white, fontupper=\large\bfseries,
  left=8pt, right=8pt, top=6pt, bottom=6pt, sharp corners}
\newtcolorbox{infobox}[1]{enhanced, breakable, colback=black!3, colframe=black!45, fonttitle=\bfseries,
  title={#1}, left=6pt, right=6pt, top=4pt, bottom=4pt}
\newtcolorbox{solbox}[1]{enhanced, breakable, colback=okgreen!3, colframe=okgreen!60!black, fonttitle=\bfseries,
  title={Answer: (#1) \quad Detailed solution}, left=6pt, right=6pt, top=3pt, bottom=4pt, before skip=4pt, after skip=6pt}
\newcommand{\QFTop}{\mathrm{QFT}}
\allowdisplaybreaks

\begin{document}
\pagestyle{fancy}\fancyhf{}
\fancyhead[L]{\small\color{black!70}Quantum Fourier Transform --- Assessment}
\fancyhead[R]{\small\color{black!70}\textbf{Answer Book}}
\fancyfoot[C]{\small Page \thepage\ of \pageref*{LastPage}}
\renewcommand{\headrulewidth}{0.4pt}

\thispagestyle{empty}
\begin{center}
{\color{brand}\rule{\linewidth}{1.2pt}}\\[6pt]
{\Huge\bfseries\color{brand} Quantum Fourier Transform}\\[4pt]
{\Large Assessment of the QFT algorithm}\\[4pt]
{\large\bfseries ANSWER BOOK}\\[2pt]
{\color{brand}\rule{\linewidth}{1.2pt}}
\end{center}
This book accompanies the separate \emph{Question Paper} (80 questions, 3 hours, 240 marks). It contains
\begin{enumerate}[leftmargin=1.4em, itemsep=0pt, topsep=2pt]
\item the \textbf{answer key} with the topic tested by every question (page~\pageref{key}), and
\item the \textbf{Answer Bank}: every question restated with the correct option marked \textcolor{okgreen}{\ding{51}} and a
      complete solution with derivations (from page~\pageref{bank}).
\end{enumerate}
\begin{infobox}{Marking scheme (every question)}
Each question has \textbf{four} options (A), (B), (C) and (D). \textbf{ONLY ONE} of these four options is correct.\\[2pt]
\begin{tabular}{@{}ll@{}}
\emph{Full marks}: & $+3$ if ONLY the correct option is chosen;\\
\emph{Zero marks}: & $\phantom{+}0$ if none of the options is chosen (question unanswered);\\
\emph{Negative marks}: & $-1$ in all other cases (wrong option, or more than one option chosen).
\end{tabular}
\end{infobox}

\begin{infobox}{How to evaluate}
For each section count $C$ = number of correct answers and $W$ = number of questions answered wrongly (or with more than one option marked). Section marks $= 3C - W$; unanswered questions score $0$. The total is out of $240$.
Every numerical answer was re-computed independently (NumPy, and CUDA-Q state vectors for the QFT circuits) by
\texttt{verify\_answers.py} in the source folder.
\end{infobox}
\begin{infobox}{Conventions (identical to the Question Paper)}

\begin{itemize}[leftmargin=1.2em, itemsep=1pt, topsep=2pt]
\item $n$ qubits, $N = 2^n$, $\omega = e^{2\pi i/N}$. A basis state is written $\ket{x} = \ket{x_1x_2\dots x_n}$ with
      $x = \sum_{j} x_j 2^{\,n-j}$: \textbf{qubit 1 is the most significant bit} and is written first.
\item $\displaystyle \mathrm{QFT}\ket{x} = \frac{1}{\sqrt N}\sum_{y=0}^{N-1}\omega^{xy}\ket{y}$, \quad
      $\displaystyle \mathrm{QFT}^{-1}\ket{y} = \frac{1}{\sqrt N}\sum_{x=0}^{N-1}\omega^{-xy}\ket{x}$.
      Binary fraction: $0.x_lx_{l+1}\dots x_m = \frac{x_l}{2} + \frac{x_{l+1}}{4} + \dots + \frac{x_m}{2^{m-l+1}}$.
\item $H = \frac{1}{\sqrt2}\begin{psmallmatrix}1&1\\1&-1\end{psmallmatrix}$; $X, Y, Z$ are the Pauli matrices;
      $R_k = \operatorname{diag}\big(1, e^{2\pi i/2^k}\big)$, so $R_1 = Z$, $S = R_2$, $T = R_3$;
      $\ket{\pm} = \frac{1}{\sqrt2}(\ket0 \pm \ket1)$. Two-qubit matrices use the basis order $\ket{00}, \ket{01}, \ket{10}, \ket{11}$.
\item \textbf{Textbook QFT circuit:} for $j = 1, \dots, n$: apply $H$ to qubit $j$, then for $k = 2, \dots, n-j+1$ a controlled-$R_k$
      with control qubit $j+k-1$ and target qubit $j$; finally SWAP qubits $j$ and $n+1-j$ for $j \le n/2$.
\item \textbf{Phase estimation (QPE):} $U\ket{u} = e^{2\pi i\varphi}\ket{u}$, $0 \le \varphi < 1$; $t$ counting qubits in $\ket{+}$ control
      $U^{2^{t-1}}, \dots, U^{2}, U$; the inverse QFT is applied and the outcome $y$ gives the estimate $\varphi \approx y/2^t$.
\item All measurements are in the computational basis unless stated otherwise; $\log$ means $\log_2$.
\end{itemize}
\end{infobox}
\newpage
\section*{Answer key}\label{key}
{\small Answer distribution: (A) 20, (B) 20, (C) 20, (D) 20.}

\needspace{14\baselineskip}
\subsection*{Section A --- Mathematics for the QFT}
{\small\renewcommand{\arraystretch}{1.15}
\begin{tabularx}{\linewidth}{@{}r c L @{\hspace{1em}} r c L@{}}
\toprule
\textbf{Q.} & \textbf{Ans.} & \textbf{Topic tested} & \textbf{Q.} & \textbf{Ans.} & \textbf{Topic tested}\\ \midrule
1 & \textbf{A} & Weighted sum of roots of unity & 11 & \textbf{C} & FFT vs.\ QFT cost (and its limits) \\
2 & \textbf{B} & Orthogonality of roots of unity & 12 & \textbf{A} & Cyclotomic product identity \\
3 & \textbf{A} & Binary fractions, periodicity of $e^{2\pi i\theta}$ & 13 & \textbf{D} & DFT of a vector; Parseval \\
4 & \textbf{D} & Unitarity; the $4\times4$ Fourier matrix & 14 & \textbf{D} & Binary expansion of $2^k\varphi$ \\
5 & \textbf{A} & Gauss sum, trace of $F_4$ & 15 & \textbf{C} & Algebra of 8th roots of unity \\
6 & \textbf{D} & Interference pattern $|\sum e^{ij\theta}|^2$ & 16 & \textbf{B} & Properties of $F$, $F^2$, $H\otimes H$, C$T$ \\
7 & \textbf{B} & Continued fractions, convergents & 17 & \textbf{C} & Root sums when $\gcd(k,N)>1$ \\
8 & \textbf{C} & Multiplicative order; Shor's gcd step & 18 & \textbf{C} & Determinants of unitary matrices \\
9 & \textbf{B} & Parameters of an $n$-qubit state & 19 & \textbf{A} & Parseval for the un-normalised DFT \\
10 & \textbf{B} & Walsh--Hadamard signs & 20 & \textbf{D} & Spectral leakage ($r\nmid N$) \\
\bottomrule
\end{tabularx}}

\needspace{14\baselineskip}
\subsection*{Section B --- Quantum Mechanics}
{\small\renewcommand{\arraystretch}{1.15}
\begin{tabularx}{\linewidth}{@{}r c L @{\hspace{1em}} r c L@{}}
\toprule
\textbf{Q.} & \textbf{Ans.} & \textbf{Topic tested} & \textbf{Q.} & \textbf{Ans.} & \textbf{Topic tested}\\ \midrule
21 & \textbf{B} & Born rule, partial measurement & 31 & \textbf{D} & Commutation of gates \\
22 & \textbf{A} & Global vs.\ relative phase & 32 & \textbf{B} & State-vector memory cost \\
23 & \textbf{A} & Interference with a Hadamard & 33 & \textbf{C} & Bell state in the $X$ basis \\
24 & \textbf{B} & Bloch-sphere expectation values & 34 & \textbf{B} & Mixture vs.\ superposition \\
25 & \textbf{C} & Entanglement test & 35 & \textbf{C} & Single-qubit gate actions \\
26 & \textbf{D} & Phase kickback (CNOT) & 36 & \textbf{D} & Two-qubit correlations \\
27 & \textbf{A} & Controlled-phase matrices & 37 & \textbf{C} & Distinguishability; relative phase \\
28 & \textbf{B} & Kickback of controlled-$U^{2^j}$ & 38 & \textbf{D} & Product-state criterion \\
29 & \textbf{A} & Post-measurement state & 39 & \textbf{D} & Hadamard test \\
30 & \textbf{A} & Unitarity of gates & 40 & \textbf{C} & Repetition and confidence \\
\bottomrule
\end{tabularx}}

\needspace{14\baselineskip}
\subsection*{Section C --- QFT Basics}
{\small\renewcommand{\arraystretch}{1.15}
\begin{tabularx}{\linewidth}{@{}r c L @{\hspace{1em}} r c L@{}}
\toprule
\textbf{Q.} & \textbf{Ans.} & \textbf{Topic tested} & \textbf{Q.} & \textbf{Ans.} & \textbf{Topic tested}\\ \midrule
41 & \textbf{A} & QFT matrix element ($n = 2$) & 51 & \textbf{A} & Constructing the inverse QFT \\
42 & \textbf{D} & Output distribution of $\mathrm{QFT}\ket{x}$ & 52 & \textbf{B} & Phase of a QFT amplitude \\
43 & \textbf{C} & Product formula ($n = 3$) & 53 & \textbf{B} & QFT vs.\ NumPy FFT conventions \\
44 & \textbf{C} & Controlled-rotation count & 54 & \textbf{D} & $\mathrm{QFT}\ket{x}$ is a product state \\
45 & \textbf{B} & Two-qubit gates incl.\ SWAPs & 55 & \textbf{C} & Gate count $n(n-1)/2$ \\
46 & \textbf{D} & Smallest rotation angle & 56 & \textbf{D} & Two-qubit QFT outputs \\
47 & \textbf{A} & Phase on qubit 1; role of SWAPs & 57 & \textbf{A} & Why measuring $\mathrm{QFT}\ket{x}$ reveals nothing \\
48 & \textbf{C} & $\mathrm{QFT}^2$ = negation & 58 & \textbf{C} & What the QFT does \emph{not} speed up \\
49 & \textbf{B} & QFT of the uniform state & 59 & \textbf{B} & Circuit trace: phase build-up \\
50 & \textbf{A} & QFT of $\ket{10\dots0}$ & 60 & \textbf{D} & Circuit trace: final amplitude, SWAP \\
\bottomrule
\end{tabularx}}

\needspace{14\baselineskip}
\subsection*{Section D --- Advanced QFT}
{\small\renewcommand{\arraystretch}{1.15}
\begin{tabularx}{\linewidth}{@{}r c L @{\hspace{1em}} r c L@{}}
\toprule
\textbf{Q.} & \textbf{Ans.} & \textbf{Topic tested} & \textbf{Q.} & \textbf{Ans.} & \textbf{Topic tested}\\ \midrule
61 & \textbf{D} & QFT of a periodic state ($r\mid N$) & 71 & \textbf{D} & Approximate QFT gate savings \\
62 & \textbf{B} & Phases of QFT peaks & 72 & \textbf{A} & Periodic state with $r\nmid N$ \\
63 & \textbf{D} & Shift theorem & 73 & \textbf{D} & Eigenvalues of the QFT \\
64 & \textbf{B} & QPE precision vs.\ counting qubits & 74 & \textbf{C} & Semiclassical QFT \\
65 & \textbf{B} & QPE with exact $t$-bit phase & 75 & \textbf{A} & Applications of the QFT \\
66 & \textbf{C} & The $4/\pi^2$ success bound & 76 & \textbf{C} & Period $\leftrightarrow$ peak spacing \\
67 & \textbf{D} & QPE distribution for $\varphi = 1/3$ & 77 & \textbf{A} & Approximate QFT error bound \\
68 & \textbf{A} & Qubit budget of Shor's algorithm & 78 & \textbf{A} & Non-terminating binary phases \\
69 & \textbf{B} & Order finding, $a = 7$, $N = 15$ & 79 & \textbf{C} & Shor, $N = 21$: continued fractions \\
70 & \textbf{B} & Failure mode $\gcd(k, r) > 1$ & 80 & \textbf{C} & Shor, $N = 21$: peak probability \\
\bottomrule
\end{tabularx}}

\subsection*{Quick-marking grid}
\begin{center}\renewcommand{\arraystretch}{1.3}
\begin{tabular}{|c|c|c|c|c|c|c|c|c|c|}\hline
1\,\textbf{A} & 2\,\textbf{B} & 3\,\textbf{A} & 4\,\textbf{D} & 5\,\textbf{A} & 6\,\textbf{D} & 7\,\textbf{B} & 8\,\textbf{C} & 9\,\textbf{B} & 10\,\textbf{B} \\
\hline
11\,\textbf{C} & 12\,\textbf{A} & 13\,\textbf{D} & 14\,\textbf{D} & 15\,\textbf{C} & 16\,\textbf{B} & 17\,\textbf{C} & 18\,\textbf{C} & 19\,\textbf{A} & 20\,\textbf{D} \\
\hline
21\,\textbf{B} & 22\,\textbf{A} & 23\,\textbf{A} & 24\,\textbf{B} & 25\,\textbf{C} & 26\,\textbf{D} & 27\,\textbf{A} & 28\,\textbf{B} & 29\,\textbf{A} & 30\,\textbf{A} \\
\hline
31\,\textbf{D} & 32\,\textbf{B} & 33\,\textbf{C} & 34\,\textbf{B} & 35\,\textbf{C} & 36\,\textbf{D} & 37\,\textbf{C} & 38\,\textbf{D} & 39\,\textbf{D} & 40\,\textbf{C} \\
\hline
41\,\textbf{A} & 42\,\textbf{D} & 43\,\textbf{C} & 44\,\textbf{C} & 45\,\textbf{B} & 46\,\textbf{D} & 47\,\textbf{A} & 48\,\textbf{C} & 49\,\textbf{B} & 50\,\textbf{A} \\
\hline
51\,\textbf{A} & 52\,\textbf{B} & 53\,\textbf{B} & 54\,\textbf{D} & 55\,\textbf{C} & 56\,\textbf{D} & 57\,\textbf{A} & 58\,\textbf{C} & 59\,\textbf{B} & 60\,\textbf{D} \\
\hline
61\,\textbf{D} & 62\,\textbf{B} & 63\,\textbf{D} & 64\,\textbf{B} & 65\,\textbf{B} & 66\,\textbf{C} & 67\,\textbf{D} & 68\,\textbf{A} & 69\,\textbf{B} & 70\,\textbf{B} \\
\hline
71\,\textbf{D} & 72\,\textbf{A} & 73\,\textbf{D} & 74\,\textbf{C} & 75\,\textbf{A} & 76\,\textbf{C} & 77\,\textbf{A} & 78\,\textbf{A} & 79\,\textbf{C} & 80\,\textbf{C} \\
\hline\end{tabular}\end{center}
\newpage
\section*{Answer Bank --- questions with detailed solutions}\label{bank}

\begin{sectionbanner}SECTION A \enspace$\cdot$\enspace Mathematics for the QFT \hfill Q.1--20\end{sectionbanner}

\needspace{13\baselineskip}
\begin{qitem}{1}
\qmeta{Single-correct MCQ}{Weighted sum of roots of unity}
Let $\omega = e^{2\pi i/8}$. The value of $\displaystyle S = \sum_{k=0}^{7} k\,\omega^{k}$ is
\par\smallskip\noindent\begin{tabularx}{\linewidth}{@{}*{4}{L}@{}}
\correctlabel{A}\,$\dfrac{8}{\omega - 1}$ & \optlabel{B}\,$\dfrac{8}{1-\omega}$ & \optlabel{C}\,$0$ & \optlabel{D}\,$\dfrac{8\omega}{\omega-1}$
\end{tabularx}
\end{qitem}
\begin{solbox}{A}
\textbf{\color{brand}What you need to know.} Geometric series $\sum_{k=0}^{M-1}z^k = \frac{1-z^M}{1-z}$; differentiating it gives weighted sums. For $\omega = e^{2\pi i/8}$, $\omega^8 = 1$ and $\omega^7 = \omega^{-1}$.

\textbf{\color{brand}Step-by-step solution.} Use $\sum_{k=0}^{M-1} k z^k = \dfrac{z\big(1 - M z^{M-1} + (M-1) z^{M}\big)}{(1-z)^2}$ (differentiate the geometric
series $\sum z^k = \frac{1-z^M}{1-z}$ and multiply by $z$). With $z = \omega$, $M = 8$, $\omega^8 = 1$ and $\omega^7 = \omega^{-1}$:
\[ S = \frac{\omega\,(1 - 8\omega^{-1} + 7)}{(1-\omega)^2} = \frac{8\omega - 8}{(1-\omega)^2} = \frac{-8(1-\omega)}{(1-\omega)^2}
     = \frac{8}{\omega - 1}. \]
\emph{Check:} $S \neq 0$ because the weights $k$ break the perfect cancellation of $\sum_k \omega^k = 0$;
$\frac{8}{1-\omega}$ has the wrong sign. Numerically $S = -4 + 9.657i$.

\textbf{\color{brand}Why the other options are wrong.}
\begin{itemize}[leftmargin=2.4em, itemsep=1pt, topsep=2pt]
\item[\textbf{(B)}] $\frac{8}{1-\omega}$ is the negative of the correct value: it comes from a sign slip when simplifying $\frac{8\omega - 8}{(1-\omega)^2}$ (note $8\omega - 8 = -8(1-\omega)$).
\item[\textbf{(C)}] $0$ would be right for the \emph{unweighted} sum $\sum_k\omega^k$; the weights $k$ destroy the symmetric cancellation (numerically $S = -4 + 9.66i \neq 0$).
\item[\textbf{(D)}] $\frac{8\omega}{\omega - 1}$ has an extra factor $\omega$: it forgets that $\omega\cdot\omega^{-1} = 1$ inside $\omega(1 - 8\omega^{-1} + 7)$.
\end{itemize}
\textbf{\color{brand}Remember.} Weighted root-of-unity sums: differentiate the geometric series, then use $\omega^N = 1$. Always check against the numerical value.
\end{solbox}
\qsep

\needspace{13\baselineskip}
\begin{qitem}{2}
\qmeta{Single-correct MCQ}{Orthogonality of roots of unity}
The number of ordered pairs $(x, y)$ with $x, y \in \{0, 1, \dots, 15\}$ for which
$\displaystyle \sum_{z=0}^{15} e^{2\pi i\,(x+y)z/16} \neq 0$ is
\par\smallskip\noindent\begin{tabularx}{\linewidth}{@{}*{4}{L}@{}}
\optlabel{A}\,$1$ & \correctlabel{B}\,$16$ & \optlabel{C}\,$15$ & \optlabel{D}\,$31$
\end{tabularx}
\end{qitem}
\begin{solbox}{B}
\textbf{\color{brand}What you need to know.} Orthogonality: $\sum_{z=0}^{N-1}\omega^{kz} = N$ if $N \mid k$, otherwise $0$.

\textbf{\color{brand}Step-by-step solution.} The sum is a sum of 16th roots of unity: $\sum_{z=0}^{15} \omega^{kz} = 16$ if $16 \mid k$ and $0$ otherwise
($\omega = e^{2\pi i/16}$, geometric series). Here $k = x + y \in \{0, \dots, 30\}$, so we need $x + y \in \{0, 16\}$.
$x + y = 0$: one pair $(0,0)$. $x + y = 16$: $x = 1, \dots, 15$ with $y = 16 - x$: 15 pairs. Total $1 + 15 = 16$.
(Equivalently: $y \equiv -x \pmod{16}$ has exactly one solution $y$ for every $x$.)

\textbf{\color{brand}Why the other options are wrong.}
\begin{itemize}[leftmargin=2.4em, itemsep=1pt, topsep=2pt]
\item[\textbf{(A)}] $1$ counts only $x + y = 0$ and forgets the pairs with $x + y = 16$.
\item[\textbf{(C)}] $15$ counts only $x + y = 16$ and forgets $(0, 0)$.
\item[\textbf{(D)}] $31$ treats every sum value $0, \dots, 30$ as non-zero, i.e.\ ignores the orthogonality rule.
\end{itemize}
\textbf{\color{brand}Remember.} For each $x$ there is exactly one $y \equiv -x \pmod N$, so the count is always $N$.
\end{solbox}
\qsep

\needspace{13\baselineskip}
\begin{qitem}{3}
\qmeta{Single-correct MCQ}{Binary fractions, periodicity of $e^{2\pi i\theta}$}
Let $x = 45$, written with $n = 6$ bits as $x_1x_2\dots x_6$. The value of $e^{2\pi i\,x/2^4}$ is
\par\smallskip\noindent\begin{tabularx}{\linewidth}{@{}*{4}{L}@{}}
\correctlabel{A}\,$e^{i\,13\pi/8}$ & \optlabel{B}\,$e^{i\,5\pi/8}$ & \optlabel{C}\,$e^{i\,3\pi/8}$ & \optlabel{D}\,$e^{-i\,13\pi/8}$
\end{tabularx}
\end{qitem}
\begin{solbox}{A}
\textbf{\color{brand}What you need to know.} $e^{2\pi i\,x/2^l}$ depends only on the fractional part of $x/2^l$, i.e.\ on the last $l$ bits of $x$.

\textbf{\color{brand}Step-by-step solution.} $45 = 101101_2$. Dividing by $2^4$ moves the binary point four places: $\dfrac{45}{16} = 10.1101_2 = 2 + 0.8125$.
Since $e^{2\pi i\cdot 2} = 1$, only the fractional part survives:
\[ e^{2\pi i\,45/16} = e^{2\pi i\,0.1101_2} = e^{2\pi i\,\frac{13}{16}} = e^{i\,13\pi/8}. \]
In QFT language: $e^{2\pi i\,x/2^l}$ depends only on the last $l$ bits $0.x_{n-l+1}\dots x_n$ (here $0.x_3x_4x_5x_6 = 0.1101_2$).

\textbf{\color{brand}Why the other options are wrong.}
\begin{itemize}[leftmargin=2.4em, itemsep=1pt, topsep=2pt]
\item[\textbf{(B)}] $e^{i5\pi/8}$ corresponds to $\frac{5}{16} = 0.0101_2$: it uses the wrong bits ($x_2\dots$) instead of the last four bits $1101$.
\item[\textbf{(C)}] $e^{i3\pi/8}$ corresponds to $\frac{3}{16}$, i.e.\ the bits $0011$ --- the bit string read wrongly.
\item[\textbf{(D)}] $e^{-i13\pi/8}$ has the wrong sign of the exponent; the convention uses $e^{+2\pi i\dots}$.
\end{itemize}
\textbf{\color{brand}Remember.} Divide by $2^l$ = move the binary point; drop the integer part; convert the fractional part to an angle.
\end{solbox}
\qsep

\needspace{13\baselineskip}
\begin{qitem}{4}
\qmeta{Single-correct MCQ}{Unitarity; the $4\times4$ Fourier matrix}
The matrix $U = \dfrac12\begin{pmatrix}1&1&1&1\\ 1&i&-1&-i\\ 1&-1&1&-1\\ 1&c&-1&d\end{pmatrix}$ is unitary. Then $(c, d)$ equals
\par\smallskip\noindent\begin{tabularx}{\linewidth}{@{}*{4}{L}@{}}
\optlabel{A}\,$(i,\ -i)$ & \optlabel{B}\,$(1,\ -1)$ & \optlabel{C}\,$(-1,\ 1)$ & \correctlabel{D}\,$(-i,\ i)$
\end{tabularx}
\end{qitem}
\begin{solbox}{D}
\textbf{\color{brand}What you need to know.} A matrix is unitary iff its rows (or columns) are orthonormal; the complex inner product conjugates the first vector.

\textbf{\color{brand}Step-by-step solution.} For a unitary matrix the rows are orthonormal. Let $r_4 = (1, c, -1, d)$.
\[ \langle r_1, r_4\rangle = 1 + c - 1 + d = c + d = 0 \;\Rightarrow\; d = -c. \]
\[ \langle r_2, r_4\rangle = \bar1\cdot1 + \bar i\,c + \overline{(-1)}(-1) + \overline{(-i)}\,d = 2 - ic + id = 2 - 2ic = 0 \;\Rightarrow\; c = \tfrac{1}{i} = -i, \]
hence $d = i$. Check $\langle r_3, r_4\rangle = 1 - c - 1 - d = -(c+d) = 0$ and $\|r_4\|^2 = \tfrac14(1+1+1+1) = 1$.
$U$ is exactly the 2-qubit QFT matrix $F_4$ with entries $i^{xy}/2$ (row 4: $i^{3x} = 1, -i, -1, i$).

\textbf{\color{brand}Why the other options are wrong.}
\begin{itemize}[leftmargin=2.4em, itemsep=1pt, topsep=2pt]
\item[\textbf{(A)}] $(i, -i)$ satisfies $c + d = 0$ but gives $\langle r_2, r_4\rangle = 2 - ic + id = 2 + 2 = 4 \neq 0$; it comes from forgetting the conjugate on $i$.
\item[\textbf{(B)}] $(1, -1)$ gives $\langle r_2, r_4\rangle = 2 - i - i \neq 0$.
\item[\textbf{(C)}] $(-1, 1)$ gives $\langle r_2, r_4\rangle = 2 + i + i \neq 0$.
\end{itemize}
\textbf{\color{brand}Remember.} Use two orthogonality conditions to solve for two unknowns, then check the norm. Recognise the result as $F_4$.
\end{solbox}
\qsep

\needspace{13\baselineskip}
\begin{qitem}{5}
\qmeta{Single-correct MCQ}{Gauss sum, trace of $F_4$}
Let $F$ be the $4\times4$ unitary Fourier matrix with entries $F_{yx} = \tfrac12\, i^{xy}$ ($x, y = 0, \dots, 3$).
The trace of $F$ is
\par\smallskip\noindent\begin{tabularx}{\linewidth}{@{}*{4}{L}@{}}
\correctlabel{A}\,$1 + i$ & \optlabel{B}\,$2$ & \optlabel{C}\,$0$ & \optlabel{D}\,$1 - i$
\end{tabularx}
\end{qitem}
\begin{solbox}{A}
\textbf{\color{brand}What you need to know.} The trace is the sum of diagonal entries; $F_{xx} = \frac12 i^{x^2}$ (a Gauss sum). Also trace = sum of eigenvalues.

\textbf{\color{brand}Step-by-step solution.} $\operatorname{tr} F = \sum_x F_{xx} = \tfrac12\sum_{x=0}^{3} i^{x^2}$ (a \emph{Gauss sum}).
$x^2 = 0, 1, 4, 9$ give $i^0 = 1$, $i^1 = i$, $i^4 = 1$, $i^9 = i$. Hence $\operatorname{tr} F = \tfrac12(2 + 2i) = 1 + i$.
Consistency: $F^4 = I$, so its eigenvalues lie in $\{1, -1, i, -i\}$; the eigenvalues of this $F$ are $1, 1, -1, i$ and their
sum is $1 + i$.

\textbf{\color{brand}Why the other options are wrong.}
\begin{itemize}[leftmargin=2.4em, itemsep=1pt, topsep=2pt]
\item[\textbf{(B)}] $2$ treats $i^{x^2}$ as real $\pm 1$ values; $i^1$ and $i^9$ are $i$, not $1$.
\item[\textbf{(C)}] $0$ confuses the trace (a sum of $i^{x^2}$) with the vanishing sum $\sum_x i^x$.
\item[\textbf{(D)}] $1 - i$ is the complex conjugate: it would be the trace of $F^\dagger = F^{-1}$.
\end{itemize}
\textbf{\color{brand}Remember.} List $x^2 \bmod 4$ and add the corresponding powers of $i$.
\end{solbox}
\qsep

\needspace{13\baselineskip}
\begin{qitem}{6}
\qmeta{Single-correct MCQ}{Interference pattern $|\sum e^{ij\theta}|^2$}
Let $f(\theta) = \big|\,1 + e^{i\theta} + e^{2i\theta} + e^{3i\theta}\,\big|^2$. The maximum value of $f$ and the smallest
positive $\theta$ with $f(\theta) = 0$ are, respectively,
\par\smallskip\noindent\begin{tabularx}{\linewidth}{@{}*{4}{L}@{}}
\optlabel{A}\,$4$ and $\pi/2$ & \optlabel{B}\,$4$ and $\pi$ & \optlabel{C}\,$16$ and $\pi/4$ & \correctlabel{D}\,$16$ and $\pi/2$
\end{tabularx}
\end{qitem}
\begin{solbox}{D}
\textbf{\color{brand}What you need to know.} $|\sum_{j=0}^{K-1}e^{ij\theta}|^2 = \frac{\sin^2(K\theta/2)}{\sin^2(\theta/2)}$: peak $K^2$ at $\theta = 0$, zeros at $\theta = 2\pi m/K$.

\textbf{\color{brand}Step-by-step solution.} Geometric series: $\sum_{j=0}^{3} e^{ij\theta} = \dfrac{1 - e^{4i\theta}}{1 - e^{i\theta}}
= e^{3i\theta/2}\,\dfrac{\sin 2\theta}{\sin(\theta/2)}$, so
\[ f(\theta) = \frac{\sin^2 2\theta}{\sin^2(\theta/2)} . \]
Maximum: all four arrows aligned at $\theta = 0$, $f = |4|^2 = 16$ (by the triangle inequality $f \le 16$).
Zeros: $\sin 2\theta = 0$ with $\sin(\theta/2) \neq 0$, i.e. $\theta = \pi/2, \pi, 3\pi/2, \dots$ The smallest is $\pi/2$
(four arrows pointing in the four directions $1, i, -1, -i$ cancel). This $\sin^2(K\theta/2)/\sin^2(\theta/2)$ shape is the
peak shape of every QFT output.

\textbf{\color{brand}Why the other options are wrong.}
\begin{itemize}[leftmargin=2.4em, itemsep=1pt, topsep=2pt]
\item[\textbf{(A)}] $4$ and $\pi/2$: $4$ is the maximum of $|\sum|$ (the modulus), not of $|\sum|^2$.
\item[\textbf{(B)}] $4$ and $\pi$: wrong maximum, and $\pi$ is a zero but not the smallest positive one.
\item[\textbf{(C)}] $16$ and $\pi/4$: at $\theta = \pi/4$ the four arrows span only $135^\circ$ and do not cancel ($f(\pi/4) = \sin^2(\pi/2)/\sin^2(\pi/8) \approx 6.83 \neq 0$; the zeros are at multiples of $2\pi/4$).
\end{itemize}
\textbf{\color{brand}Remember.} $K$ arrows cancel first when they are spread evenly round the circle: $\theta = 2\pi/K$.
\end{solbox}
\qsep

\needspace{13\baselineskip}
\begin{qitem}{7}
\qmeta{Single-correct MCQ}{Continued fractions, convergents}
The best rational approximation $p/q$ to $\dfrac{11}{64}$ with $1 \le q \le 10$ (as produced by the continued-fraction
algorithm) is
\par\smallskip\noindent\begin{tabularx}{\linewidth}{@{}*{4}{L}@{}}
\optlabel{A}\,$2/9$ & \correctlabel{B}\,$1/6$ & \optlabel{C}\,$1/7$ & \optlabel{D}\,$1/5$
\end{tabularx}
\end{qitem}
\begin{solbox}{B}
\textbf{\color{brand}What you need to know.} Continued fraction expansion and its convergents give the best rational approximations with bounded denominator.

\textbf{\color{brand}Step-by-step solution.} Continued fraction of $\frac{11}{64}$: $64 = 5\cdot 11 + 9$, $11 = 1\cdot 9 + 2$, $9 = 4\cdot 2 + 1$, $2 = 2\cdot 1$, so
$\frac{11}{64} = [0;\,5,\,1,\,4,\,2]$. Convergents: $\frac{0}{1}, \frac15, \frac16, \frac{5}{29}, \frac{11}{64}$.
The last convergent with $q \le 10$ is $\frac16$. Direct check: $\left|\frac{11}{64} - \frac16\right| = \frac{2}{384} \approx 0.0052$,
while $\left|\frac{11}{64} - \frac15\right| \approx 0.028$, $\frac17$ differs by $0.029$ and $\frac29$ by $0.050$.
This is exactly the step that turns a QFT measurement $y = 11$ out of $M = 64$ into the period $r = 6$.

\textbf{\color{brand}Why the other options are wrong.}
\begin{itemize}[leftmargin=2.4em, itemsep=1pt, topsep=2pt]
\item[\textbf{(A)}] $2/9$ is not a convergent and is far away ($0.050$).
\item[\textbf{(C)}] $1/7$ is not a convergent; its error $0.029$ is much larger than that of $1/6$.
\item[\textbf{(D)}] $1/5$ is an earlier convergent; $1/6$ has an allowed denominator and is closer ($0.0052$ vs $0.028$).
\end{itemize}
\textbf{\color{brand}Remember.} Compute $[a_0; a_1, \dots]$, build convergents, keep the last one with denominator below the bound.
\end{solbox}
\qsep

\needspace{13\baselineskip}
\begin{qitem}{8}
\qmeta{Single-correct MCQ}{Multiplicative order; Shor's gcd step}
The order $r$ of $a = 3$ modulo $N = 35$ (the least $r > 0$ with $a^r \equiv 1 \pmod{35}$), and the numbers
$\gcd(a^{r/2} - 1, N)$ and $\gcd(a^{r/2} + 1, N)$, are
\par\smallskip\noindent\begin{tabularx}{\linewidth}{@{}*{2}{L}@{}}
\optlabel{A}\,$r = 6$; $\;5$ and $7$ & \optlabel{B}\,$r = 12$; $\;1$ and $35$ \\[3pt]
\correctlabel{C}\,$r = 12$; $\;7$ and $5$ & \optlabel{D}\,$r = 4$; $\;5$ and $7$
\end{tabularx}
\end{qitem}
\begin{solbox}{C}
\textbf{\color{brand}What you need to know.} Order $r$: least power with $a^r \equiv 1$; $r \mid \varphi(N)$. Shor's step: if $r$ even and $a^{r/2} \not\equiv -1$, $\gcd(a^{r/2}\pm1, N)$ are factors.

\textbf{\color{brand}Step-by-step solution.} Powers of 3 modulo 35: $3, 9, 27, 11, 33, 29, 17, 16, 13, 4, 12, 1$, so $r = 12$ (by Euler, $r \mid \varphi(35) = 24$;
$3^{6} \equiv 29 \not\equiv 1$ and $3^{4} \equiv 11 \not\equiv 1$, so $r$ is neither 6 nor 4).
$a^{r/2} = 3^6 \equiv 29 \pmod{35}$, which is not $\equiv -1 = 34$, so Shor's reduction works:
\[ \gcd(29 - 1, 35) = \gcd(28, 35) = 7, \qquad \gcd(29 + 1, 35) = \gcd(30, 35) = 5 . \]
Indeed $35 = 5 \times 7$.

\textbf{\color{brand}Why the other options are wrong.}
\begin{itemize}[leftmargin=2.4em, itemsep=1pt, topsep=2pt]
\item[\textbf{(A)}] $r = 6$ is wrong because $3^6 \equiv 29 \not\equiv 1 \pmod{35}$.
\item[\textbf{(B)}] $r = 12$ with gcds $1$ and $35$ would mean $a^{r/2} \equiv \pm1$; but $3^6 \equiv 29$, and the gcds are $7$ and $5$.
\item[\textbf{(D)}] $r = 4$ is wrong because $3^4 = 81 \equiv 11 \not\equiv 1$.
\end{itemize}
\textbf{\color{brand}Remember.} List powers until you return to 1; test candidates that divide $\varphi(N)$; then take the two gcds.
\end{solbox}
\qsep

\needspace{13\baselineskip}
\begin{qitem}{9}
\qmeta{Single-correct MCQ}{Parameters of an $n$-qubit state}
A general pure state of 10 qubits is described by $D$ complex amplitudes. After removing normalisation and
the unobservable global phase, the number of independent \emph{real} parameters is $P$. Then $(D, P)$ is
\par\smallskip\noindent\begin{tabularx}{\linewidth}{@{}*{4}{L}@{}}
\optlabel{A}\,$(1024,\ 2048)$ & \correctlabel{B}\,$(1024,\ 2046)$ & \optlabel{C}\,$(1024,\ 2047)$ & \optlabel{D}\,$(20,\ 40)$
\end{tabularx}
\end{qitem}
\begin{solbox}{B}
\textbf{\color{brand}What you need to know.} $n$ qubits $\Rightarrow 2^n$ complex amplitudes; normalisation and global phase each remove one real parameter.

\textbf{\color{brand}Step-by-step solution.} $n$ qubits span a $2^n$-dimensional space: $D = 2^{10} = 1024$ complex numbers $= 2048$ real numbers.
Normalisation $\sum|\alpha_x|^2 = 1$ removes one real parameter and the global phase $e^{i\gamma}$ removes another:
$P = 2048 - 2 = 2046$. (Option $(1024, 2047)$ forgets the global phase; $(20, 40)$ counts as if qubits were independent.)

\textbf{\color{brand}Why the other options are wrong.}
\begin{itemize}[leftmargin=2.4em, itemsep=1pt, topsep=2pt]
\item[\textbf{(A)}] $(1024, 2048)$ removes neither normalisation nor global phase.
\item[\textbf{(C)}] $(1024, 2047)$ removes normalisation but forgets that the global phase is unobservable.
\item[\textbf{(D)}] $(20, 40)$ counts as if the qubits were independent (product states only); it ignores entanglement.
\end{itemize}
\textbf{\color{brand}Remember.} $2^{n+1} - 2$ real parameters: exponential in $n$, the root of quantum computing's power.
\end{solbox}
\qsep

\needspace{13\baselineskip}
\begin{qitem}{10}
\qmeta{Single-correct MCQ}{Walsh--Hadamard signs}
For $n = 4$, $H^{\otimes 4}|x\rangle = \frac14\sum_{y}(-1)^{x\cdot y}|y\rangle$ where $x\cdot y = \sum_j x_jy_j$.
For $x = 1011_2$, the number of $y \in \{0,\dots,15\}$ with a \emph{negative} amplitude is
\par\smallskip\noindent\begin{tabularx}{\linewidth}{@{}*{4}{L}@{}}
\optlabel{A}\,$4$ & \correctlabel{B}\,$8$ & \optlabel{C}\,$12$ & \optlabel{D}\,$3$
\end{tabularx}
\end{qitem}
\begin{solbox}{B}
\textbf{\color{brand}What you need to know.} $H^{\otimes n}\ket{x} = \frac{1}{\sqrt N}\sum_y (-1)^{x\cdot y}\ket{y}$; the sign is negative iff $x\cdot y$ is odd.

\textbf{\color{brand}Step-by-step solution.} The amplitude is negative iff $x\cdot y = y_1 + y_3 + y_4$ is odd (because $x = 1011$ has $x_1 = x_3 = x_4 = 1$).
$y_2$ is free (2 choices) and among the 8 choices of $(y_1, y_3, y_4)$ exactly half have odd sum, so the count is $2 \times 4 = 8$.
In general, for any $x \neq 0$ exactly $2^{n-1}$ of the $2^n$ signs are negative (the characters $(-1)^{x\cdot y}$ of
$\mathbb{Z}_2^n$ sum to zero).

\textbf{\color{brand}Why the other options are wrong.}
\begin{itemize}[leftmargin=2.4em, itemsep=1pt, topsep=2pt]
\item[\textbf{(A)}] $4$ ignores the free bit $y_2$ (which doubles the count).
\item[\textbf{(C)}] $12$ counts the wrong parity or adds the free bit twice.
\item[\textbf{(D)}] $3$ is just the number of 1-bits in $x$, unrelated to the sign count.
\end{itemize}
\textbf{\color{brand}Remember.} For any $x \neq 0$, exactly half of the $2^n$ Hadamard signs are negative.
\end{solbox}
\qsep

\needspace{13\baselineskip}
\begin{qitem}{11}
\qmeta{Single-correct MCQ}{FFT vs.\ QFT cost (and its limits)}
A classical FFT on a vector of length $2^n$ needs about $n\,2^n$ operations, while the QFT circuit needs
$\frac{n(n+1)}{2} + \lfloor n/2\rfloor$ gates. For $n = 20$ the ratio (FFT operations)/(QFT gates) is closest to
\par\smallskip\noindent\begin{tabularx}{\linewidth}{@{}*{4}{L}@{}}
\optlabel{A}\,$10^{9}$ & \optlabel{B}\,$10^{3}$ & \correctlabel{C}\,$10^{5}$ & \optlabel{D}\,$10^{7}$
\end{tabularx}
\end{qitem}
\begin{solbox}{C}
\textbf{\color{brand}What you need to know.} FFT cost $n2^n$; QFT gates $\frac{n(n+1)}{2} + \lfloor n/2\rfloor$. Compare orders of magnitude.

\textbf{\color{brand}Step-by-step solution.} FFT: $20 \times 2^{20} = 20 \times 1\,048\,576 = 20\,971\,520$. QFT: $\frac{20\cdot21}{2} + 10 = 210 + 10 = 220$ gates.
Ratio $= 20\,971\,520 / 220 \approx 9.5 \times 10^{4} \approx 10^5$. Caveat (important for interpretation): the QFT output is
stored in amplitudes and cannot be read out completely, so this ratio is \emph{not} a speed-up for computing all coefficients.

\textbf{\color{brand}Why the other options are wrong.}
\begin{itemize}[leftmargin=2.4em, itemsep=1pt, topsep=2pt]
\item[\textbf{(A)}] $10^9$ uses $N^2$ or similar instead of $n2^n$.
\item[\textbf{(B)}] $10^3$ underestimates: $2.1\times10^7/220 \approx 9.5\times10^4$.
\item[\textbf{(D)}] $10^7$ forgets to divide by the QFT gate count.
\end{itemize}
\textbf{\color{brand}Remember.} Exponentially fewer gates --- but the QFT output cannot be read out completely, so this is not a speed-up for computing all coefficients.
\end{solbox}
\qsep

\needspace{13\baselineskip}
\begin{qitem}{12}
\qmeta{Single-correct MCQ}{Cyclotomic product identity}
With $\omega = e^{2\pi i/8}$, the value of $\displaystyle\prod_{k=1}^{7}\big|1 - \omega^{k}\big|$ is
\par\smallskip\noindent\begin{tabularx}{\linewidth}{@{}*{4}{L}@{}}
\correctlabel{A}\,$8$ & \optlabel{B}\,$0$ & \optlabel{C}\,$1$ & \optlabel{D}\,$7$
\end{tabularx}
\end{qitem}
\begin{solbox}{A}
\textbf{\color{brand}What you need to know.} $z^N - 1 = \prod_k(z - \omega^k)$, so $1 + z + \dots + z^{N-1} = \prod_{k\ge1}(z - \omega^k)$; set $z = 1$.

\textbf{\color{brand}Step-by-step solution.} The roots of $z^8 - 1$ are $\omega^0, \dots, \omega^7$, so $z^8 - 1 = \prod_{k=0}^{7}(z - \omega^k)$ and
\[ \frac{z^8-1}{z-1} = 1 + z + \dots + z^7 = \prod_{k=1}^{7}(z - \omega^k). \]
Put $z = 1$: $\prod_{k=1}^{7}(1-\omega^k) = 8$, hence the product of the moduli is $8$. (Geometrically: the product of the
distances from one vertex of a regular octagon inscribed in the unit circle to the other seven vertices.)

\textbf{\color{brand}Why the other options are wrong.}
\begin{itemize}[leftmargin=2.4em, itemsep=1pt, topsep=2pt]
\item[\textbf{(B)}] $0$ would require some $\omega^k = 1$ for $k = 1, \dots, 7$, which is false.
\item[\textbf{(C)}] $1$ wrongly assumes the product of distances is 1 because each $|\omega^k| = 1$.
\item[\textbf{(D)}] $7$ counts the factors instead of evaluating the polynomial at $z = 1$.
\end{itemize}
\textbf{\color{brand}Remember.} Products of $(1 - \omega^k)$ come from the factorisation of $z^N - 1$.
\end{solbox}
\qsep

\needspace{13\baselineskip}
\begin{qitem}{13}
\qmeta{Single-correct MCQ}{DFT of a vector; Parseval}
The normalised vector $\alpha = \frac13(1, 2, 2, 0)$ is transformed by the unitary $4\times4$ Fourier matrix
$\beta_y = \frac12\sum_{x}\alpha_x\, i^{xy}$. The value of $|\beta_0|^2$ is
\par\smallskip\noindent\begin{tabularx}{\linewidth}{@{}*{4}{L}@{}}
\optlabel{A}\,$\dfrac56$ & \optlabel{B}\,$\dfrac{25}{12}$ & \optlabel{C}\,$\dfrac14$ & \correctlabel{D}\,$\dfrac{25}{36}$
\end{tabularx}
\end{qitem}
\begin{solbox}{D}
\textbf{\color{brand}What you need to know.} $\beta_0 = \frac{1}{\sqrt N}\sum_x\alpha_x$; probability $= |\text{amplitude}|^2$; Parseval $\sum|\beta|^2 = 1$.

\textbf{\color{brand}Step-by-step solution.} $\beta_0 = \frac12\sum_x \alpha_x = \frac12\cdot\frac{1+2+2+0}{3} = \frac56$, so $|\beta_0|^2 = \frac{25}{36}$.
Check the other options: $\frac56$ is the amplitude, not the probability; $\frac{25}{12} > 1$ is impossible because
$\sum_y|\beta_y|^2 = \sum_x|\alpha_x|^2 = 1$ (Parseval, $F$ is unitary). The remaining probability $\frac{11}{36}$ is spread over $y = 1,2,3$.

\textbf{\color{brand}Why the other options are wrong.}
\begin{itemize}[leftmargin=2.4em, itemsep=1pt, topsep=2pt]
\item[\textbf{(A)}] $\frac56$ is the amplitude $\beta_0$, not its squared modulus.
\item[\textbf{(B)}] $\frac{25}{12} > 1$ cannot be a probability; it forgets the factor $\frac12$.
\item[\textbf{(C)}] $\frac14$ assumes a uniform output, true only for basis-state inputs.
\end{itemize}
\textbf{\color{brand}Remember.} Zero frequency = average of the amplitudes; always square at the end.
\end{solbox}
\qsep

\needspace{13\baselineskip}
\begin{qitem}{14}
\qmeta{Single-correct MCQ}{Binary expansion of $2^k\varphi$}
Let $\varphi = 0.\varphi_1\varphi_2\varphi_3$ (binary, $\varphi_j \in \{0,1\}$). Then $e^{2\pi i\,4\varphi}$ equals
\par\smallskip\noindent\begin{tabularx}{\linewidth}{@{}*{4}{L}@{}}
\optlabel{A}\,$e^{2\pi i\,0.\varphi_2\varphi_3}$ & \optlabel{B}\,$i^{\varphi_2}$ & \optlabel{C}\,$(-1)^{\varphi_1}$ & \correctlabel{D}\,$(-1)^{\varphi_3}$
\end{tabularx}
\end{qitem}
\begin{solbox}{D}
\textbf{\color{brand}What you need to know.} Multiplying a binary fraction by $2^k$ moves the binary point $k$ places right; whole turns vanish.

\textbf{\color{brand}Step-by-step solution.} $4\varphi = \varphi_1\varphi_2.\varphi_3$ in binary, i.e. $4\varphi = (2\varphi_1 + \varphi_2) + 0.\varphi_3$. The integer part is a
whole number of turns, so $e^{2\pi i\,4\varphi} = e^{2\pi i\,0.\varphi_3} = e^{i\pi\varphi_3} = (-1)^{\varphi_3}$.
This is why, in phase estimation, the counting qubit controlling $U^{4}$ carries only the \emph{last} bit $\varphi_3$, and the
inverse QFT reads the bits from there.

\textbf{\color{brand}Why the other options are wrong.}
\begin{itemize}[leftmargin=2.4em, itemsep=1pt, topsep=2pt]
\item[\textbf{(A)}] $e^{2\pi i\,0.\varphi_2\varphi_3}$ is $e^{2\pi i\,2\varphi}$ (shift by one place), not by two.
\item[\textbf{(B)}] $i^{\varphi_2}$ mixes the wrong bit and the wrong angle.
\item[\textbf{(C)}] $(-1)^{\varphi_1}$ would be $e^{2\pi i\varphi}$ for $\varphi = 0.\varphi_1$ only; after multiplying by 4 the bit $\varphi_1$ is in the integer part.
\end{itemize}
\textbf{\color{brand}Remember.} $U^{2^k}$ in phase estimation sees only the bits after position $k$.
\end{solbox}
\qsep

\needspace{13\baselineskip}
\begin{qitem}{15}
\qmeta{Matching lists}{Algebra of 8th roots of unity}
Let $\omega = e^{2\pi i/8}$. Match List-I with List-II.
\par\smallskip\noindent{\renewcommand{\arraystretch}{1.35}%
\begin{tabularx}{\linewidth}{@{}c>{\raggedright\arraybackslash}p{0.33\linewidth}@{\hspace{1.2em}}cL@{}}
\toprule
\multicolumn{2}{@{}l}{\textbf{List-I}} & \multicolumn{2}{l@{}}{\textbf{List-II}} \\ \midrule
\textbf{(P)} & $\omega^{4}$ & \textbf{(1)} & $\sqrt2$ \\
\textbf{(Q)} & $\sum_{k=0}^{7}\omega^{2k}$ & \textbf{(2)} & $-1$ \\
\textbf{(R)} & $\omega\cdot\omega^{7}$ & \textbf{(3)} & $0$ \\
\textbf{(S)} & $\omega + \omega^{7}$ & \textbf{(4)} & $1$ \\
\bottomrule
\end{tabularx}}
\par\smallskip\noindent\begin{tabularx}{\linewidth}{@{}*{2}{L}@{}}
\optlabel{A}\,P$\to$1, Q$\to$3, R$\to$4, S$\to$2 & \optlabel{B}\,P$\to$2, Q$\to$4, R$\to$3, S$\to$1 \\[3pt]
\correctlabel{C}\,P$\to$2, Q$\to$3, R$\to$4, S$\to$1 & \optlabel{D}\,P$\to$3, Q$\to$2, R$\to$1, S$\to$4
\end{tabularx}
\end{qitem}
\begin{solbox}{C}
\textbf{\color{brand}What you need to know.} Powers of $\omega = e^{i\pi/4}$: $\omega^2 = i$, $\omega^4 = -1$, $\omega^8 = 1$, $\omega + \bar\omega = 2\cos(\pi/4)$.

\textbf{\color{brand}Step-by-step solution.} P: $\omega^4 = e^{i\pi} = -1$ (2). Q: $\omega^2 = i$ is a primitive 4th root, and $\sum_{k=0}^{7} i^k = 2(1 + i - 1 - i) = 0$ (3).
R: $\omega\,\omega^7 = \omega^8 = 1$ (4). S: $\omega + \omega^7 = \omega + \bar\omega = 2\cos\frac{\pi}{4} = \sqrt2$ (1).
Hence P$\to$2, Q$\to$3, R$\to$4, S$\to$1.

\textbf{\color{brand}Why the other options are wrong.}
\begin{itemize}[leftmargin=2.4em, itemsep=1pt, topsep=2pt]
\item[\textbf{(A)}] P$\to$1 says $\omega^4 = \sqrt2$; it confuses a power with the sum $\omega + \omega^7$.
\item[\textbf{(B)}] Q$\to$4 says $\sum\omega^{2k} = 1$; the sum of a full cycle of powers of $i$ is $0$.
\item[\textbf{(D)}] P$\to$3 says $\omega^4 = 0$; a root of unity has modulus 1, never 0.
\end{itemize}
\textbf{\color{brand}Remember.} In matching questions solve the easiest row (here R: $\omega^8 = 1$) and eliminate options.
\end{solbox}
\qsep

\needspace{13\baselineskip}
\begin{qitem}{16}
\qmeta{Matching lists}{Properties of $F$, $F^2$, $H\otimes H$, C$T$}
$F$ is the $4\times4$ unitary Fourier matrix ($F_{yx} = i^{xy}/2$), $H$ the Hadamard matrix and
$CT = \operatorname{diag}(1, 1, 1, e^{i\pi/4})$. Match each matrix in List-I with the \emph{only} property in List-II that it has.
\par\smallskip\noindent{\renewcommand{\arraystretch}{1.35}%
\begin{tabularx}{\linewidth}{@{}c>{\raggedright\arraybackslash}p{0.33\linewidth}@{\hspace{1.2em}}cL@{}}
\toprule
\multicolumn{2}{@{}l}{\textbf{List-I}} & \multicolumn{2}{l@{}}{\textbf{List-II}} \\ \midrule
\textbf{(P)} & $F$ & \textbf{(1)} & a permutation matrix that is its own inverse \\
\textbf{(Q)} & $F^2$ & \textbf{(2)} & real, symmetric, its own inverse, but not a permutation matrix \\
\textbf{(R)} & $H\otimes H$ & \textbf{(3)} & diagonal with determinant $e^{i\pi/4}$ \\
\textbf{(S)} & $CT$ & \textbf{(4)} & $M^4 = I$ but $M^2 \neq I$ \\
\bottomrule
\end{tabularx}}
\par\smallskip\noindent\begin{tabularx}{\linewidth}{@{}*{2}{L}@{}}
\optlabel{A}\,P$\to$4, Q$\to$2, R$\to$1, S$\to$3 & \correctlabel{B}\,P$\to$4, Q$\to$1, R$\to$2, S$\to$3 \\[3pt]
\optlabel{C}\,P$\to$3, Q$\to$1, R$\to$2, S$\to$4 & \optlabel{D}\,P$\to$1, Q$\to$4, R$\to$2, S$\to$3
\end{tabularx}
\end{qitem}
\begin{solbox}{B}
\textbf{\color{brand}What you need to know.} $F^2\ket{x} = \ket{-x}$, $F^4 = I$; $H\otimes H$ is real symmetric self-inverse; $\det$ of a diagonal = product of entries.

\textbf{\color{brand}Step-by-step solution.} $(F^2)_{zx} = \frac14\sum_y i^{y(x+z)} = \delta_{z,\,-x \bmod 4}$, so $F^2$ is the permutation $|x\rangle \mapsto |{-x}\rangle$
($0\to0$, $1\leftrightarrow3$, $2\to2$), which squares to $I$: Q$\to$1. Hence $F^4 = I$ while $F^2 \neq I$: P$\to$4.
$H\otimes H = \frac12\begin{psmallmatrix}1&1&1&1\\1&-1&1&-1\\1&1&-1&-1\\1&-1&-1&1\end{psmallmatrix}$ is real, symmetric and
$(H\otimes H)^2 = I$, but its entries are $\pm\frac12$, so it is not a permutation: R$\to$2.
$CT$ is diagonal with $\det = e^{i\pi/4}$; note $CT^4 = \operatorname{diag}(1,1,1,-1) \neq I$, so it does not have property 4: S$\to$3.

\textbf{\color{brand}Why the other options are wrong.}
\begin{itemize}[leftmargin=2.4em, itemsep=1pt, topsep=2pt]
\item[\textbf{(A)}] Q$\to$2 says $F^2$ is not a permutation; but $F^2$ maps $\ket{x}\mapsto\ket{-x}$, a permutation.
\item[\textbf{(C)}] P$\to$3 says $F$ is diagonal; $F$ is a full matrix.
\item[\textbf{(D)}] P$\to$1 says $F$ is a permutation; $F$ has entries of modulus $\frac12$, not 0/1.
\end{itemize}
\textbf{\color{brand}Remember.} Squaring the Fourier matrix gives the reflection $x \mapsto -x$.
\end{solbox}
\qsep

\needspace{13\baselineskip}
\begin{qitem}{17}
\qmeta{Assertion--Reason}{Root sums when $\gcd(k,N)>1$}
\textbf{Assertion (A):} For every integer $k$ that is not a multiple of $N$, $\displaystyle\sum_{j=0}^{N-1} e^{2\pi i\,jk/N} = 0$.\\[1mm]
\textbf{Reason (R):} For every integer $k$ that is not a multiple of $N$, the $N$ numbers $e^{2\pi i\,jk/N}$
($j = 0, \dots, N-1$) are the $N$ distinct vertices of a regular $N$-gon centred at the origin.
\begin{itemize}[leftmargin=2.6em, labelsep=0.4em, itemsep=2pt, topsep=3pt]
\item[\optlabel{A}] Both Assertion (A) and Reason (R) are true, and R is the correct explanation of A
\item[\optlabel{B}] Both A and R are true, but R is \emph{not} the correct explanation of A
\item[\correctlabel{C}] A is true, but R is false
\item[\optlabel{D}] A is false, but R is true
\end{itemize}
\end{qitem}
\begin{solbox}{C}
\textbf{\color{brand}What you need to know.} Assertion-Reason: decide A and R separately, then whether R explains A. Roots repeat when $\gcd(k,N) > 1$.

\textbf{\color{brand}Step-by-step solution.} \textbf{A is true:} with $z = e^{2\pi i k/N} \neq 1$, $\sum_{j=0}^{N-1} z^j = \frac{1-z^N}{1-z} = \frac{1 - e^{2\pi i k}}{1-z} = 0$.
\textbf{R is false:} if $g = \gcd(k, N) > 1$ the numbers repeat. Example $N = 8$, $k = 2$: the values are
$1, i, -1, -i, 1, i, -1, -i$ --- a square traversed twice, not 8 distinct vertices of an octagon. (The sum is still zero,
because a regular $N/g$-gon is traversed $g$ times.) Answer: (C).

\textbf{\color{brand}Why the other options are wrong.}
\begin{itemize}[leftmargin=2.4em, itemsep=1pt, topsep=2pt]
\item[\textbf{(A)}] needs A true, R true \emph{and} R explaining A; R is false (counter-example $N=8$, $k=2$).
\item[\textbf{(B)}] needs A and R both true but R not explaining A; R is false.
\item[\textbf{(D)}] needs A false and R true; A is true (geometric series).
\end{itemize}
\textbf{\color{brand}Remember.} A true statement can be explained by a false-looking ``intuitive'' reason; check the reason on an example.
\end{solbox}
\qsep

\needspace{13\baselineskip}
\begin{qitem}{18}
\qmeta{Statement I / II}{Determinants of unitary matrices}
\textbf{Statement I:} Every unitary matrix $U$ satisfies $|\det U| = 1$.\\[1mm]
\textbf{Statement II:} The unitary Fourier matrix $F_N$ ($N = 2^n$, entries $\omega^{xy}/\sqrt N$) has $\det F_N = \pm 1$ for every $n$.
\begin{itemize}[leftmargin=2.6em, labelsep=0.4em, itemsep=2pt, topsep=3pt]
\item[\optlabel{A}] Both Statement I and Statement II are true
\item[\optlabel{B}] Both Statement I and Statement II are false
\item[\correctlabel{C}] Statement I is true, but Statement II is false
\item[\optlabel{D}] Statement I is false, but Statement II is true
\end{itemize}
\end{qitem}
\begin{solbox}{C}
\textbf{\color{brand}What you need to know.} $|\det U| = 1$ for unitaries; $F^4 = I$ so the eigenvalues lie in $\{\pm1, \pm i\}$.

\textbf{\color{brand}Step-by-step solution.} \textbf{I is true:} $|\det U|^2 = \det(U^\dagger U) = \det I = 1$.
\textbf{II is false:} the determinant is the product of the eigenvalues, which lie in $\{1,-1,i,-i\}$ because $F^4 = I$.
For $N = 4$ the eigenvalues are $1, 1, -1, i$ (compare the trace $1 + i$ in Q.5), so $\det F_4 = 1\cdot1\cdot(-1)\cdot i = -i \neq \pm1$.
(For $N = 2$, $F_2 = H$ and $\det H = -1$, so the claim happens to hold there --- one counterexample is enough.) Answer: (C).

\textbf{\color{brand}Why the other options are wrong.}
\begin{itemize}[leftmargin=2.4em, itemsep=1pt, topsep=2pt]
\item[\textbf{(A)}] needs both statements true; II fails for $N = 4$ ($\det F_4 = -i$).
\item[\textbf{(B)}] needs both statements false; I is true.
\item[\textbf{(D)}] needs I false and II true; I is true and II is false.
\end{itemize}
\textbf{\color{brand}Remember.} Modulus 1 does not mean real: a determinant can be $\pm i$.
\end{solbox}
\qsep

\needspace{20\baselineskip}
\begin{passagebox}{Paragraph for Questions 19 and 20}
Let $N = 32$ and define the sequence $a_x = 1$ if $x \equiv 2 \pmod 5$ and $a_x = 0$ otherwise, for
$x = 0, 1, \dots, 31$. Its (un-normalised) Fourier coefficients are
$A_y = \sum_{x=0}^{31} a_x\, e^{2\pi i\,xy/32}$, $\;y = 0, 1, \dots, 31$. (The sequence is ``periodic with period 5'', but
5 does not divide 32.)
\end{passagebox}
\needspace{13\baselineskip}
\begin{qitem}{19}
\qmeta{Paragraph based}{Parseval for the un-normalised DFT}
The value of $\displaystyle\sum_{y=0}^{31} |A_y|^2$ is
\par\smallskip\noindent\begin{tabularx}{\linewidth}{@{}*{4}{L}@{}}
\correctlabel{A}\,$192$ & \optlabel{B}\,$6$ & \optlabel{C}\,$36$ & \optlabel{D}\,$1152$
\end{tabularx}
\end{qitem}
\begin{solbox}{A}
\textbf{\color{brand}What you need to know.} Parseval for the un-normalised transform: $\sum_y|A_y|^2 = N\sum_x|a_x|^2$.

\textbf{\color{brand}Step-by-step solution.} The non-zero samples are $x = 2, 7, 12, 17, 22, 27$, i.e. $K = 6$ ones. For the un-normalised transform
$A = \sqrt{N}\,F a$ with unitary $F$, Parseval gives
\[ \sum_y |A_y|^2 = N\sum_x |a_x|^2 = 32 \times 6 = 192 . \]
(36 is $|A_0|^2$ alone; 6 is $\sum_x|a_x|^2$; 1152 would be $N \cdot 36$.)

\textbf{\color{brand}Why the other options are wrong.}
\begin{itemize}[leftmargin=2.4em, itemsep=1pt, topsep=2pt]
\item[\textbf{(B)}] $6$ is $\sum_x|a_x|^2$ without the factor $N$.
\item[\textbf{(C)}] $36$ is only $|A_0|^2$.
\item[\textbf{(D)}] $1152 = 32\times36$ multiplies the wrong quantity.
\end{itemize}
\textbf{\color{brand}Remember.} Count the non-zero samples, multiply by $N$.
\end{solbox}
\qsep

\needspace{13\baselineskip}
\begin{qitem}{20}
\qmeta{Paragraph based}{Spectral leakage ($r\nmid N$)}
Among $y = 1, 2, \dots, 15$, the value of $y$ for which $|A_y|$ is largest is
\begin{itemize}[leftmargin=2.6em, labelsep=0.4em, itemsep=2pt, topsep=3pt]
\item[\optlabel{A}] $6$
\item[\optlabel{B}] $16$ is the only peak; all of $1,\dots,15$ give $|A_y| = 0$
\item[\optlabel{C}] $7$
\item[\correctlabel{D}] $13$
\end{itemize}
\end{qitem}
\begin{solbox}{D}
\textbf{\color{brand}What you need to know.} With period $r \nmid N$ the peaks leak; $|A_y| = |\sin(\pi Kry/N)/\sin(\pi ry/N)|$ is largest when $ry/N$ is closest to an integer.

\textbf{\color{brand}Step-by-step solution.} $A_y = \sum_{j=0}^{5} e^{2\pi i (5j+2)y/32} = e^{4\pi i y/32}\sum_{j=0}^{5} e^{2\pi i\,5jy/32}$, so
\[ |A_y| = \left|\frac{\sin(6\cdot 5\pi y/32)}{\sin(5\pi y/32)}\right| = \left|\frac{\sin(30\pi y/32)}{\sin(5\pi y/32)}\right| . \]
Peaks occur where $5y/32$ is close to an integer, i.e. $y \approx 32k/5 = 6.4,\ 12.8,\ \dots$ The closer $5y/32$ is to an
integer, the higher the peak: $5\cdot13/32 = 2.03$ (distance $0.03$) beats $5\cdot6/32 = 0.94$ (distance $0.06$) and
$5\cdot7/32 = 1.09$. Numerically $|A_{13}| = 5.67$, $|A_6| = 4.74$, $|A_7| = 3.38$, $|A_{12}| = 1.85$; and $|A_{16}| = 0$ since
$\sin(15\pi) = 0$. Hence $y = 13$. Lesson: with a period that does not divide $N$, the QFT peaks ``leak'' and their
heights depend on how close $ky/N$ is to an integer.

\textbf{\color{brand}Why the other options are wrong.}
\begin{itemize}[leftmargin=2.4em, itemsep=1pt, topsep=2pt]
\item[\textbf{(A)}] $y = 6$: $5\cdot6/32 = 0.94$ is $0.06$ from an integer, twice as far as $y = 13$ ($0.03$); $|A_6| = 4.74 < 5.67$.
\item[\textbf{(B)}] $|A_y|$ is not zero for $y = 1..15$; and $y = 16$ actually gives $|A_{16}| = 0$ ($\sin 15\pi = 0$).
\item[\textbf{(C)}] $y = 7$: $5\cdot7/32 = 1.09$ is $0.09$ away; $|A_7| = 3.38$.
\end{itemize}
\textbf{\color{brand}Remember.} Rank peaks by the distance of $ry/N$ from the nearest integer.
\end{solbox}
\qsep

\newpage

\begin{sectionbanner}SECTION B \enspace$\cdot$\enspace Quantum Mechanics \hfill Q.21--40\end{sectionbanner}

\needspace{13\baselineskip}
\begin{qitem}{21}
\qmeta{Single-correct MCQ}{Born rule, partial measurement}
A two-qubit state is $\ket{\psi} = c\,\big(\ket{00} + 2i\ket{01} - \ket{10} + \sqrt2\,\ket{11}\big)$ with $c > 0$.
If only the \emph{first} qubit is measured in the computational basis, the probability of obtaining $1$ is
\par\smallskip\noindent\begin{tabularx}{\linewidth}{@{}*{4}{L}@{}}
\optlabel{A}\,$\dfrac12$ & \correctlabel{B}\,$\dfrac38$ & \optlabel{C}\,$\dfrac14$ & \optlabel{D}\,$\dfrac58$
\end{tabularx}
\end{qitem}
\begin{solbox}{B}
\textbf{\color{brand}What you need to know.} Normalise first; partial measurement probability = sum of $|c_x|^2$ over basis states with the required bit.

\textbf{\color{brand}Step-by-step solution.} Normalisation: $c^2\big(|1|^2 + |2i|^2 + |-1|^2 + |\sqrt2|^2\big) = c^2(1 + 4 + 1 + 2) = 8c^2 = 1$, so $c = \frac{1}{2\sqrt2}$.
By the Born rule, the probability that the first qubit gives 1 is the total weight of the basis states whose first bit is 1:
\[ P(1) = c^2\big(|-1|^2 + |\sqrt2|^2\big) = \frac{1 + 2}{8} = \frac38 . \]
(The phases $i$ and $-1$ do not affect probabilities in this basis; $\frac58$ is $P(0)$; $\frac14$ and $\frac12$ come from
forgetting to normalise or from counting states instead of weights.)

\textbf{\color{brand}Why the other options are wrong.}
\begin{itemize}[leftmargin=2.4em, itemsep=1pt, topsep=2pt]
\item[\textbf{(A)}] $\frac12$ counts basis states (2 of 4) instead of weights.
\item[\textbf{(C)}] $\frac14$ forgets to normalise or uses only one of the two terms.
\item[\textbf{(D)}] $\frac58$ is $P(\text{first qubit} = 0)$.
\end{itemize}
\textbf{\color{brand}Remember.} Probabilities come from squared moduli of normalised amplitudes; phases do not matter here.
\end{solbox}
\qsep

\needspace{13\baselineskip}
\begin{qitem}{22}
\qmeta{Single-correct MCQ}{Global vs.\ relative phase}
Which pair of single-qubit state vectors describes the \emph{same} physical state (i.e.\ gives identical
measurement statistics in every basis)?
\par\smallskip\noindent\begin{tabularx}{\linewidth}{@{}*{2}{L}@{}}
\correctlabel{A}\,$\ket{0} + i\ket{1}\;$ and $\;i\ket{0} - \ket{1}$ & \optlabel{B}\,$\ket{0} + i\ket{1}\;$ and $\;i\ket{0} + \ket{1}$ \\[3pt]
\optlabel{C}\,$\ket{0} + i\ket{1}\;$ and $\;\ket{0} - i\ket{1}$ & \optlabel{D}\,$\ket{0} + \ket{1}\;$ and $\;\ket{0} - \ket{1}$
\end{tabularx}
\end{qitem}
\begin{solbox}{A}
\textbf{\color{brand}What you need to know.} Two vectors describe the same state iff one is a constant phase times the other (component ratios all equal).

\textbf{\color{brand}Step-by-step solution.} Two vectors represent the same physical state iff they differ by an overall (global) factor $e^{i\gamma}$.
$i\ket{0} - \ket{1} = i\big(\ket{0} + i\ket{1}\big)$, since $i\cdot i = -1$: global phase $i$, so the pair is physically identical.
The other pairs differ by a \emph{relative} phase: $\ket{0} \pm i\ket{1}$ (relative phase $\pm i$), $\ket{0} \pm \ket{1}$
(relative phase $\pm1$), and $i\ket{0} + \ket{1} = i(\ket{0} - i\ket{1})$, which is $\ket{0} - i\ket{1}$ again.
A relative phase is measurable: e.g.\ $\ket{0}+\ket{1} \propto \ket{+}$ and $\ket{0} - \ket{1} \propto \ket{-}$ are perfectly
distinguishable in the $X$ basis. Relative phases are exactly what the QFT writes information into.

\textbf{\color{brand}Why the other options are wrong.}
\begin{itemize}[leftmargin=2.4em, itemsep=1pt, topsep=2pt]
\item[\textbf{(B)}] $i\ket0 + \ket1 = i(\ket0 - i\ket1)$, which equals $\ket0 - i\ket1$, not $\ket0 + i\ket1$.
\item[\textbf{(C)}] $\ket0 \pm i\ket1$: ratios $1$ and $-1$ differ (relative phase) --- they are orthogonal states.
\item[\textbf{(D)}] $\ket0 \pm \ket1$ are $\ket{+}$ and $\ket{-}$, orthogonal.
\end{itemize}
\textbf{\color{brand}Remember.} Global phase: invisible. Relative phase: physical.
\end{solbox}
\qsep

\needspace{13\baselineskip}
\begin{qitem}{23}
\qmeta{Single-correct MCQ}{Interference with a Hadamard}
The state $\frac{1}{\sqrt2}\big(\ket{0} + e^{i\phi}\ket{1}\big)$ passes through a Hadamard gate and is then measured
in the computational basis. The probability of outcome $0$ is
\par\smallskip\noindent\begin{tabularx}{\linewidth}{@{}*{4}{L}@{}}
\correctlabel{A}\,$\cos^2(\phi/2)$ & \optlabel{B}\,$\sin^2(\phi/2)$ & \optlabel{C}\,$\dfrac12$ & \optlabel{D}\,$\cos^2\phi$
\end{tabularx}
\end{qitem}
\begin{solbox}{A}
\textbf{\color{brand}What you need to know.} $H$ maps relative phase to probability: amplitudes $\frac{1\pm e^{i\phi}}{2}$; $|1+e^{i\phi}|^2 = 4\cos^2(\phi/2)$.

\textbf{\color{brand}Step-by-step solution.} With $H\ket{0} = \frac{\ket0 + \ket1}{\sqrt2}$ and $H\ket{1} = \frac{\ket0 - \ket1}{\sqrt2}$:
\[ H\,\frac{\ket{0} + e^{i\phi}\ket{1}}{\sqrt2} = \frac{1 + e^{i\phi}}{2}\ket{0} + \frac{1 - e^{i\phi}}{2}\ket{1}. \]
$|1 + e^{i\phi}|^2 = (1 + e^{i\phi})(1 + e^{-i\phi}) = 2 + 2\cos\phi = 4\cos^2(\phi/2)$, so
$P(0) = \frac{4\cos^2(\phi/2)}{4} = \cos^2(\phi/2)$ (and $P(1) = \sin^2(\phi/2)$).
The Hadamard converts an invisible \emph{relative phase} into a visible \emph{probability}: $\phi = 0$ gives $0$ surely,
$\phi = \pi$ gives $1$ surely. A single-qubit QFT is exactly this $H$.

\textbf{\color{brand}Why the other options are wrong.}
\begin{itemize}[leftmargin=2.4em, itemsep=1pt, topsep=2pt]
\item[\textbf{(B)}] $\sin^2(\phi/2)$ is $P(1)$.
\item[\textbf{(C)}] $\frac12$ ignores the interference produced by $H$ (true only for $\phi = \pm\pi/2$).
\item[\textbf{(D)}] $\cos^2\phi$ drops the half-angle.
\end{itemize}
\textbf{\color{brand}Remember.} Interference converts phase into measurable probability.
\end{solbox}
\qsep

\needspace{13\baselineskip}
\begin{qitem}{24}
\qmeta{Single-correct MCQ}{Bloch-sphere expectation values}
For the qubit state $\ket{\psi} = \cos\frac{\theta}{2}\ket{0} + e^{i\phi}\sin\frac{\theta}{2}\ket{1}$ the quantity
$\langle X\rangle + i\langle Y\rangle$, where $\langle A\rangle = \bra{\psi}A\ket{\psi}$ and $X, Y$ are Pauli matrices, equals
\par\smallskip\noindent\begin{tabularx}{\linewidth}{@{}*{4}{L}@{}}
\optlabel{A}\,$\cos\theta\; e^{i\phi}$ & \correctlabel{B}\,$\sin\theta\; e^{i\phi}$ & \optlabel{C}\,$\sin\frac{\theta}{2}\, e^{i\phi}$ & \optlabel{D}\,$\sin\theta\; e^{-i\phi}$
\end{tabularx}
\end{qitem}
\begin{solbox}{B}
\textbf{\color{brand}What you need to know.} $\langle X\rangle = 2\mathrm{Re}(\bar a b)$, $\langle Y\rangle = 2\mathrm{Im}(\bar a b)$, so $\langle X\rangle + i\langle Y\rangle = 2\bar a b$.

\textbf{\color{brand}Step-by-step solution.} Write $a = \cos\frac\theta2$, $b = e^{i\phi}\sin\frac\theta2$. With $X = \begin{psmallmatrix}0&1\\1&0\end{psmallmatrix}$,
$Y = \begin{psmallmatrix}0&-i\\i&0\end{psmallmatrix}$:
\[ \langle X\rangle = \bar a b + \bar b a = 2\,\mathrm{Re}(\bar a b), \qquad \langle Y\rangle = -i\bar a b + i\bar b a = 2\,\mathrm{Im}(\bar a b). \]
Hence $\langle X\rangle + i\langle Y\rangle = 2\bar a b = 2\cos\frac\theta2\sin\frac\theta2\,e^{i\phi} = \sin\theta\,e^{i\phi}$.
So $\langle X\rangle = \sin\theta\cos\phi$, $\langle Y\rangle = \sin\theta\sin\phi$ (and $\langle Z\rangle = \cos\theta$): the Bloch vector
$(\sin\theta\cos\phi, \sin\theta\sin\phi, \cos\theta)$. The \emph{azimuth} $\phi$ is the phase that QFT rotations change.

\textbf{\color{brand}Why the other options are wrong.}
\begin{itemize}[leftmargin=2.4em, itemsep=1pt, topsep=2pt]
\item[\textbf{(A)}] $\cos\theta\,e^{i\phi}$ uses $\langle Z\rangle$'s factor instead of $2\sin\frac\theta2\cos\frac\theta2$.
\item[\textbf{(C)}] $\sin\frac\theta2\,e^{i\phi}$ forgets the factor $2\cos\frac\theta2$.
\item[\textbf{(D)}] $\sin\theta\,e^{-i\phi}$ is $\langle X\rangle - i\langle Y\rangle$ (conjugate).
\end{itemize}
\textbf{\color{brand}Remember.} Bloch vector $(\sin\theta\cos\phi, \sin\theta\sin\phi, \cos\theta)$.
\end{solbox}
\qsep

\needspace{13\baselineskip}
\begin{qitem}{25}
\qmeta{Single-correct MCQ}{Entanglement test}
Which of the following two-qubit states is \emph{entangled} (cannot be written as $\ket{a}\otimes\ket{b}$)?
\par\smallskip\noindent\begin{tabularx}{\linewidth}{@{}*{2}{L}@{}}
\optlabel{A}\,$\frac12\big(\ket{00} - \ket{01} + \ket{10} - \ket{11}\big)$ & \optlabel{B}\,$\frac12\big(\ket{00} + \ket{01} + \ket{10} + \ket{11}\big)$ \\[3pt]
\correctlabel{C}\,$\frac12\big(\ket{00} + \ket{01} + \ket{10} - \ket{11}\big)$ & \optlabel{D}\,$\frac12\big(\ket{00} + i\ket{01} - \ket{10} - i\ket{11}\big)$
\end{tabularx}
\end{qitem}
\begin{solbox}{C}
\textbf{\color{brand}What you need to know.} Product test: $c_{00}c_{11} = c_{01}c_{10}$.

\textbf{\color{brand}Step-by-step solution.} For $\sum c_{jk}\ket{jk}$, a product $(a_0\ket0 + a_1\ket1)(b_0\ket0 + b_1\ket1)$ has $c_{jk} = a_jb_k$, so
$c_{00}c_{11} = a_0a_1b_0b_1 = c_{01}c_{10}$. Conversely, $c_{00}c_{11} = c_{01}c_{10}$ (the $2\times2$ coefficient matrix has
determinant 0, i.e.\ rank 1) guarantees a product. Test each option (factor $\frac12$ dropped):
\begin{itemize}\itemsep0pt
\item $(1, 1, 1, -1)$: $c_{00}c_{11} = -1 \neq 1 = c_{01}c_{10}$ $\Rightarrow$ \textbf{entangled}.
\item $(1, 1, 1, 1)$: $1 = 1$, product $\ket{+}\ket{+}$.
\item $(1, -1, 1, -1)$: $-1 = -1$, product $\ket{+}\ket{-}$.
\item $(1, i, -1, -i)$: $-i = -i$, product $\frac{1}{2}(\ket0 - \ket1)(\ket0 + i\ket1)$.
\end{itemize}
The entangled state is $\mathrm{CZ}\ket{+}\ket{+}$: a controlled-phase gate (like those in the QFT) \emph{can} entangle.
(This is not in conflict with Q.54: the QFT acting on a \emph{basis} state happens to give a product state.)

\textbf{\color{brand}Why the other options are wrong.}
\begin{itemize}[leftmargin=2.4em, itemsep=1pt, topsep=2pt]
\item[\textbf{(A)}] $(1,-1,1,-1)/2 = \ket+\ket-$: $-1 = -1$, product.
\item[\textbf{(B)}] $(1,1,1,1)/2 = \ket+\ket+$: $1 = 1$, product.
\item[\textbf{(D)}] $(1,i,-1,-i)/2 = \frac12(\ket0-\ket1)(\ket0+i\ket1)$: $-i = -i$, product.
\end{itemize}
\textbf{\color{brand}Remember.} One sign change in the uniform state creates entanglement ($CZ\ket+\ket+$).
\end{solbox}
\qsep

\needspace{13\baselineskip}
\begin{qitem}{26}
\qmeta{Single-correct MCQ}{Phase kickback (CNOT)}
A CNOT gate (qubit 1 = control, qubit 2 = target) acts on $\ket{+}\ket{-}$, where
$\ket{\pm} = \frac{1}{\sqrt2}(\ket0 \pm \ket1)$. The output is
\par\smallskip\noindent\begin{tabularx}{\linewidth}{@{}*{2}{L}@{}}
\optlabel{A}\,$\ket{+}\ket{-}$ & \optlabel{B}\,$\ket{-}\ket{+}$ \\[3pt]
\optlabel{C}\,$\frac{1}{\sqrt2}\big(\ket{0}\ket{-} + \ket{1}\ket{+}\big)$ & \correctlabel{D}\,$\ket{-}\ket{-}$
\end{tabularx}
\end{qitem}
\begin{solbox}{D}
\textbf{\color{brand}What you need to know.} $\ket-$ is an eigenvector of $X$ with eigenvalue $-1$; controlled gates kick eigenvalues back to the control.

\textbf{\color{brand}Step-by-step solution.} $\ket{-}$ is an eigenvector of $X$ with eigenvalue $-1$: $X\ket{-} = \frac{\ket1 - \ket0}{\sqrt2} = -\ket{-}$. Therefore
\[ \mathrm{CNOT}\,\ket{+}\ket{-} = \frac{1}{\sqrt2}\big(\ket0\ket{-} + \ket1\,X\ket{-}\big) = \frac{1}{\sqrt2}\big(\ket0 - \ket1\big)\ket{-}
   = \ket{-}\ket{-}. \]
The target is unchanged; the eigenvalue $-1$ has been ``kicked back'' onto the control as a relative phase. This
\emph{phase kickback} is the mechanism by which controlled-$U^{2^j}$ gates load the eigenphase into the counting
register in phase estimation (see Q.28).

\textbf{\color{brand}Why the other options are wrong.}
\begin{itemize}[leftmargin=2.4em, itemsep=1pt, topsep=2pt]
\item[\textbf{(A)}] $\ket+\ket-$ assumes CNOT does nothing because the target is unchanged --- but the control acquires a phase.
\item[\textbf{(B)}] $\ket-\ket+$ flips the wrong qubit.
\item[\textbf{(C)}] $\frac{1}{\sqrt2}(\ket0\ket- + \ket1\ket+)$ treats $X\ket- = \ket+$, which is false ($X\ket- = -\ket-$).
\end{itemize}
\textbf{\color{brand}Remember.} Phase kickback: the target stays, the control changes.
\end{solbox}
\qsep

\needspace{13\baselineskip}
\begin{qitem}{27}
\qmeta{Single-correct MCQ}{Controlled-phase matrices}
In the basis $\ket{00}, \ket{01}, \ket{10}, \ket{11}$ (first qubit = control), the matrix of controlled-$R_3$, where
$R_k = \operatorname{diag}\big(1, e^{2\pi i/2^k}\big)$, is
\par\smallskip\noindent\begin{tabularx}{\linewidth}{@{}*{2}{L}@{}}
\correctlabel{A}\,$\operatorname{diag}\big(1, 1, 1, e^{i\pi/4}\big)$ & \optlabel{B}\,$\operatorname{diag}\big(1, e^{i\pi/4}, 1, e^{i\pi/4}\big)$ \\[3pt]
\optlabel{C}\,$\operatorname{diag}\big(1, 1, e^{i\pi/4}, e^{i\pi/4}\big)$ & \optlabel{D}\,$\operatorname{diag}\big(1, 1, 1, e^{i\pi/8}\big)$
\end{tabularx}
\end{qitem}
\begin{solbox}{A}
\textbf{\color{brand}What you need to know.} $CU = \ket0\!\bra0\otimes I + \ket1\!\bra1\otimes U$; $R_3 = \operatorname{diag}(1, e^{2\pi i/8})$.

\textbf{\color{brand}Step-by-step solution.} $R_3 = \operatorname{diag}(1, e^{2\pi i/8}) = \operatorname{diag}(1, e^{i\pi/4})$ (the $T$ gate). Controlled-$R_3$ applies $R_3$
to the target only when the control is $\ket1$:
\[ \mathrm{C}R_3 = \ket0\!\bra0 \otimes I + \ket1\!\bra1 \otimes R_3 = \operatorname{diag}(1, 1, 1, e^{i\pi/4}). \]
Only $\ket{11}$ picks up a phase, so the gate is \emph{symmetric}: it does not matter which qubit is called the control.
$\operatorname{diag}(1, e^{i\pi/4}, 1, e^{i\pi/4}) = I\otimes R_3$ and $\operatorname{diag}(1,1,e^{i\pi/4},e^{i\pi/4}) = R_3 \otimes I$ are
uncontrolled; $e^{i\pi/8}$ confuses $2\pi/2^3$ with $\pi/2^3$.

\textbf{\color{brand}Why the other options are wrong.}
\begin{itemize}[leftmargin=2.4em, itemsep=1pt, topsep=2pt]
\item[\textbf{(B)}] $\operatorname{diag}(1,e^{i\pi/4},1,e^{i\pi/4}) = I\otimes R_3$ (uncontrolled).
\item[\textbf{(C)}] $\operatorname{diag}(1,1,e^{i\pi/4},e^{i\pi/4}) = R_3\otimes I$ (uncontrolled).
\item[\textbf{(D)}] $e^{i\pi/8}$ uses $\pi/2^3$ instead of $2\pi/2^3$.
\end{itemize}
\textbf{\color{brand}Remember.} A controlled phase changes only $\ket{11}$, and is symmetric in the two qubits.
\end{solbox}
\qsep

\needspace{13\baselineskip}
\begin{qitem}{28}
\qmeta{Single-correct MCQ}{Kickback of controlled-$U^{2^j}$}
$U\ket{u} = e^{2\pi i\varphi}\ket{u}$. A control qubit in $\ket{+}$ and the register in $\ket{u}$ pass through a
controlled-$U^{2^j}$ gate. The output state is
\par\smallskip\noindent\begin{tabularx}{\linewidth}{@{}*{2}{L}@{}}
\optlabel{A}\,$\frac{1}{\sqrt2}\,e^{2\pi i\,2^j\varphi}\big(\ket0 + \ket1\big)\ket{u}$ & \correctlabel{B}\,$\frac{1}{\sqrt2}\big(\ket0 + e^{2\pi i\,2^j\varphi}\ket1\big)\ket{u}$ \\[3pt]
\optlabel{C}\,$\frac{1}{\sqrt2}\big(\ket0 + e^{2\pi i\,j\varphi}\ket1\big)\ket{u}$ & \optlabel{D}\,$\frac{1}{\sqrt2}\big(\ket0 + e^{2\pi i\,\varphi/2^j}\ket1\big)\ket{u}$
\end{tabularx}
\end{qitem}
\begin{solbox}{B}
\textbf{\color{brand}What you need to know.} $U^{m}\ket u = e^{2\pi i m\varphi}\ket u$; the controlled gate acts only on the $\ket1$ branch.

\textbf{\color{brand}Step-by-step solution.} $U^{2^j}\ket{u} = \big(e^{2\pi i\varphi}\big)^{2^j}\ket{u} = e^{2\pi i\,2^j\varphi}\ket{u}$. The gate acts only on the $\ket1$ branch:
\[ \frac{1}{\sqrt2}\big(\ket0\ket u + \ket1\,U^{2^j}\ket u\big) = \frac{1}{\sqrt2}\big(\ket0 + e^{2\pi i\,2^j\varphi}\ket1\big)\ket u. \]
The phase is \emph{relative} (only on $\ket1$), not global, so it is measurable after a Hadamard / inverse QFT. If
$\varphi = 0.\varphi_1\varphi_2\dots$ in binary, $2^j\varphi$ shifts the binary point $j$ places and the phase is
$e^{2\pi i\,0.\varphi_{j+1}\varphi_{j+2}\dots}$ --- exactly the factor in the QFT product formula.

\textbf{\color{brand}Why the other options are wrong.}
\begin{itemize}[leftmargin=2.4em, itemsep=1pt, topsep=2pt]
\item[\textbf{(A)}] A phase on both branches is a global phase: it would be unobservable and is not what happens.
\item[\textbf{(C)}] $j\varphi$ multiplies by $j$ instead of $2^j$.
\item[\textbf{(D)}] $\varphi/2^j$ divides instead of multiplying by $2^j$.
\end{itemize}
\textbf{\color{brand}Remember.} $U^{2^j}$ shifts the binary point of $\varphi$ by $j$ places.
\end{solbox}
\qsep

\needspace{13\baselineskip}
\begin{qitem}{29}
\qmeta{Single-correct MCQ}{Post-measurement state}
The state $\frac{1}{\sqrt3}\big(\ket{00} + \ket{01} + \ket{11}\big)$ is prepared and only the \emph{first} qubit is
measured, giving $0$. The probability of this outcome and the post-measurement state of the second qubit are
\par\smallskip\noindent\begin{tabularx}{\linewidth}{@{}*{4}{L}@{}}
\correctlabel{A}\,$\frac23$ and $\ket{+}$ & \optlabel{B}\,$\frac13$ and $\ket{0}$ & \optlabel{C}\,$\frac23$ and $\ket{0}$ & \optlabel{D}\,$\frac13$ and $\ket{+}$
\end{tabularx}
\end{qitem}
\begin{solbox}{A}
\textbf{\color{brand}What you need to know.} Group by the measured qubit, take the weight, renormalise the surviving branch.

\textbf{\color{brand}Step-by-step solution.} Group by the first qubit: $\frac{1}{\sqrt3}\Big(\ket0\big(\ket0 + \ket1\big) + \ket1\ket1\Big)
= \sqrt{\tfrac23}\,\ket0\ket{+} + \sqrt{\tfrac13}\,\ket1\ket1$.
Probability of first qubit $= 0$: $\big|\sqrt{2/3}\big|^2 = \frac23$. By the projection postulate the state collapses to the
normalised $\ket{0}\ket{+}$, so the second qubit is $\ket{+} = \frac{1}{\sqrt2}(\ket0 + \ket1)$.
(Partial measurement keeps the \emph{relative} amplitudes of the surviving branch; this is why in Shor's algorithm measuring
the work register leaves the counting register in a periodic superposition.)

\textbf{\color{brand}Why the other options are wrong.}
\begin{itemize}[leftmargin=2.4em, itemsep=1pt, topsep=2pt]
\item[\textbf{(B)}] Both probability and state are wrong.
\item[\textbf{(C)}] $\ket0$: the surviving branch is $\ket0 + \ket1$, not $\ket0$ alone.
\item[\textbf{(D)}] $\frac13$ is $P(\text{first qubit} = 1)$.
\end{itemize}
\textbf{\color{brand}Remember.} Partial measurement keeps the relative amplitudes of the surviving branch.
\end{solbox}
\qsep

\needspace{13\baselineskip}
\begin{qitem}{30}
\qmeta{Single-correct MCQ}{Unitarity of gates}
Which of the following matrices can \emph{not} be a quantum gate?
\par\smallskip\noindent\begin{tabularx}{\linewidth}{@{}*{4}{L}@{}}
\correctlabel{A}\,$\frac{1}{\sqrt2}\begin{pmatrix}1&1\\ i&i\end{pmatrix}$ & \optlabel{B}\,$\frac{1}{\sqrt2}\begin{pmatrix}1&1\\ 1&-1\end{pmatrix}$ & \optlabel{C}\,$\frac{1}{\sqrt2}\begin{pmatrix}1&i\\ i&1\end{pmatrix}$ & \optlabel{D}\,$\begin{pmatrix}0&1\\ i&0\end{pmatrix}$
\end{tabularx}
\end{qitem}
\begin{solbox}{A}
\textbf{\color{brand}What you need to know.} Gates are unitary: orthonormal columns, $|\det| = 1$.

\textbf{\color{brand}Step-by-step solution.} Closed-system evolution is unitary: $U^\dagger U = I$, i.e.\ the columns are orthonormal.
For $\frac{1}{\sqrt2}\begin{psmallmatrix}1&1\\ i&i\end{psmallmatrix}$ both columns equal $\frac{1}{\sqrt2}(1, i)^T$, so their inner
product is $1 \neq 0$ (also $\det = 0$): it maps $\ket0$ and $\ket1$ to the same state and would destroy information --- not unitary.
Checks: $H$ is unitary; for $\frac{1}{\sqrt2}\begin{psmallmatrix}1&i\\ i&1\end{psmallmatrix}$ the column inner product is
$\frac12(\bar1\cdot i + \bar i\cdot 1) = \frac12(i - i) = 0$ and each column has norm 1; $\begin{psmallmatrix}0&1\\ i&0\end{psmallmatrix}$ has
columns $(0, i)$ and $(1, 0)$, orthonormal. (Unitarity is also why the QFT is invertible: $F^{-1} = F^\dagger$.)

\textbf{\color{brand}Why the other options are wrong.}
\begin{itemize}[leftmargin=2.4em, itemsep=1pt, topsep=2pt]
\item[\textbf{(B)}] $H$ is unitary (orthonormal columns).
\item[\textbf{(C)}] $\frac{1}{\sqrt2}\begin{psmallmatrix}1&i\\i&1\end{psmallmatrix}$: columns orthogonal ($\frac12(i - i) = 0$), norms 1 --- unitary.
\item[\textbf{(D)}] $\begin{psmallmatrix}0&1\\i&0\end{psmallmatrix}$: columns $(0,i)$, $(1,0)$ orthonormal --- unitary.
\end{itemize}
\textbf{\color{brand}Remember.} Equal columns $\Rightarrow$ $\det = 0$ $\Rightarrow$ not a gate.
\end{solbox}
\qsep

\needspace{13\baselineskip}
\begin{qitem}{31}
\qmeta{Single-correct MCQ}{Commutation of gates}
Which two-qubit gate does \emph{not} commute with controlled-$R_2 = \operatorname{diag}(1, 1, 1, i)$?
\par\smallskip\noindent\begin{tabularx}{\linewidth}{@{}*{2}{L}@{}}
\optlabel{A}\,$Z\otimes Z$ & \optlabel{B}\,$S\otimes I$, where $S = \operatorname{diag}(1, i)$ \\[3pt]
\optlabel{C}\,$\mathrm{SWAP}$ & \correctlabel{D}\,$H\otimes I$
\end{tabularx}
\end{qitem}
\begin{solbox}{D}
\textbf{\color{brand}What you need to know.} Diagonal matrices commute; a permutation commutes with a diagonal iff the swapped entries are equal.

\textbf{\color{brand}Step-by-step solution.} $Z\otimes Z$ and $S \otimes I$ are diagonal; diagonal matrices always commute.
SWAP exchanges $\ket{01} \leftrightarrow \ket{10}$; $\mathrm{C}R_2$ has the same entry ($1$) on both, so
$\mathrm{SWAP}\cdot\mathrm{C}R_2\cdot\mathrm{SWAP} = \mathrm{C}R_2$: they commute (controlled-phase gates are symmetric).
$H\otimes I$ mixes $\ket{0x}$ with $\ket{1x}$. For $x = 1$ the diagonal entries of $\mathrm{C}R_2$ on $\ket{01}$ and $\ket{11}$ are
$1$ and $i$ --- different --- so the two operators do not commute. Explicitly
$(H\otimes I)\,\mathrm{C}R_2\ket{01} = \frac{1}{\sqrt2}(\ket{01} + \ket{11})$ but $\mathrm{C}R_2(H\otimes I)\ket{01} = \frac{1}{\sqrt2}(\ket{01} + i\ket{11})$.
This is why the order ``$H$ first, then the controlled rotations'' in the QFT circuit matters.

\textbf{\color{brand}Why the other options are wrong.}
\begin{itemize}[leftmargin=2.4em, itemsep=1pt, topsep=2pt]
\item[\textbf{(A)}] $Z\otimes Z$ is diagonal $\Rightarrow$ commutes.
\item[\textbf{(B)}] $S\otimes I$ is diagonal $\Rightarrow$ commutes.
\item[\textbf{(C)}] SWAP exchanges $\ket{01},\ket{10}$, both with entry 1 $\Rightarrow$ commutes.
\end{itemize}
\textbf{\color{brand}Remember.} Order matters only for gates that mix states with different phases.
\end{solbox}
\qsep

\needspace{13\baselineskip}
\begin{qitem}{32}
\qmeta{Single-correct MCQ}{State-vector memory cost}
A state-vector simulator stores each amplitude as a single-precision complex number (two 4-byte floats).
The memory needed for the full state of $32$ qubits is
\par\smallskip\noindent\begin{tabularx}{\linewidth}{@{}*{4}{L}@{}}
\optlabel{A}\,$256$ bytes & \correctlabel{B}\,$32$ GiB & \optlabel{C}\,$16$ GiB & \optlabel{D}\,$64$ GiB
\end{tabularx}
\end{qitem}
\begin{solbox}{B}
\textbf{\color{brand}What you need to know.} Memory $= 2^n\times$ bytes per amplitude; single-precision complex = 8 bytes.

\textbf{\color{brand}Step-by-step solution.} $n$ qubits $\Rightarrow 2^n$ amplitudes. Each amplitude occupies $2 \times 4 = 8$ bytes. Memory $= 2^{32}\times 8 = 2^{35}$ bytes
$= 2^{5}\times 2^{30}$ bytes $= 32$ GiB. Double precision (16 bytes) would need 64 GiB; 256 bytes ($32\times8$) is the
wrong idea that each qubit needs one number. Every extra qubit doubles the memory --- the reason the course notebooks
default to 15 qubits ($2^{15}\times 8$ B $= 256$ KiB) on a laptop.

\textbf{\color{brand}Why the other options are wrong.}
\begin{itemize}[leftmargin=2.4em, itemsep=1pt, topsep=2pt]
\item[\textbf{(A)}] 256 bytes counts one number per qubit, ignoring superposition.
\item[\textbf{(C)}] 16 GiB uses 4 bytes per amplitude (real only).
\item[\textbf{(D)}] 64 GiB is double precision (16 bytes).
\end{itemize}
\textbf{\color{brand}Remember.} Each extra qubit doubles the memory.
\end{solbox}
\qsep

\needspace{13\baselineskip}
\begin{qitem}{33}
\qmeta{Single-correct MCQ}{Bell state in the $X$ basis}
$H\otimes H$ is applied to the Bell state $\frac{1}{\sqrt2}\big(\ket{00} - \ket{11}\big)$, followed by a measurement of both
qubits. The probability of the outcome $01$ is
\par\smallskip\noindent\begin{tabularx}{\linewidth}{@{}*{4}{L}@{}}
\optlabel{A}\,$0$ & \optlabel{B}\,$\dfrac14$ & \correctlabel{C}\,$\dfrac12$ & \optlabel{D}\,$1$
\end{tabularx}
\end{qitem}
\begin{solbox}{C}
\textbf{\color{brand}What you need to know.} Apply $H\otimes H$ term by term using signs $(-1)^{x\cdot y}$, then add.

\textbf{\color{brand}Step-by-step solution.} $H\otimes H\ket{00} = \frac12(\ket{00} + \ket{01} + \ket{10} + \ket{11})$ and
$H\otimes H\ket{11} = \frac12(\ket{00} - \ket{01} - \ket{10} + \ket{11})$ (sign $(-1)^{x\cdot y}$). Subtract and divide by $\sqrt2$:
\[ \frac{1}{2\sqrt2}\big(2\ket{01} + 2\ket{10}\big) = \frac{1}{\sqrt2}\big(\ket{01} + \ket{10}\big). \]
Hence $P(01) = P(10) = \frac12$ and $P(00) = P(11) = 0$. The perfect ``equal bits'' correlation in the $Z$ basis became a
perfect ``opposite bits'' correlation in the $X$ basis --- interference of four amplitudes, not a product-state result
($\frac14$).

\textbf{\color{brand}Why the other options are wrong.}
\begin{itemize}[leftmargin=2.4em, itemsep=1pt, topsep=2pt]
\item[\textbf{(A)}] $0$ assumes the Bell correlation (equal bits) survives in the new basis.
\item[\textbf{(B)}] $\frac14$ assumes a uniform distribution (product-state thinking).
\item[\textbf{(D)}] $1$ would need all weight on $\ket{01}$; it is shared with $\ket{10}$.
\end{itemize}
\textbf{\color{brand}Remember.} Entangled correlations change form, not strength, in another basis.
\end{solbox}
\qsep

\needspace{13\baselineskip}
\begin{qitem}{34}
\qmeta{Single-correct MCQ}{Mixture vs.\ superposition}
A source emits a qubit that is $\ket0$ with probability $\frac12$ and $\ket1$ with probability $\frac12$ (a classical
coin decides; nobody knows which). A Hadamard gate is applied and the qubit is measured. The probability of outcome $0$ is
\par\smallskip\noindent\begin{tabularx}{\linewidth}{@{}*{4}{L}@{}}
\optlabel{A}\,$1$ & \correctlabel{B}\,$\dfrac12$ & \optlabel{C}\,$\dfrac14$ & \optlabel{D}\,$0$
\end{tabularx}
\end{qitem}
\begin{solbox}{B}
\textbf{\color{brand}What you need to know.} A mixture is averaged over cases; amplitudes of different cases never interfere.

\textbf{\color{brand}Step-by-step solution.} This is a \emph{mixture}, not a superposition. Treat the two cases separately and average:
$H\ket0 = \ket{+}$ gives $P(0) = \frac12$; $H\ket1 = \ket{-}$ gives $P(0) = \frac12$. Hence $P(0) = \frac12\cdot\frac12 + \frac12\cdot\frac12 = \frac12$.
(Density matrix: $\rho = \frac12 I$, and $H\rho H^\dagger = \frac12 I$.) Contrast: the superposition
$\ket{+} = \frac{1}{\sqrt2}(\ket0 + \ket1)$ gives $H\ket{+} = \ket0$ and $P(0) = 1$. Amplitudes interfere, classical probabilities do
not --- QFT peaks need coherence, which is why decoherence (dephasing) destroys them.

\textbf{\color{brand}Why the other options are wrong.}
\begin{itemize}[leftmargin=2.4em, itemsep=1pt, topsep=2pt]
\item[\textbf{(A)}] $1$ treats the mixture as the superposition $\ket+$.
\item[\textbf{(C)}] $\frac14$ multiplies probabilities instead of averaging.
\item[\textbf{(D)}] $0$ treats it as $\ket-$.
\end{itemize}
\textbf{\color{brand}Remember.} Interference requires coherence; decoherence kills QFT peaks.
\end{solbox}
\qsep

\needspace{13\baselineskip}
\begin{qitem}{35}
\qmeta{Matching lists}{Single-qubit gate actions}
Match each single-qubit gate in List-I with its action in List-II. ($R_k = \operatorname{diag}(1, e^{2\pi i/2^k})$.)
\par\smallskip\noindent{\renewcommand{\arraystretch}{1.35}%
\begin{tabularx}{\linewidth}{@{}c>{\raggedright\arraybackslash}p{0.33\linewidth}@{\hspace{1.2em}}cL@{}}
\toprule
\multicolumn{2}{@{}l}{\textbf{List-I}} & \multicolumn{2}{l@{}}{\textbf{List-II}} \\ \midrule
\textbf{(P)} & $X$ & \textbf{(1)} & $\ket{1} \mapsto \frac{1}{\sqrt2}\big(\ket0 - \ket1\big)$ \\
\textbf{(Q)} & $H$ & \textbf{(2)} & $\ket{1} \mapsto e^{i\pi/4}\ket{1}$ \\
\textbf{(R)} & $R_2$ & \textbf{(3)} & $\ket{0} \leftrightarrow \ket{1}$ \\
\textbf{(S)} & $R_3$ & \textbf{(4)} & $\ket{1} \mapsto i\ket{1}$ \\
\bottomrule
\end{tabularx}}
\par\smallskip\noindent\begin{tabularx}{\linewidth}{@{}*{2}{L}@{}}
\optlabel{A}\,P$\to$3, Q$\to$1, R$\to$2, S$\to$4 & \optlabel{B}\,P$\to$1, Q$\to$3, R$\to$4, S$\to$2 \\[3pt]
\correctlabel{C}\,P$\to$3, Q$\to$1, R$\to$4, S$\to$2 & \optlabel{D}\,P$\to$4, Q$\to$1, R$\to$3, S$\to$2
\end{tabularx}
\end{qitem}
\begin{solbox}{C}
\textbf{\color{brand}What you need to know.} $X$ flips, $H\ket1 = \ket-$, $R_2 = S$ ($\ket1\to i\ket1$), $R_3 = T$ ($\ket1\to e^{i\pi/4}\ket1$).

\textbf{\color{brand}Step-by-step solution.} $X$ is the bit flip: $\ket0 \leftrightarrow \ket1$ (3). $H\ket1 = \frac{1}{\sqrt2}(\ket0 - \ket1) = \ket{-}$ (1).
$R_2 = \operatorname{diag}(1, e^{2\pi i/4}) = \operatorname{diag}(1, i) = S$, so $\ket1 \mapsto i\ket1$ (4).
$R_3 = \operatorname{diag}(1, e^{2\pi i/8}) = T$, so $\ket1 \mapsto e^{i\pi/4}\ket1$ (2).
Hence P$\to$3, Q$\to$1, R$\to$4, S$\to$2. Note $R_1 = Z$, $R_2 = S$, $R_3 = T$, and $R_{k}^{2} = R_{k-1}$.

\textbf{\color{brand}Why the other options are wrong.}
\begin{itemize}[leftmargin=2.4em, itemsep=1pt, topsep=2pt]
\item[\textbf{(A)}] R$\to$2, S$\to$4 swaps $S$ and $T$: $R_2$ is the quarter turn $i$.
\item[\textbf{(B)}] P$\to$1, Q$\to$3 swaps $X$ and $H$.
\item[\textbf{(D)}] P$\to$4 says $X$ is a phase gate.
\end{itemize}
\textbf{\color{brand}Remember.} $R_k$ is a rotation by $360^\circ/2^k$.
\end{solbox}
\qsep

\needspace{13\baselineskip}
\begin{qitem}{36}
\qmeta{Matching lists}{Two-qubit correlations}
Match each two-qubit state in List-I with its statistics when \emph{both} qubits are measured in the computational
basis (List-II).
\par\smallskip\noindent{\renewcommand{\arraystretch}{1.35}%
\begin{tabularx}{\linewidth}{@{}c>{\raggedright\arraybackslash}p{0.33\linewidth}@{\hspace{1.2em}}cL@{}}
\toprule
\multicolumn{2}{@{}l}{\textbf{List-I}} & \multicolumn{2}{l@{}}{\textbf{List-II}} \\ \midrule
\textbf{(P)} & $\frac{1}{\sqrt2}\big(\ket{00} + \ket{11}\big)$ & \textbf{(1)} & all four outcomes, each with probability $\frac14$ \\
\textbf{(Q)} & $\frac{1}{\sqrt2}\big(\ket{01} - \ket{10}\big)$ & \textbf{(2)} & the two bits are always equal; $00$ and $11$ each $\frac12$ \\
\textbf{(R)} & $\frac12\big(\ket0 + i\ket1\big)\big(\ket0 - \ket1\big)$ & \textbf{(3)} & first bit always $0$; second bit $0$ or $1$ with $\frac12$ each \\
\textbf{(S)} & $\ket{0}\ket{+}$ & \textbf{(4)} & the two bits always differ; $01$ and $10$ each $\frac12$ \\
\bottomrule
\end{tabularx}}
\par\smallskip\noindent\begin{tabularx}{\linewidth}{@{}*{2}{L}@{}}
\optlabel{A}\,P$\to$1, Q$\to$4, R$\to$2, S$\to$3 & \optlabel{B}\,P$\to$4, Q$\to$2, R$\to$1, S$\to$3 \\[3pt]
\optlabel{C}\,P$\to$2, Q$\to$4, R$\to$3, S$\to$1 & \correctlabel{D}\,P$\to$2, Q$\to$4, R$\to$1, S$\to$3
\end{tabularx}
\end{qitem}
\begin{solbox}{D}
\textbf{\color{brand}What you need to know.} Born rule on each listed state; phases do not affect computational-basis statistics.

\textbf{\color{brand}Step-by-step solution.} P ($\Phi^+$): amplitudes only on $\ket{00}$, $\ket{11}$, each $\frac{1}{\sqrt2}$ $\Rightarrow$ equal bits, $\frac12$ each (2).
Q ($\Psi^-$): only $\ket{01}$, $\ket{10}$ $\Rightarrow$ bits always differ (4); the minus sign is invisible in this basis.
R: expand $\frac12(\ket{00} - \ket{01} + i\ket{10} - i\ket{11})$; every amplitude has modulus $\frac12$ $\Rightarrow$ all four outcomes
with probability $\frac14$ (1). The phases $i$ and $-1$ do not show up --- just like the output of a QFT on a basis state.
S: $\ket0\ket{+} = \frac{1}{\sqrt2}(\ket{00} + \ket{01})$ $\Rightarrow$ first bit 0, second random (3).
Hence P$\to$2, Q$\to$4, R$\to$1, S$\to$3.

\textbf{\color{brand}Why the other options are wrong.}
\begin{itemize}[leftmargin=2.4em, itemsep=1pt, topsep=2pt]
\item[\textbf{(A)}] P$\to$1 says $\Phi^+$ gives all four outcomes; it has only $00$ and $11$.
\item[\textbf{(B)}] P$\to$4, Q$\to$2 swaps $\Phi^+$ (equal bits) and $\Psi^-$ (different bits).
\item[\textbf{(C)}] R$\to$3, S$\to$1 swaps the uniform product and $\ket0\ket+$.
\end{itemize}
\textbf{\color{brand}Remember.} Read off which basis states have non-zero amplitude.
\end{solbox}
\qsep

\needspace{13\baselineskip}
\begin{qitem}{37}
\qmeta{Assertion--Reason}{Distinguishability; relative phase}
\textbf{Assertion (A):} The states $\ket{+}$ and $\ket{-}$ cannot be distinguished by a measurement in the computational
basis $\{\ket0, \ket1\}$.\\[1mm]
\textbf{Reason (R):} $\ket{+}$ and $\ket{-}$ differ only by a global phase.
\begin{itemize}[leftmargin=2.6em, labelsep=0.4em, itemsep=2pt, topsep=3pt]
\item[\optlabel{A}] Both Assertion (A) and Reason (R) are true, and R is the correct explanation of A
\item[\optlabel{B}] Both A and R are true, but R is \emph{not} the correct explanation of A
\item[\correctlabel{C}] A is true, but R is false
\item[\optlabel{D}] A is false, but R is true
\end{itemize}
\end{qitem}
\begin{solbox}{C}
\textbf{\color{brand}What you need to know.} $\ket\pm$ differ by a relative phase; same statistics in $Z$, opposite in $X$.

\textbf{\color{brand}Step-by-step solution.} \textbf{A is true:} for both states $P(0) = P(1) = \frac12$, because $|\pm\frac{1}{\sqrt2}|^2 = \frac12$.
\textbf{R is false:} $\ket{-} = \frac{1}{\sqrt2}(\ket0 - \ket1)$ is not $e^{i\gamma}\ket{+}$ for any $\gamma$; the two differ by a
\emph{relative} phase $-1$ on $\ket1$. They are orthogonal ($\braket{+|-} = 0$) and therefore perfectly distinguishable in the
$X$ basis (apply $H$, then measure: $\ket{+}\to\ket0$, $\ket-\to\ket1$). States that differ by a global phase could never be
distinguished in \emph{any} basis. Answer: (C).

\textbf{\color{brand}Why the other options are wrong.}
\begin{itemize}[leftmargin=2.4em, itemsep=1pt, topsep=2pt]
\item[\textbf{(A)}] needs A true, R true \emph{and} R explaining A; R is false.
\item[\textbf{(B)}] needs A and R both true but R not explaining A; R is false.
\item[\textbf{(D)}] needs A false and R true; A is true, R is false.
\end{itemize}
\textbf{\color{brand}Remember.} Indistinguishable in one basis $\neq$ identical.
\end{solbox}
\qsep

\needspace{13\baselineskip}
\begin{qitem}{38}
\qmeta{Statement I / II}{Product-state criterion}
\textbf{Statement I:} The Bell state $\frac{1}{\sqrt2}\big(\ket{00} + \ket{11}\big)$ can be written as $\ket{a}\otimes\ket{b}$ for
suitable single-qubit states $\ket a$, $\ket b$.\\[1mm]
\textbf{Statement II:} A two-qubit state $c_{00}\ket{00} + c_{01}\ket{01} + c_{10}\ket{10} + c_{11}\ket{11}$ is a product state if and only
if $c_{00}c_{11} = c_{01}c_{10}$.
\begin{itemize}[leftmargin=2.6em, labelsep=0.4em, itemsep=2pt, topsep=3pt]
\item[\optlabel{A}] Both Statement I and Statement II are true
\item[\optlabel{B}] Both Statement I and Statement II are false
\item[\optlabel{C}] Statement I is true, but Statement II is false
\item[\correctlabel{D}] Statement I is false, but Statement II is true
\end{itemize}
\end{qitem}
\begin{solbox}{D}
\textbf{\color{brand}What you need to know.} Rank-1 coefficient matrix $\Leftrightarrow$ product state.

\textbf{\color{brand}Step-by-step solution.} \textbf{II is true.} ($\Rightarrow$) If the state is $(a_0\ket0 + a_1\ket1)(b_0\ket0 + b_1\ket1)$ then $c_{jk} = a_jb_k$ and
$c_{00}c_{11} = a_0b_0a_1b_1 = c_{01}c_{10}$. ($\Leftarrow$) If $c_{00}c_{11} = c_{01}c_{10}$, the matrix $C = \begin{psmallmatrix}c_{00}&c_{01}\\ c_{10}&c_{11}\end{psmallmatrix}$
has $\det C = 0$, so it has rank 1 and $C = \mathbf{a}\mathbf{b}^T$, i.e.\ $c_{jk} = a_jb_k$.
\textbf{I is false:} for $\Phi^+$, $c_{00}c_{11} = \frac12 \neq 0 = c_{01}c_{10}$. Answer: (D).

\textbf{\color{brand}Why the other options are wrong.}
\begin{itemize}[leftmargin=2.4em, itemsep=1pt, topsep=2pt]
\item[\textbf{(A)}] needs both statements true; I is false.
\item[\textbf{(B)}] needs both statements false; II is true.
\item[\textbf{(C)}] needs I true and II false; the truth values are the other way round.
\end{itemize}
\textbf{\color{brand}Remember.} $\det C = 0$ is the two-qubit entanglement test.
\end{solbox}
\qsep

\needspace{20\baselineskip}
\begin{passagebox}{Paragraph for Questions 39 and 40}
\textbf{Hadamard test.} A control qubit starts in $\ket{0}$ and a target register starts in an eigenstate $\ket{u}$
of a unitary $U$ with $U\ket{u} = e^{2\pi i\varphi}\ket{u}$, $0 \le \varphi < 1$. The circuit applies, in order:
(i) $H$ on the control; (ii) controlled-$U$ (control qubit $\to$ target register); (iii) $H$ on the control;
(iv) a measurement of the control qubit in the computational basis. This is the one-counting-qubit version of quantum
phase estimation.
\end{passagebox}
\needspace{13\baselineskip}
\begin{qitem}{39}
\qmeta{Paragraph based}{Hadamard test}
The probability that the control qubit is found in state $\ket1$ is
\par\smallskip\noindent\begin{tabularx}{\linewidth}{@{}*{4}{L}@{}}
\optlabel{A}\,$\cos^2(\pi\varphi)$ & \optlabel{B}\,$\sin^2(\pi\varphi/2)$ & \optlabel{C}\,$\sin^2(2\pi\varphi)$ & \correctlabel{D}\,$\sin^2(\pi\varphi)$
\end{tabularx}
\end{qitem}
\begin{solbox}{D}
\textbf{\color{brand}What you need to know.} Hadamard test = $H$, kickback, $H$: $P(1) = \sin^2(\pi\varphi)$.

\textbf{\color{brand}Step-by-step solution.} After (i): $\frac{1}{\sqrt2}(\ket0 + \ket1)\ket u$. After (ii), by phase kickback (Q.28 with $j = 0$):
$\frac{1}{\sqrt2}\big(\ket0 + e^{2\pi i\varphi}\ket1\big)\ket u$. After (iii) (Q.23 with $\phi = 2\pi\varphi$):
\[ \Big(\frac{1 + e^{2\pi i\varphi}}{2}\ket0 + \frac{1 - e^{2\pi i\varphi}}{2}\ket1\Big)\ket u . \]
$|1 - e^{2\pi i\varphi}|^2 = 2 - 2\cos(2\pi\varphi) = 4\sin^2(\pi\varphi)$, so $P(1) = \sin^2(\pi\varphi)$ and $P(0) = \cos^2(\pi\varphi)$.
$\varphi = 0 \Rightarrow$ always 0; $\varphi = \frac12 \Rightarrow$ always 1: one qubit resolves exactly one binary digit
$\varphi = 0.\varphi_1$. More digits need more counting qubits and an inverse QFT.

\textbf{\color{brand}Why the other options are wrong.}
\begin{itemize}[leftmargin=2.4em, itemsep=1pt, topsep=2pt]
\item[\textbf{(A)}] $\cos^2(\pi\varphi)$ is $P(0)$.
\item[\textbf{(B)}] $\sin^2(\pi\varphi/2)$ halves the angle twice.
\item[\textbf{(C)}] $\sin^2(2\pi\varphi)$ uses the full angle instead of the half angle.
\end{itemize}
\textbf{\color{brand}Remember.} One counting qubit resolves one binary digit of $\varphi$.
\end{solbox}
\qsep

\needspace{13\baselineskip}
\begin{qitem}{40}
\qmeta{Paragraph based}{Repetition and confidence}
Take $\varphi = \frac13$. The circuit is run repeatedly (fresh $\ket0$ control each time). The minimum number of runs
needed so that the probability of observing outcome $1$ \emph{at least once} exceeds $0.99$ is
\par\smallskip\noindent\begin{tabularx}{\linewidth}{@{}*{4}{L}@{}}
\optlabel{A}\,$2$ & \optlabel{B}\,$3$ & \correctlabel{C}\,$4$ & \optlabel{D}\,$7$
\end{tabularx}
\end{qitem}
\begin{solbox}{C}
\textbf{\color{brand}What you need to know.} Independent repetitions: $P(\text{at least one success}) = 1 - (1-p)^k$.

\textbf{\color{brand}Step-by-step solution.} From Q.39, $P(1) = \sin^2(\pi/3) = \left(\frac{\sqrt3}{2}\right)^2 = \frac34$, so a single run gives 0 with probability $\frac14$.
Runs are independent: $P(\text{no } 1 \text{ in } k \text{ runs}) = \left(\frac14\right)^k$. We need $1 - 4^{-k} > 0.99$, i.e.\ $4^k > 100$.
$4^3 = 64 < 100$ (failure probability $0.0156$), $4^4 = 256 > 100$ (failure probability $0.0039$). Hence $k = 4$.
(Repetition amplifies confidence; it does not make the single-shot estimate of $\varphi = \frac13 = 0.0101\ldots_2$ exact.)

\textbf{\color{brand}Why the other options are wrong.}
\begin{itemize}[leftmargin=2.4em, itemsep=1pt, topsep=2pt]
\item[\textbf{(A)}] $2$ runs give failure $0.0625$.
\item[\textbf{(B)}] $3$ runs give failure $4^{-3} = 0.0156 > 0.01$.
\item[\textbf{(D)}] $7$ is the answer for $p = \frac12$, not $p = \frac34$.
\end{itemize}
\textbf{\color{brand}Remember.} Find $p$ first, then the smallest $k$ with $(1-p)^k < \delta$.
\end{solbox}
\qsep

\newpage

\begin{sectionbanner}SECTION C \enspace$\cdot$\enspace QFT Basics \hfill Q.41--60\end{sectionbanner}

\needspace{13\baselineskip}
\begin{qitem}{41}
\qmeta{Single-correct MCQ}{QFT matrix element ($n = 2$)}
For $n = 2$, the amplitude of $\ket{10}$ in $\mathrm{QFT}\ket{11}$ is
\par\smallskip\noindent\begin{tabularx}{\linewidth}{@{}*{4}{L}@{}}
\correctlabel{A}\,$-\dfrac12$ & \optlabel{B}\,$\dfrac12$ & \optlabel{C}\,$-\dfrac{i}{2}$ & \optlabel{D}\,$\dfrac{i}{2}$
\end{tabularx}
\end{qitem}
\begin{solbox}{A}
\textbf{\color{brand}What you need to know.} Amplitude $\braket{y|\QFTop|x} = \omega^{xy}/\sqrt N$; convert bit strings to integers first.

\textbf{\color{brand}Step-by-step solution.} $N = 4$, $\omega = e^{2\pi i/4} = i$, $x = \ket{11} = 3$, $y = \ket{10} = 2$. From the definition
\[ \braket{y|\mathrm{QFT}|x} = \frac{1}{\sqrt N}\,\omega^{xy} = \frac12\, i^{6} = \frac12\,(i^4)(i^2) = -\frac12 . \]
Full output: $\mathrm{QFT}\ket{3} = \frac12\big(\ket0 - i\ket1 - \ket2 + i\ket3\big)$, i.e.\ $\frac12(\ket{00} - i\ket{01} - \ket{10} + i\ket{11})$.

\textbf{\color{brand}Why the other options are wrong.}
\begin{itemize}[leftmargin=2.4em, itemsep=1pt, topsep=2pt]
\item[\textbf{(B)}] $\frac12$ drops the sign of $i^6 = -1$.
\item[\textbf{(C)}] $-\frac{i}{2}$ uses $xy = 3$ (wrong integers).
\item[\textbf{(D)}] $\frac{i}{2}$ uses $i^{xy}$ with $xy = 1$ or 5 (wrong integers).
\end{itemize}
\textbf{\color{brand}Remember.} $\ket{10} = 2$, $\ket{11} = 3$; reduce $xy$ mod $N$.
\end{solbox}
\qsep

\needspace{13\baselineskip}
\begin{qitem}{42}
\qmeta{Single-correct MCQ}{Output distribution of $\mathrm{QFT}\ket{x}$}
A 5-qubit QFT is applied to the basis state $\ket{x}$ with $x = 19$ and the output is measured. The probability of
obtaining $y = 7$ is
\par\smallskip\noindent\begin{tabularx}{\linewidth}{@{}*{4}{L}@{}}
\optlabel{A}\,$\dfrac{7}{32}$ & \optlabel{B}\,$\dfrac{1}{\sqrt{32}}$ & \optlabel{C}\,$0$ & \correctlabel{D}\,$\dfrac{1}{32}$
\end{tabularx}
\end{qitem}
\begin{solbox}{D}
\textbf{\color{brand}What you need to know.} $|\omega^{xy}| = 1$, so every outcome has probability $1/N$ for a basis input.

\textbf{\color{brand}Step-by-step solution.} $\mathrm{QFT}\ket{x} = \frac{1}{\sqrt{32}}\sum_y \omega^{xy}\ket y$ and $|\omega^{xy}| = 1$, so every outcome has probability
$\left|\frac{1}{\sqrt{32}}\right|^2 = \frac{1}{32}$, whatever $x$ and $y$ are. ($\frac{1}{\sqrt{32}}$ is the modulus of the amplitude,
not a probability.) All the information about $x$ sits in the phases $\omega^{19y}$.

\textbf{\color{brand}Why the other options are wrong.}
\begin{itemize}[leftmargin=2.4em, itemsep=1pt, topsep=2pt]
\item[\textbf{(A)}] $\frac{7}{32}$ invents a dependence on $y$.
\item[\textbf{(B)}] $\frac{1}{\sqrt{32}}$ is the modulus of the amplitude, not the probability.
\item[\textbf{(C)}] $0$: no outcome is forbidden for a basis input.
\end{itemize}
\textbf{\color{brand}Remember.} QFT of a basis state = flat distribution, information in phases.
\end{solbox}
\qsep

\needspace{13\baselineskip}
\begin{qitem}{43}
\qmeta{Single-correct MCQ}{Product formula ($n = 3$)}
For $n = 3$, $\mathrm{QFT}\ket{110}$ equals (qubit 1 written first)
\par\smallskip\noindent\begin{tabularx}{\linewidth}{@{}*{2}{L}@{}}
\optlabel{A}\,$\frac{1}{\sqrt8}\big(\ket0 - \ket1\big)\big(\ket0 - \ket1\big)\big(\ket0 - i\ket1\big)$ & \optlabel{B}\,$\frac{1}{\sqrt8}\big(\ket0 + \ket1\big)\big(\ket0 - \ket1\big)\big(\ket0 + i\ket1\big)$ \\[3pt]
\correctlabel{C}\,$\frac{1}{\sqrt8}\big(\ket0 + \ket1\big)\big(\ket0 - \ket1\big)\big(\ket0 - i\ket1\big)$ & \optlabel{D}\,$\frac{1}{\sqrt8}\big(\ket0 - i\ket1\big)\big(\ket0 - \ket1\big)\big(\ket0 + \ket1\big)$
\end{tabularx}
\end{qitem}
\begin{solbox}{C}
\textbf{\color{brand}What you need to know.} Product formula: output qubit $l$ carries $0.x_{n-l+1}\dots x_n$.

\textbf{\color{brand}Step-by-step solution.} \emph{Product formula.} Write $y = \sum_{l=1}^{n} y_l 2^{\,n-l}$; then $\omega^{xy} = \prod_{l}e^{2\pi i\,x y_l/2^{l}}$ and
$e^{2\pi i\,x/2^l} = e^{2\pi i\,0.x_{n-l+1}\dots x_n}$ (only the last $l$ bits of $x$ survive), so the sum over $y$ factorises into one
factor per output qubit. For $n = 3$: $\mathrm{QFT}\ket{x_1x_2x_3} = \frac{1}{\sqrt8}
\big(\ket0 + e^{2\pi i\,0.x_3}\ket1\big)\big(\ket0 + e^{2\pi i\,0.x_2x_3}\ket1\big)\big(\ket0 + e^{2\pi i\,0.x_1x_2x_3}\ket1\big)$.
With $x_1x_2x_3 = 110$: $0.x_3 = 0.0 = 0 \Rightarrow e^{0} = 1$; $\;0.x_2x_3 = 0.10_2 = \frac12 \Rightarrow e^{i\pi} = -1$;
$\;0.x_1x_2x_3 = 0.110_2 = \frac34 \Rightarrow e^{3\pi i/2} = -i$. Hence
$\frac{1}{\sqrt8}(\ket0 + \ket1)(\ket0 - \ket1)(\ket0 - i\ket1)$.
\emph{Check} with the definition: amplitude of $y$ is $\omega^{6y}/\sqrt8 = (-i)^{y}/\sqrt8$; for $y = 4y_1 + 2y_2 + y_3$,
$(-i)^{4y_1}(-i)^{2y_2}(-i)^{y_3} = 1\cdot(-1)^{y_2}(-i)^{y_3}$ --- the same three factors. The option
$\frac{1}{\sqrt8}(\ket0 - i\ket1)(\ket0 - \ket1)(\ket0 + \ket1)$ is the bit-reversed order (what you get without the final SWAPs).

\textbf{\color{brand}Why the other options are wrong.}
\begin{itemize}[leftmargin=2.4em, itemsep=1pt, topsep=2pt]
\item[\textbf{(A)}] $(\ket0 - \ket1)$ in the first factor uses $0.x_3 = \frac12$, but $x_3 = 0$.
\item[\textbf{(B)}] $(\ket0 + i\ket1)$ in the last factor uses the wrong sign ($0.110 = \frac34$ gives $-i$).
\item[\textbf{(D)}] Reversed factor order: the state before the final SWAPs.
\end{itemize}
\textbf{\color{brand}Remember.} Check one amplitude with the definition after using the product formula.
\end{solbox}
\qsep

\needspace{13\baselineskip}
\begin{qitem}{44}
\qmeta{Single-correct MCQ}{Controlled-rotation count}
The textbook QFT circuit on $n = 12$ qubits contains how many controlled-$R_k$ gates?
\par\smallskip\noindent\begin{tabularx}{\linewidth}{@{}*{4}{L}@{}}
\optlabel{A}\,$144$ & \optlabel{B}\,$72$ & \correctlabel{C}\,$66$ & \optlabel{D}\,$78$
\end{tabularx}
\end{qitem}
\begin{solbox}{C}
\textbf{\color{brand}What you need to know.} Qubit $j$ receives $n - j$ controlled rotations; total $\binom n2$.

\textbf{\color{brand}Step-by-step solution.} Qubit $j$ receives one controlled rotation from each qubit below it, i.e.\ $n - j$ of them. Total
\[ \sum_{j=1}^{n}(n - j) = (n-1) + (n-2) + \dots + 0 = \frac{n(n-1)}{2} = \frac{12\cdot11}{2} = 66 . \]
($78 = \frac{n(n+1)}{2}$ also counts the 12 Hadamards; $144 = n^2$.)

\textbf{\color{brand}Why the other options are wrong.}
\begin{itemize}[leftmargin=2.4em, itemsep=1pt, topsep=2pt]
\item[\textbf{(A)}] $144 = n^2$ overcounts.
\item[\textbf{(B)}] $72$ is not a formula of the circuit.
\item[\textbf{(D)}] $78 = \frac{n(n+1)}{2}$ also counts the $n$ Hadamards.
\end{itemize}
\textbf{\color{brand}Remember.} $\frac{n(n-1)}{2}$ controlled rotations.
\end{solbox}
\qsep

\needspace{13\baselineskip}
\begin{qitem}{45}
\qmeta{Single-correct MCQ}{Two-qubit gates incl.\ SWAPs}
On hardware a SWAP is implemented with 3 CNOTs. Counting every controlled-$R_k$ as one two-qubit gate, the total number
of two-qubit gates in the 12-qubit textbook QFT (including the final SWAPs) is
\par\smallskip\noindent\begin{tabularx}{\linewidth}{@{}*{4}{L}@{}}
\optlabel{A}\,$72$ & \correctlabel{B}\,$84$ & \optlabel{C}\,$102$ & \optlabel{D}\,$66$
\end{tabularx}
\end{qitem}
\begin{solbox}{B}
\textbf{\color{brand}What you need to know.} Two-qubit gates = controlled rotations + 3 per SWAP; $\lfloor n/2\rfloor$ SWAPs.

\textbf{\color{brand}Step-by-step solution.} Controlled rotations: $\frac{12\cdot11}{2} = 66$ (Q.44). SWAPs: $\lfloor 12/2\rfloor = 6$, each costing 3 CNOTs $= 18$.
Total $66 + 18 = 84$. ($72 = 66 + 6$ counts each SWAP as one gate; $102 = 66 + 36$ uses 6 CNOTs per SWAP.)
In practice the SWAPs are often removed by simply relabelling the output qubits, leaving 66.

\textbf{\color{brand}Why the other options are wrong.}
\begin{itemize}[leftmargin=2.4em, itemsep=1pt, topsep=2pt]
\item[\textbf{(A)}] $72$ counts each SWAP as one gate.
\item[\textbf{(C)}] $102$ uses 6 CNOTs per SWAP or 12 SWAPs.
\item[\textbf{(D)}] $66$ forgets the SWAPs.
\end{itemize}
\textbf{\color{brand}Remember.} $\binom n2 + 3\lfloor n/2\rfloor$.
\end{solbox}
\qsep

\needspace{13\baselineskip}
\begin{qitem}{46}
\qmeta{Single-correct MCQ}{Smallest rotation angle}
The smallest rotation angle appearing in the textbook QFT circuit on $6$ qubits is
\par\smallskip\noindent\begin{tabularx}{\linewidth}{@{}*{4}{L}@{}}
\optlabel{A}\,$11.25^\circ$ & \optlabel{B}\,$6^\circ$ & \optlabel{C}\,$2.8125^\circ$ & \correctlabel{D}\,$5.625^\circ$
\end{tabularx}
\end{qitem}
\begin{solbox}{D}
\textbf{\color{brand}What you need to know.} The smallest rotation is $R_n$ with angle $360^\circ/2^n$.

\textbf{\color{brand}Step-by-step solution.} The rotations used are $R_2, \dots, R_n$ with angle $\frac{2\pi}{2^k}$. The smallest is $R_n = R_6$:
$\frac{360^\circ}{2^6} = \frac{360^\circ}{64} = 5.625^\circ$. ($11.25^\circ$ is $R_5$; $2.8125^\circ$ would be $R_7$, which needs 7 qubits.)
These tiny angles are why the approximate QFT drops the smallest rotations (Section D).

\textbf{\color{brand}Why the other options are wrong.}
\begin{itemize}[leftmargin=2.4em, itemsep=1pt, topsep=2pt]
\item[\textbf{(A)}] $11.25^\circ$ is $R_5$.
\item[\textbf{(B)}] $6^\circ$ divides $360^\circ$ by $n\cdot10$ --- not a QFT angle.
\item[\textbf{(C)}] $2.8125^\circ$ is $R_7$, which needs 7 qubits.
\end{itemize}
\textbf{\color{brand}Remember.} Angles halve with each extra qubit.
\end{solbox}
\qsep

\needspace{13\baselineskip}
\begin{qitem}{47}
\qmeta{Single-correct MCQ}{Phase on qubit 1; role of SWAPs}
In the textbook QFT circuit acting on $\ket{x_1x_2\dots x_n}$, \emph{before} the final SWAPs, qubit 1 is in the state
$\frac{1}{\sqrt2}\big(\ket0 + e^{2\pi i\,\phi}\ket1\big)$ with $\phi$ equal to
\par\smallskip\noindent\begin{tabularx}{\linewidth}{@{}*{4}{L}@{}}
\correctlabel{A}\,$0.x_1x_2\dots x_n$ & \optlabel{B}\,$0.x_nx_{n-1}\dots x_1$ & \optlabel{C}\,$0.x_n$ & \optlabel{D}\,$0.x_1$
\end{tabularx}
\end{qitem}
\begin{solbox}{A}
\textbf{\color{brand}What you need to know.} $H$ gives $0.x_1$; each controlled $R_k$ appends bit $x_k$.

\textbf{\color{brand}Step-by-step solution.} $H$ on qubit 1 gives $\frac{1}{\sqrt2}(\ket0 + e^{2\pi i\,0.x_1}\ket1)$ because $(-1)^{x_1} = e^{2\pi i\,x_1/2}$.
Each controlled-$R_k$ from qubit $k$ (for $k = 2, \dots, n$) multiplies the $\ket1$ part by $e^{2\pi i\,x_k/2^k}$, appending the
binary digit $x_k$ at position $k$: $0.x_1 \to 0.x_1x_2 \to \dots \to 0.x_1x_2\dots x_n$.
This factor belongs to the \emph{last} output qubit in the product formula, so the SWAPs move it from qubit 1 to qubit $n$.

\textbf{\color{brand}Why the other options are wrong.}
\begin{itemize}[leftmargin=2.4em, itemsep=1pt, topsep=2pt]
\item[\textbf{(B)}] $0.x_n\dots x_1$ reads the bits in reverse.
\item[\textbf{(C)}] $0.x_n$ is the phase of qubit $n$ before the SWAPs.
\item[\textbf{(D)}] $0.x_1$ is the phase right after $H$, before the rotations.
\end{itemize}
\textbf{\color{brand}Remember.} Before the SWAPs, qubit $j$ holds $0.x_j\dots x_n$.
\end{solbox}
\qsep

\needspace{13\baselineskip}
\begin{qitem}{48}
\qmeta{Single-correct MCQ}{$\mathrm{QFT}^2$ = negation}
For any $n$-qubit basis state, $\mathrm{QFT}^2\ket{x}$ equals ($N = 2^n$)
\par\smallskip\noindent\begin{tabularx}{\linewidth}{@{}*{4}{L}@{}}
\optlabel{A}\,$\ket{x}$ & \optlabel{B}\,$\ket{2x \bmod N}$ & \correctlabel{C}\,$\ket{(-x) \bmod N}$ & \optlabel{D}\,$\ket{N - 1 - x}$
\end{tabularx}
\end{qitem}
\begin{solbox}{C}
\textbf{\color{brand}What you need to know.} $\QFTop^2\ket x = \frac1N\sum_z\sum_y\omega^{y(x+z)}\ket z$; use orthogonality.

\textbf{\color{brand}Step-by-step solution.} \[ \mathrm{QFT}^2\ket x = \frac{1}{N}\sum_{z}\Big(\sum_y \omega^{y(x+z)}\Big)\ket z . \]
The inner geometric sum equals $N$ if $x + z \equiv 0 \pmod N$ and $0$ otherwise. Hence $\mathrm{QFT}^2\ket x = \ket{-x \bmod N}$.
Examples ($n = 2$): $\ket0 \to \ket0$, $\ket1 \to \ket3$, $\ket2 \to \ket2$. Consequently $\mathrm{QFT}^4 = I$.
($\ket{N-1-x}$ is the bitwise complement, a different map: it sends $\ket0 \to \ket{N-1}$.)

\textbf{\color{brand}Why the other options are wrong.}
\begin{itemize}[leftmargin=2.4em, itemsep=1pt, topsep=2pt]
\item[\textbf{(A)}] $\ket x$ would mean $\QFTop^2 = I$; that is only true for $\QFTop^4$.
\item[\textbf{(B)}] $\ket{2x}$ adds the exponents wrongly.
\item[\textbf{(D)}] $\ket{N-1-x}$ is the bit complement, a different map ($0\to N-1$).
\end{itemize}
\textbf{\color{brand}Remember.} $\QFTop^2$ is the reflection $x\mapsto -x$; $\QFTop^4 = I$.
\end{solbox}
\qsep

\needspace{13\baselineskip}
\begin{qitem}{49}
\qmeta{Single-correct MCQ}{QFT of the uniform state}
The QFT is applied to the uniform superposition $\frac{1}{\sqrt N}\sum_{x=0}^{N-1}\ket{x}$. The output is
\par\smallskip\noindent\begin{tabularx}{\linewidth}{@{}*{2}{L}@{}}
\optlabel{A}\,$\ket{N-1}$ & \correctlabel{B}\,$\ket{0}$ \\[3pt]
\optlabel{C}\,the uniform superposition itself & \optlabel{D}\,$\ket{N/2}$
\end{tabularx}
\end{qitem}
\begin{solbox}{B}
\textbf{\color{brand}What you need to know.} $\sum_x\omega^{xy} = N\delta_{y0}$: a constant signal has only zero frequency.

\textbf{\color{brand}Step-by-step solution.} \[ \frac{1}{\sqrt N}\sum_x \frac{1}{\sqrt N}\sum_y\omega^{xy}\ket y = \frac1N\sum_y\Big(\sum_x\omega^{xy}\Big)\ket y = \ket0, \]
since $\sum_x \omega^{xy} = N\,\delta_{y,0}$. A constant (``DC'') signal has only the zero-frequency component.
Alternatively: the uniform superposition is $\mathrm{QFT}\ket0$ (product formula with $x = 0$), so applying the QFT once more gives
$\mathrm{QFT}^2\ket0 = \ket{-0 \bmod N} = \ket0$ by Q.48.

\textbf{\color{brand}Why the other options are wrong.}
\begin{itemize}[leftmargin=2.4em, itemsep=1pt, topsep=2pt]
\item[\textbf{(A)}] $\ket{N-1}$ confuses the output with the bit complement.
\item[\textbf{(C)}] The uniform state is the QFT of $\ket0$, not its own QFT.
\item[\textbf{(D)}] $\ket{N/2}$ is the highest frequency, produced by an alternating input.
\end{itemize}
\textbf{\color{brand}Remember.} Uniform $\leftrightarrow\ket0$ under the QFT.
\end{solbox}
\qsep

\needspace{13\baselineskip}
\begin{qitem}{50}
\qmeta{Single-correct MCQ}{QFT of $\ket{10\dots0}$}
For $n$ qubits, $\mathrm{QFT}\ket{10\dots0}$ (i.e.\ $x = 2^{n-1}$) equals
\par\smallskip\noindent\begin{tabularx}{\linewidth}{@{}*{4}{L}@{}}
\correctlabel{A}\,$\ket{+}^{\otimes(n-1)}\otimes\ket{-}$ & \optlabel{B}\,$\ket{-}^{\otimes n}$ & \optlabel{C}\,$\ket{+}^{\otimes n}$ & \optlabel{D}\,$\ket{-}\otimes\ket{+}^{\otimes(n-1)}$
\end{tabularx}
\end{qitem}
\begin{solbox}{A}
\textbf{\color{brand}What you need to know.} $\omega^{2^{n-1}y} = (-1)^y$ depends only on the last bit of $y$.

\textbf{\color{brand}Step-by-step solution.} Amplitude of $y$: $\frac{1}{\sqrt N}\omega^{2^{n-1}y} = \frac{1}{\sqrt N}e^{2\pi i\,y/2} = \frac{1}{\sqrt N}(-1)^{y}$. The sign depends only on
the parity of $y$, i.e.\ on its last bit $y_n$ (qubit $n$). So the state factorises as
$\ket{+}\otimes\dots\otimes\ket{+}\otimes\ket{-}$. Product-formula check: $x_1 = 1$, others $0$, so every factor has phase
$e^{2\pi i\,0.0\dots}=1$ except the last, $e^{2\pi i\,0.x_1\dots x_n} = e^{2\pi i\,0.1} = -1$. (This is the ``highest frequency''
alternating signal $+,-,+,-,\dots$.)

\textbf{\color{brand}Why the other options are wrong.}
\begin{itemize}[leftmargin=2.4em, itemsep=1pt, topsep=2pt]
\item[\textbf{(B)}] $\ket-^{\otimes n}$ would need signs depending on every bit.
\item[\textbf{(C)}] $\ket+^{\otimes n}$ is $\QFTop\ket0$.
\item[\textbf{(D)}] $\ket-$ on qubit 1 puts the sign on the most significant bit.
\end{itemize}
\textbf{\color{brand}Remember.} The highest bit of $x$ only affects the last output qubit.
\end{solbox}
\qsep

\needspace{13\baselineskip}
\begin{qitem}{51}
\qmeta{Single-correct MCQ}{Constructing the inverse QFT}
The inverse QFT on 3 qubits is built by reversing the textbook QFT circuit and conjugating every gate. Ignoring the
SWAP, the \emph{first} gate of the inverse circuit is
\par\smallskip\noindent\begin{tabularx}{\linewidth}{@{}*{2}{L}@{}}
\correctlabel{A}\,$H$ on qubit 3 & \optlabel{B}\,controlled-$R_3^\dagger$ between qubits 1 and 3 \\[3pt]
\optlabel{C}\,controlled-$R_2^\dagger$ between qubits 1 and 2 & \optlabel{D}\,$H$ on qubit 1
\end{tabularx}
\end{qitem}
\begin{solbox}{A}
\textbf{\color{brand}What you need to know.} Inverse circuit: reverse the order, invert each gate.

\textbf{\color{brand}Step-by-step solution.} QFT $= \mathrm{SWAP}_{13}\cdot H_3\cdot \mathrm{C}R_2^{(3\to2)}\cdot H_2\cdot\mathrm{C}R_3^{(3\to1)}\cdot \mathrm{C}R_2^{(2\to1)}\cdot H_1$
(rightmost acts first). Since $(ABC)^{-1} = C^{-1}B^{-1}A^{-1}$:
\[ \mathrm{QFT}^{-1} = H_1\cdot \mathrm{C}R_2^{\dagger\,(2\to1)}\cdot\mathrm{C}R_3^{\dagger\,(3\to1)}\cdot H_2\cdot \mathrm{C}R_2^{\dagger\,(3\to2)}\cdot H_3\cdot \mathrm{SWAP}_{13}, \]
so the order of application is SWAP, then $H$ on qubit 3, then $\mathrm{C}R_2^\dagger$ (3$\to$2), $H_2$, \dots, finally $H_1$.
$H^\dagger = H$ and $R_k^\dagger = \operatorname{diag}(1, e^{-2\pi i/2^k})$.

\textbf{\color{brand}Why the other options are wrong.}
\begin{itemize}[leftmargin=2.4em, itemsep=1pt, topsep=2pt]
\item[\textbf{(B)}] $CR_3^\dagger$ (1,3) comes later, after $H_2$.
\item[\textbf{(C)}] $CR_2^\dagger$ (1,2) is the second-to-last gate.
\item[\textbf{(D)}] $H$ on qubit 1 is the \emph{last} gate of the inverse, not the first: the inverse starts where the QFT ended (qubit 3).
\end{itemize}
\textbf{\color{brand}Remember.} $(ABC)^{-1} = C^{-1}B^{-1}A^{-1}$.
\end{solbox}
\qsep

\needspace{13\baselineskip}
\begin{qitem}{52}
\qmeta{Single-correct MCQ}{Phase of a QFT amplitude}
For $n = 4$, the phase (in degrees, in $[0^\circ, 360^\circ)$) of the amplitude of $\ket{y}$, $y = 3$, in $\mathrm{QFT}\ket{x}$, $x = 5$, is
\par\smallskip\noindent\begin{tabularx}{\linewidth}{@{}*{4}{L}@{}}
\optlabel{A}\,$22.5^\circ$ & \correctlabel{B}\,$337.5^\circ$ & \optlabel{C}\,$168.75^\circ$ & \optlabel{D}\,$15^\circ$
\end{tabularx}
\end{qitem}
\begin{solbox}{B}
\textbf{\color{brand}What you need to know.} Phase $= 360^\circ\cdot\frac{xy \bmod N}{N}$.

\textbf{\color{brand}Step-by-step solution.} The amplitude is $\frac14\omega^{xy}$ with $\omega = e^{2\pi i/16}$, so the phase is
$\frac{2\pi\,xy}{16} = 2\pi\cdot\frac{15}{16}$, i.e.\ $\frac{15}{16}\times360^\circ = 337.5^\circ$ (equivalently $-22.5^\circ$; the option
$22.5^\circ$ has the sign convention reversed).

\textbf{\color{brand}Why the other options are wrong.}
\begin{itemize}[leftmargin=2.4em, itemsep=1pt, topsep=2pt]
\item[\textbf{(A)}] $22.5^\circ$ uses the opposite sign ($e^{-2\pi i}$).
\item[\textbf{(C)}] $168.75^\circ$ divides by 32 instead of 16.
\item[\textbf{(D)}] $15^\circ$ is $xy$ read as degrees.
\end{itemize}
\textbf{\color{brand}Remember.} Compute $xy \bmod N$, then scale to $360^\circ$.
\end{solbox}
\qsep

\needspace{13\baselineskip}
\begin{qitem}{53}
\qmeta{Single-correct MCQ}{QFT vs.\ NumPy FFT conventions}
NumPy's two transforms of a length-$N$ vector $a$ are
\[ \texttt{fft}(a)_y = \sum_x a_x\, e^{-2\pi i\,xy/N}, \qquad \texttt{ifft}(a)_y = \frac1N\sum_x a_x\, e^{+2\pi i\,xy/N}. \]
If $\mathrm{QFT}\sum_x\alpha_x\ket x = \sum_y\beta_y\ket y$ (convention box), then the vector $\beta$ equals
\par\smallskip\noindent\begin{tabularx}{\linewidth}{@{}*{4}{L}@{}}
\optlabel{A}\,$\texttt{fft}(\alpha)$ & \correctlabel{B}\,$\sqrt N\cdot\texttt{ifft}(\alpha)$ & \optlabel{C}\,$\texttt{ifft}(\alpha)/\sqrt N$ & \optlabel{D}\,$\texttt{fft}(\alpha)/\sqrt N$
\end{tabularx}
\end{qitem}
\begin{solbox}{B}
\textbf{\color{brand}What you need to know.} Match the exponent sign and normalisation: QFT uses $+$ and $\frac{1}{\sqrt N}$.

\textbf{\color{brand}Step-by-step solution.} $\beta_y = \frac{1}{\sqrt N}\sum_x\alpha_x e^{+2\pi i\,xy/N}$. The sign of the exponent is $+$, which matches \texttt{ifft}, and
$\frac{1}{\sqrt N} = \sqrt N\cdot\frac1N$, so $\beta = \sqrt N\,\texttt{ifft}(\alpha)$ (\texttt{numpy.fft.ifft(alpha, norm="ortho")}).
$\texttt{fft}(\alpha)/\sqrt N$ is the \emph{inverse} QFT. This is how the course notebooks verify CUDA-Q's state vectors.

\textbf{\color{brand}Why the other options are wrong.}
\begin{itemize}[leftmargin=2.4em, itemsep=1pt, topsep=2pt]
\item[\textbf{(A)}] $\texttt{fft}$ has the wrong sign and no normalisation.
\item[\textbf{(C)}] $\texttt{ifft}/\sqrt N$ has normalisation $N^{-3/2}$.
\item[\textbf{(D)}] $\texttt{fft}/\sqrt N$ has the minus sign: it is the inverse QFT.
\end{itemize}
\textbf{\color{brand}Remember.} $\beta = \sqrt N\,\texttt{ifft}(\alpha)$ = \texttt{ifft(alpha, norm=''ortho'')}.
\end{solbox}
\qsep

\needspace{13\baselineskip}
\begin{qitem}{54}
\qmeta{Single-correct MCQ}{$\mathrm{QFT}\ket{x}$ is a product state}
For an $n$-qubit computational basis state $\ket x$, the state $\mathrm{QFT}\ket x$ is
\par\smallskip\noindent\begin{tabularx}{\linewidth}{@{}*{2}{L}@{}}
\optlabel{A}\,a product state only when $x$ is even & \optlabel{B}\,always maximally entangled \\[3pt]
\optlabel{C}\,entangled unless $x = 0$ & \correctlabel{D}\,always a product (unentangled) state
\end{tabularx}
\end{qitem}
\begin{solbox}{D}
\textbf{\color{brand}What you need to know.} The product formula writes $\QFTop\ket x$ as a tensor product.

\textbf{\color{brand}Step-by-step solution.} The product formula writes $\mathrm{QFT}\ket{x}$ as a tensor product of $n$ single-qubit states
$\frac{1}{\sqrt2}(\ket0 + e^{2\pi i\,0.x_{n-l+1}\dots x_n}\ket1)$, $l = 1, \dots, n$. Hence it is never entangled,
for \emph{every} $x$. (The QFT \emph{can} entangle superposition inputs, because the controlled rotations are entangling
gates --- compare Q.25.) This product structure is also why the textbook circuit needs only $O(n^2)$ gates.

\textbf{\color{brand}Why the other options are wrong.}
\begin{itemize}[leftmargin=2.4em, itemsep=1pt, topsep=2pt]
\item[\textbf{(A)}] Product for every $x$, not only even $x$.
\item[\textbf{(B)}] Never entangled for basis inputs, so ``maximally entangled'' is false.
\item[\textbf{(C)}] Product even for $x \neq 0$.
\end{itemize}
\textbf{\color{brand}Remember.} Basis input $\Rightarrow$ product output; superposition input can entangle.
\end{solbox}
\qsep

\needspace{13\baselineskip}
\begin{qitem}{55}
\qmeta{Matching lists}{Gate count $n(n-1)/2$}
Match the number of qubits $n$ (List-I) with the number of controlled-$R_k$ gates in the textbook QFT (List-II).
\par\smallskip\noindent{\renewcommand{\arraystretch}{1.35}%
\begin{tabularx}{\linewidth}{@{}c>{\raggedright\arraybackslash}p{0.33\linewidth}@{\hspace{1.2em}}cL@{}}
\toprule
\multicolumn{2}{@{}l}{\textbf{List-I}} & \multicolumn{2}{l@{}}{\textbf{List-II}} \\ \midrule
\textbf{(P)} & $n = 4$ & \textbf{(1)} & $15$ \\
\textbf{(Q)} & $n = 6$ & \textbf{(2)} & $28$ \\
\textbf{(R)} & $n = 8$ & \textbf{(3)} & $45$ \\
\textbf{(S)} & $n = 10$ & \textbf{(4)} & $6$ \\
\bottomrule
\end{tabularx}}
\par\smallskip\noindent\begin{tabularx}{\linewidth}{@{}*{2}{L}@{}}
\optlabel{A}\,P$\to$1, Q$\to$4, R$\to$2, S$\to$3 & \optlabel{B}\,P$\to$4, Q$\to$2, R$\to$1, S$\to$3 \\[3pt]
\correctlabel{C}\,P$\to$4, Q$\to$1, R$\to$2, S$\to$3 & \optlabel{D}\,P$\to$4, Q$\to$1, R$\to$3, S$\to$2
\end{tabularx}
\end{qitem}
\begin{solbox}{C}
\textbf{\color{brand}What you need to know.} $\binom n2$ for each $n$.

\textbf{\color{brand}Step-by-step solution.} Count $= \binom{n}{2} = \frac{n(n-1)}{2}$: $n = 4 \to 6$ (4), $n = 6 \to 15$ (1), $n = 8 \to 28$ (2), $n = 10 \to 45$ (3).
Hence P$\to$4, Q$\to$1, R$\to$2, S$\to$3. Growth is quadratic, compared with the exponential $2^n$ amplitudes transformed.

\textbf{\color{brand}Why the other options are wrong.}
\begin{itemize}[leftmargin=2.4em, itemsep=1pt, topsep=2pt]
\item[\textbf{(A)}] P$\to$1 swaps $n = 4$ and $n = 6$.
\item[\textbf{(B)}] Q$\to$2 assigns $28$ to $n = 6$.
\item[\textbf{(D)}] R$\to$3 assigns $45$ to $n = 8$.
\end{itemize}
\textbf{\color{brand}Remember.} $6, 15, 28, 45$ for $n = 4, 6, 8, 10$.
\end{solbox}
\qsep

\needspace{13\baselineskip}
\begin{qitem}{56}
\qmeta{Matching lists}{Two-qubit QFT outputs}
For $n = 2$, match each input (List-I) with $\mathrm{QFT}$ of it (List-II); qubit 1 is written first.
\par\smallskip\noindent{\renewcommand{\arraystretch}{1.35}%
\begin{tabularx}{\linewidth}{@{}c>{\raggedright\arraybackslash}p{0.33\linewidth}@{\hspace{1.2em}}cL@{}}
\toprule
\multicolumn{2}{@{}l}{\textbf{List-I}} & \multicolumn{2}{l@{}}{\textbf{List-II}} \\ \midrule
\textbf{(P)} & $\ket{00}$ & \textbf{(1)} & $\frac12\big(\ket0 - \ket1\big)\big(\ket0 + i\ket1\big)$ \\
\textbf{(Q)} & $\ket{01}$ & \textbf{(2)} & $\frac12\big(\ket0 - \ket1\big)\big(\ket0 - i\ket1\big)$ \\
\textbf{(R)} & $\ket{10}$ & \textbf{(3)} & $\frac12\big(\ket0 + \ket1\big)\big(\ket0 + \ket1\big)$ \\
\textbf{(S)} & $\ket{11}$ & \textbf{(4)} & $\frac12\big(\ket0 + \ket1\big)\big(\ket0 - \ket1\big)$ \\
\bottomrule
\end{tabularx}}
\par\smallskip\noindent\begin{tabularx}{\linewidth}{@{}*{2}{L}@{}}
\optlabel{A}\,P$\to$3, Q$\to$4, R$\to$1, S$\to$2 & \optlabel{B}\,P$\to$4, Q$\to$1, R$\to$3, S$\to$2 \\[3pt]
\optlabel{C}\,P$\to$3, Q$\to$2, R$\to$4, S$\to$1 & \correctlabel{D}\,P$\to$3, Q$\to$1, R$\to$4, S$\to$2
\end{tabularx}
\end{qitem}
\begin{solbox}{D}
\textbf{\color{brand}What you need to know.} Product formula for $n = 2$: $(\ket0 + e^{2\pi i0.x_2}\ket1)(\ket0 + e^{2\pi i0.x_1x_2}\ket1)/2$.

\textbf{\color{brand}Step-by-step solution.} Product formula for $n = 2$: $\frac12\big(\ket0 + e^{2\pi i\,0.x_2}\ket1\big)\big(\ket0 + e^{2\pi i\,0.x_1x_2}\ket1\big)$.
\begin{itemize}\itemsep0pt
\item $00$: phases $e^0, e^0$ $\Rightarrow$ $(\ket0 + \ket1)(\ket0 + \ket1)/2$ (3).
\item $01$: $0.x_2 = \frac12 \to -1$; $0.x_1x_2 = 0.01_2 = \frac14 \to i$ $\Rightarrow$ (1).
\item $10$: $0.x_2 = 0 \to 1$; $0.10_2 = \frac12 \to -1$ $\Rightarrow$ (4).
\item $11$: $0.x_2 = \frac12 \to -1$; $0.11_2 = \frac34 \to -i$ $\Rightarrow$ (2).
\end{itemize}
Hence P$\to$3, Q$\to$1, R$\to$4, S$\to$2. Check $\ket{11}$ against Q.41: expanding (2) gives $\frac12(\ket{00} - i\ket{01} - \ket{10} + i\ket{11})$,
whose $\ket{10}$ amplitude is $-\frac12$.

\textbf{\color{brand}Why the other options are wrong.}
\begin{itemize}[leftmargin=2.4em, itemsep=1pt, topsep=2pt]
\item[\textbf{(A)}] Q$\to$4, R$\to$1 swaps $\ket{01}$ and $\ket{10}$ (bit order).
\item[\textbf{(B)}] P$\to$4 says $\QFTop\ket{00}$ has a minus sign; it is uniform.
\item[\textbf{(C)}] Q$\to$2, S$\to$1 swaps $\ket{01}$ and $\ket{11}$ (sign of $i$).
\end{itemize}
\textbf{\color{brand}Remember.} Do the four rows independently with the product formula.
\end{solbox}
\qsep

\needspace{13\baselineskip}
\begin{qitem}{57}
\qmeta{Assertion--Reason}{Why measuring $\mathrm{QFT}\ket{x}$ reveals nothing}
\textbf{Assertion (A):} Measuring $\mathrm{QFT}\ket{x}$ in the computational basis reveals no information about the
basis state $x$.\\[1mm]
\textbf{Reason (R):} All $N$ amplitudes of $\mathrm{QFT}\ket{x}$ have the same magnitude $1/\sqrt N$.
\begin{itemize}[leftmargin=2.6em, labelsep=0.4em, itemsep=2pt, topsep=3pt]
\item[\correctlabel{A}] Both Assertion (A) and Reason (R) are true, and R is the correct explanation of A
\item[\optlabel{B}] Both A and R are true, but R is \emph{not} the correct explanation of A
\item[\optlabel{C}] A is true, but R is false
\item[\optlabel{D}] A is false, but R is true
\end{itemize}
\end{qitem}
\begin{solbox}{A}
\textbf{\color{brand}What you need to know.} Equal magnitudes $\Rightarrow$ uniform outcome distribution independent of $x$.

\textbf{\color{brand}Step-by-step solution.} \textbf{R is true:} $|\omega^{xy}/\sqrt N| = 1/\sqrt N$ for every $y$. \textbf{A is true} \emph{because} of R: the outcome
distribution $P(y) = 1/N$ is the same uniform distribution for every $x$, so the result is statistically independent of $x$.
The information is entirely in the relative phases, which a computational-basis measurement discards; one must first
undo the QFT (or interfere the phases in some other way) to read it. R is the correct explanation. Answer: (A).

\textbf{\color{brand}Why the other options are wrong.}
\begin{itemize}[leftmargin=2.4em, itemsep=1pt, topsep=2pt]
\item[\textbf{(B)}] needs A and R both true but R not explaining A; here R is exactly the reason.
\item[\textbf{(C)}] needs A true and R false; R is true.
\item[\textbf{(D)}] needs A false and R true; A is true.
\end{itemize}
\textbf{\color{brand}Remember.} Information hidden in phases needs interference to be read.
\end{solbox}
\qsep

\needspace{13\baselineskip}
\begin{qitem}{58}
\qmeta{Statement I / II}{What the QFT does \emph{not} speed up}
\textbf{Statement I:} The QFT on $n$ qubits can be implemented with $O(n^2)$ one- and two-qubit gates.\\[1mm]
\textbf{Statement II:} Consequently, for an arbitrary classical input vector of length $2^n$, a quantum computer can output all
$2^n$ of its Fourier coefficients using only $O(n^2)$ operations in total.
\begin{itemize}[leftmargin=2.6em, labelsep=0.4em, itemsep=2pt, topsep=3pt]
\item[\optlabel{A}] Both Statement I and Statement II are true
\item[\optlabel{B}] Both Statement I and Statement II are false
\item[\correctlabel{C}] Statement I is true, but Statement II is false
\item[\optlabel{D}] Statement I is false, but Statement II is true
\end{itemize}
\end{qitem}
\begin{solbox}{C}
\textbf{\color{brand}What you need to know.} Gate count vs input loading and output readout.

\textbf{\color{brand}Step-by-step solution.} \textbf{I is true:} $n$ Hadamards, $\frac{n(n-1)}{2}$ controlled rotations and $\lfloor n/2\rfloor$ SWAPs.
\textbf{II is false}, for two reasons: (1) \emph{input} --- loading an arbitrary $2^n$-entry vector into amplitudes
generally needs $\Omega(2^n)$ gates; (2) \emph{output} --- the coefficients are amplitudes; a measurement returns just one
index $y$ with probability $|\beta_y|^2$, and estimating all $2^n$ values needs exponentially many repetitions.
The QFT is useful as a \emph{subroutine} (phase estimation, Shor) where only a few peaks need to be read. Answer: (C).

\textbf{\color{brand}Why the other options are wrong.}
\begin{itemize}[leftmargin=2.4em, itemsep=1pt, topsep=2pt]
\item[\textbf{(A)}] needs both statements true; II is false.
\item[\textbf{(B)}] needs both statements false; I is true.
\item[\textbf{(D)}] needs I false and II true; I is true.
\end{itemize}
\textbf{\color{brand}Remember.} The QFT is a subroutine, not a fast classical FFT.
\end{solbox}
\qsep

\needspace{20\baselineskip}
\begin{passagebox}{Paragraph for Questions 59 and 60}
The textbook 3-qubit QFT circuit is applied to the input $\ket{x_1x_2x_3} = \ket{011}$ ($x = 3$). The gates, in order, are:
$H$ on qubit 1; controlled-$R_2$ (control qubit 2, target qubit 1); controlled-$R_3$ (control qubit 3, target qubit 1);
$H$ on qubit 2; controlled-$R_2$ (control qubit 3, target qubit 2); $H$ on qubit 3; SWAP of qubits 1 and 3.
Here $R_k = \operatorname{diag}(1, e^{2\pi i/2^k})$, and ``the phase of qubit $j$'' means the angle $\theta$ in its state
$\frac{1}{\sqrt2}(\ket0 + e^{i\theta}\ket1)$, taken in $[0^\circ, 360^\circ)$.
\end{passagebox}
\needspace{13\baselineskip}
\begin{qitem}{59}
\qmeta{Paragraph based}{Circuit trace: phase build-up}
Immediately after the controlled-$R_3$ gate (the third gate), the phase of qubit 1 is
\par\smallskip\noindent\begin{tabularx}{\linewidth}{@{}*{4}{L}@{}}
\optlabel{A}\,$270^\circ$ & \correctlabel{B}\,$135^\circ$ & \optlabel{C}\,$90^\circ$ & \optlabel{D}\,$45^\circ$
\end{tabularx}
\end{qitem}
\begin{solbox}{B}
\textbf{\color{brand}What you need to know.} Gate-by-gate phase build-up on qubit 1: $180^\circ x_1 + 90^\circ x_2 + 45^\circ x_3$.

\textbf{\color{brand}Step-by-step solution.} Input $x_1 = 0$, $x_2 = 1$, $x_3 = 1$.
$H$ on qubit 1: phase $2\pi\,0.x_1 = 0$, state $\frac{1}{\sqrt2}(\ket0 + \ket1)$.
C-$R_2$ (control $x_2 = 1$): adds $\frac{2\pi}{4} = 90^\circ$. C-$R_3$ (control $x_3 = 1$): adds $\frac{2\pi}{8} = 45^\circ$.
Total $0^\circ + 90^\circ + 45^\circ = 135^\circ$, i.e.\ $2\pi\cdot 0.011_2 = 2\pi\cdot\frac38$. (Since qubits 2 and 3 are basis states,
the controls are classical here and nothing becomes entangled.)

\textbf{\color{brand}Why the other options are wrong.}
\begin{itemize}[leftmargin=2.4em, itemsep=1pt, topsep=2pt]
\item[\textbf{(A)}] $270^\circ$ treats $x_1 = 1$.
\item[\textbf{(C)}] $90^\circ$ stops after $CR_2$.
\item[\textbf{(D)}] $45^\circ$ counts only $CR_3$.
\end{itemize}
\textbf{\color{brand}Remember.} Each controlled rotation adds its angle only if its control bit is 1.
\end{solbox}
\qsep

\needspace{13\baselineskip}
\begin{qitem}{60}
\qmeta{Paragraph based}{Circuit trace: final amplitude, SWAP}
After the complete circuit (including the SWAP), the amplitude of the basis state $\ket{110}$ is
\par\smallskip\noindent\begin{tabularx}{\linewidth}{@{}*{4}{L}@{}}
\optlabel{A}\,$\dfrac{i}{8}$ & \optlabel{B}\,$\dfrac{e^{i\pi/4}}{\sqrt8}$ & \optlabel{C}\,$\dfrac{-i}{\sqrt8}$ & \correctlabel{D}\,$\dfrac{i}{\sqrt8}$
\end{tabularx}
\end{qitem}
\begin{solbox}{D}
\textbf{\color{brand}What you need to know.} With the SWAP the circuit equals the QFT; amplitude $= \omega^{xy}/\sqrt N$.

\textbf{\color{brand}Step-by-step solution.} With the SWAP, the circuit implements exactly $\mathrm{QFT}\ket{3} = \frac{1}{\sqrt8}\sum_y\omega^{3y}\ket y$, $\omega = e^{i\pi/4}$.
For $\ket{110}$, $y = 6$: $\omega^{18} = \omega^{16}\omega^{2} = e^{i\pi/2} = i$, so the amplitude is $\frac{i}{\sqrt8}$.
Gate-by-gate check: before the SWAP, qubit 1 has phase $\frac38$ turn (Q.59), qubit 2 has $2\pi\,0.x_2x_3 = 2\pi\,0.11_2 = 270^\circ$,
qubit 3 has $2\pi\,0.x_3 = 180^\circ$. After swapping qubits 1 and 3, $\ket{110}$ collects phases $180^\circ + 270^\circ = 450^\circ \equiv 90^\circ$
$\Rightarrow i$. Forgetting the SWAP gives $\frac38 + \frac34 = \frac98 \equiv \frac18$ turn, i.e.\ $e^{i\pi/4}/\sqrt8$ --- a wrong option.

\textbf{\color{brand}Why the other options are wrong.}
\begin{itemize}[leftmargin=2.4em, itemsep=1pt, topsep=2pt]
\item[\textbf{(A)}] $i/8$ uses the wrong normalisation ($1/8$ instead of $1/\sqrt8$).
\item[\textbf{(B)}] $e^{i\pi/4}/\sqrt8$ is what you get without the SWAP.
\item[\textbf{(C)}] $-i/\sqrt8$ uses $e^{-2\pi i}$.
\end{itemize}
\textbf{\color{brand}Remember.} Always include the final SWAPs (or relabel the output).
\end{solbox}
\qsep

\newpage

\begin{sectionbanner}SECTION D \enspace$\cdot$\enspace Advanced QFT \hfill Q.61--80\end{sectionbanner}

\needspace{13\baselineskip}
\begin{qitem}{61}
\qmeta{Single-correct MCQ}{QFT of a periodic state ($r\mid N$)}
On $n = 6$ qubits ($N = 64$) the state $\frac{1}{\sqrt8}\sum_{j=0}^{7}\ket{3 + 8j}$ is prepared, the QFT is applied and all
qubits are measured. The possible outcomes $y$ and their probabilities are
\begin{itemize}[leftmargin=2.6em, labelsep=0.4em, itemsep=2pt, topsep=3pt]
\item[\optlabel{A}] only $y = 8$, with probability $1$
\item[\optlabel{B}] all $64$ values of $y$, each with probability $\frac{1}{64}$
\item[\optlabel{C}] $y \in \{3, 11, 19, \dots, 59\}$, each with probability $\frac18$
\item[\correctlabel{D}] $y \in \{0, 8, 16, \dots, 56\}$, each with probability $\frac18$
\end{itemize}
\end{qitem}
\begin{solbox}{D}
\textbf{\color{brand}What you need to know.} Period $r \mid N$ $\Rightarrow$ $r$ peaks at multiples of $N/r$, probability $\frac1r$ each.

\textbf{\color{brand}Step-by-step solution.} With $\omega = e^{2\pi i/64}$,
\[ \beta_y = \frac{1}{\sqrt8}\cdot\frac{1}{\sqrt{64}}\sum_{j=0}^{7}\omega^{(3+8j)y} = \frac{\omega^{3y}}{\sqrt{512}}\sum_{j=0}^{7} e^{2\pi i\,jy/8}. \]
The inner sum is $8$ if $8 \mid y$ and $0$ otherwise. Hence $\beta_y = \frac{8\,\omega^{3y}}{\sqrt{512}} = \frac{\omega^{3y}}{\sqrt8}$ for
$y = 0, 8, \dots, 56$ (8 values, probability $\frac18$ each) and $\beta_y = 0$ elsewhere. A state with period $r$ (here $8$)
dividing $N$ produces $r$ peaks at the multiples of $N/r$; the offset $3$ only changes phases (see Q.63).

\textbf{\color{brand}Why the other options are wrong.}
\begin{itemize}[leftmargin=2.4em, itemsep=1pt, topsep=2pt]
\item[\textbf{(A)}] A single peak would require a constant (period-1) input.
\item[\textbf{(B)}] A flat distribution is the QFT of a single basis state.
\item[\textbf{(C)}] Peaks at $3, 11, \dots$ copy the input positions; the offset moves into phases, not positions.
\end{itemize}
\textbf{\color{brand}Remember.} Period in $x$ $\leftrightarrow$ spacing $N/r$ in $y$.
\end{solbox}
\qsep

\needspace{13\baselineskip}
\begin{qitem}{62}
\qmeta{Single-correct MCQ}{Phases of QFT peaks}
For the state $\frac{1}{\sqrt8}\sum_{j=0}^{7}\ket{3 + 8j}$ on $6$ qubits, the phase (in $[0^\circ, 360^\circ)$) of the amplitude of
$\ket{16}$ after the QFT is
\par\smallskip\noindent\begin{tabularx}{\linewidth}{@{}*{4}{L}@{}}
\optlabel{A}\,$90^\circ$ & \correctlabel{B}\,$270^\circ$ & \optlabel{C}\,$135^\circ$ & \optlabel{D}\,$0^\circ$
\end{tabularx}
\end{qitem}
\begin{solbox}{B}
\textbf{\color{brand}What you need to know.} Peak amplitudes are $\frac{1}{\sqrt r}\omega^{x_0y}$.

\textbf{\color{brand}Step-by-step solution.} From the Q.61 calculation, $\beta_{16} = \frac{1}{\sqrt8}\,\omega^{3\cdot16} = \frac{1}{\sqrt8}\,e^{2\pi i\,48/64} = \frac{1}{\sqrt8}\,e^{2\pi i\cdot\frac34}$.
Phase $= \frac34\times360^\circ = 270^\circ$ (the amplitude is $-\frac{i}{\sqrt8}$). The offset $3$ is encoded in these phases:
$\beta_{8k} \propto e^{2\pi i\,3k/8}$, which is invisible to a single measurement but is what the shift theorem predicts.

\textbf{\color{brand}Why the other options are wrong.}
\begin{itemize}[leftmargin=2.4em, itemsep=1pt, topsep=2pt]
\item[\textbf{(A)}] $90^\circ$ uses $e^{-2\pi i}$ or $x_0 = 1$.
\item[\textbf{(C)}] $135^\circ$ uses $x_0 y/N$ with the wrong $N$.
\item[\textbf{(D)}] $0^\circ$ forgets the offset phase.
\end{itemize}
\textbf{\color{brand}Remember.} The offset $x_0$ is encoded as a phase ramp over the peaks.
\end{solbox}
\qsep

\needspace{13\baselineskip}
\begin{qitem}{63}
\qmeta{Single-correct MCQ}{Shift theorem}
Let $\mathrm{QFT}\sum_x\alpha_x\ket{x} = \sum_y\beta_y\ket{y}$ on $n$ qubits ($N = 2^n$, $\omega = e^{2\pi i/N}$). For a fixed integer $s$,
$\mathrm{QFT}\sum_x\alpha_x\ket{(x + s) \bmod N}$ equals
\par\smallskip\noindent\begin{tabularx}{\linewidth}{@{}*{2}{L}@{}}
\optlabel{A}\,$\sum_y\omega^{-sy}\beta_y\ket{y}$ & \optlabel{B}\,$\sum_y\beta_{(y+s)\bmod N}\ket{y}$ \\[3pt]
\optlabel{C}\,$\sum_y\beta_y\ket{(y + s)\bmod N}$ & \correctlabel{D}\,$\sum_y\omega^{sy}\beta_y\ket{y}$
\end{tabularx}
\end{qitem}
\begin{solbox}{D}
\textbf{\color{brand}What you need to know.} Shift theorem: $\QFTop\ket{x + s} = \sum_y\omega^{sy}(\omega^{xy}/\sqrt N)\ket y$.

\textbf{\color{brand}Step-by-step solution.} By linearity it is enough to look at one basis state:
$\mathrm{QFT}\ket{x+s} = \frac{1}{\sqrt N}\sum_y\omega^{(x+s)y}\ket y = \sum_y\omega^{sy}\Big(\frac{\omega^{xy}}{\sqrt N}\Big)\ket y$
($\omega^{N} = 1$ makes ``mod $N$'' harmless). Summing with weights $\alpha_x$ gives $\sum_y\omega^{sy}\beta_y\ket y$.
\textbf{Shift theorem:} a shift in $x$ becomes a phase ramp in $y$, and $|\beta_y|^2$ is unchanged. This is why the random,
unknown offset $x_0$ left after measuring the work register in Shor's algorithm does not matter.

\textbf{\color{brand}Why the other options are wrong.}
\begin{itemize}[leftmargin=2.4em, itemsep=1pt, topsep=2pt]
\item[\textbf{(A)}] $\omega^{-sy}$ is the shift theorem for the inverse QFT.
\item[\textbf{(B)}] $\beta_{y+s}$ shifts the output instead of multiplying by a phase.
\item[\textbf{(C)}] $\ket{y+s}$ shifts the output labels.
\end{itemize}
\textbf{\color{brand}Remember.} Shift $\leftrightarrow$ phase ramp; magnitudes unchanged.
\end{solbox}
\qsep

\needspace{13\baselineskip}
\begin{qitem}{64}
\qmeta{Single-correct MCQ}{QPE precision vs.\ counting qubits}
To obtain $\varphi$ accurate to $n$ bits with probability at least $1 - \varepsilon$, phase estimation uses
$t = n + \big\lceil \log_2\!\big(2 + \frac{1}{2\varepsilon}\big)\big\rceil$ counting qubits. For $n = 5$ and $\varepsilon = 0.1$, $t$ equals
\par\smallskip\noindent\begin{tabularx}{\linewidth}{@{}*{4}{L}@{}}
\optlabel{A}\,$7$ & \correctlabel{B}\,$8$ & \optlabel{C}\,$10$ & \optlabel{D}\,$9$
\end{tabularx}
\end{qitem}
\begin{solbox}{B}
\textbf{\color{brand}What you need to know.} $t = n + \lceil\log_2(2 + \frac{1}{2\varepsilon})\rceil$.

\textbf{\color{brand}Step-by-step solution.} $2 + \frac{1}{2\times0.1} = 2 + 5 = 7$ and $\log_2 7 \approx 2.807$, so $\lceil\log_2 7\rceil = 3$ and $t = 5 + 3 = 8$.
The extra 3 qubits make the outcome distribution (whose peak has width about one grid point $2^{-t}$, with tails
$\propto 1/\delta^2$) concentrated enough that the first 5 bits are right with probability $\ge 0.9$.

\textbf{\color{brand}Why the other options are wrong.}
\begin{itemize}[leftmargin=2.4em, itemsep=1pt, topsep=2pt]
\item[\textbf{(A)}] $7$ rounds $\log_2 7$ down.
\item[\textbf{(C)}] $10$ doubles $n$.
\item[\textbf{(D)}] $9$ uses $\varepsilon = 0.05$ or adds one extra.
\end{itemize}
\textbf{\color{brand}Remember.} A few extra qubits buy confidence; they grow like $\log(1/\varepsilon)$.
\end{solbox}
\qsep

\needspace{13\baselineskip}
\begin{qitem}{65}
\qmeta{Single-correct MCQ}{QPE with exact $t$-bit phase}
Phase estimation with $t = 3$ counting qubits (inverse QFT readout) is run on an eigenstate with $\varphi = \frac58$.
The probability that the counting register reads $101$ is
\par\smallskip\noindent\begin{tabularx}{\linewidth}{@{}*{4}{L}@{}}
\optlabel{A}\,$\dfrac58$ & \correctlabel{B}\,$1$ & \optlabel{C}\,$\dfrac18$ & \optlabel{D}\,$\dfrac{4}{\pi^2}$
\end{tabularx}
\end{qitem}
\begin{solbox}{B}
\textbf{\color{brand}What you need to know.} If $\varphi = b/2^t$ exactly, the register is $\QFTop\ket b$ and the inverse QFT returns $b$ with certainty.

\textbf{\color{brand}Step-by-step solution.} After the controlled-$U^{2^j}$ gates the counting register holds
$\frac{1}{\sqrt8}\sum_{k=0}^{7}e^{2\pi i\varphi k}\ket{k} = \frac{1}{\sqrt8}\sum_k\omega^{5k}\ket k = \mathrm{QFT}\ket{5}$ ($\omega = e^{2\pi i/8}$),
because $\varphi = \frac58 = 0.101_2$ is exactly representable with 3 bits. The inverse QFT returns $\ket{5} = \ket{101}$ with
certainty. The bound $\frac{4}{\pi^2}$ applies only when $\varphi$ is \emph{not} a multiple of $2^{-t}$.

\textbf{\color{brand}Why the other options are wrong.}
\begin{itemize}[leftmargin=2.4em, itemsep=1pt, topsep=2pt]
\item[\textbf{(A)}] $\frac58$ confuses the phase value with a probability.
\item[\textbf{(C)}] $\frac18$ assumes a uniform output.
\item[\textbf{(D)}] $\frac{4}{\pi^2}$ is the bound for \emph{non}-exact phases.
\end{itemize}
\textbf{\color{brand}Remember.} Exact binary phases are read perfectly.
\end{solbox}
\qsep

\needspace{13\baselineskip}
\begin{qitem}{66}
\qmeta{Single-correct MCQ}{The $4/\pi^2$ success bound}
If $\varphi$ is not exactly a $t$-bit binary fraction, phase estimation returns the \emph{single} closest $t$-bit
approximation with probability at least
\par\smallskip\noindent\begin{tabularx}{\linewidth}{@{}*{4}{L}@{}}
\optlabel{A}\,$\dfrac12$ & \optlabel{B}\,$\dfrac1\pi \approx 0.318$ & \correctlabel{C}\,$\dfrac{4}{\pi^2} \approx 0.405$ & \optlabel{D}\,$\dfrac{8}{\pi^2} \approx 0.81$
\end{tabularx}
\end{qitem}
\begin{solbox}{C}
\textbf{\color{brand}What you need to know.} $\sin x \ge 2x/\pi$ on $[0,\pi/2]$ gives $|\alpha_b| \ge 2/\pi$.

\textbf{\color{brand}Step-by-step solution.} Let $b$ be the closest grid point and $\delta = \varphi - b/2^t$, with $0 < |\delta| \le 2^{-(t+1)}$. The amplitude of $b$ is
\[ \alpha_b = \frac{1}{2^t}\sum_{k=0}^{2^t-1}e^{2\pi i k\delta}, \qquad |\alpha_b| = \frac{|\sin(\pi 2^t\delta)|}{2^t|\sin(\pi\delta)|}. \]
Since $\pi2^t|\delta| \le \frac\pi2$, use $\sin x \ge \frac{2x}{\pi}$ there (concavity) in the numerator and $\sin x \le x$ in the
denominator: $|\alpha_b| \ge \frac{2\cdot2^t|\delta|}{2^t\pi|\delta|} = \frac{2}{\pi}$, so $P(b) \ge \frac{4}{\pi^2} \approx 0.405$.
($\frac{8}{\pi^2}$ is the bound for getting \emph{one of the two} nearest grid points.)

\textbf{\color{brand}Why the other options are wrong.}
\begin{itemize}[leftmargin=2.4em, itemsep=1pt, topsep=2pt]
\item[\textbf{(A)}] $\frac12$ is not the proven bound (it can be lower: $0.405$).
\item[\textbf{(B)}] $\frac1\pi$ is not a bound from the derivation.
\item[\textbf{(D)}] $\frac{8}{\pi^2}$ is for one of the \emph{two} closest points.
\end{itemize}
\textbf{\color{brand}Remember.} Worst case: $\varphi$ exactly midway between grid points.
\end{solbox}
\qsep

\needspace{13\baselineskip}
\begin{qitem}{67}
\qmeta{Single-correct MCQ}{QPE distribution for $\varphi = 1/3$}
Phase estimation with $t = 2$ counting qubits is run with $\varphi = \frac13$. The probability of the outcome $y = 1$
(i.e.\ the estimate $\varphi \approx \frac14$) is
\par\smallskip\noindent\begin{tabularx}{\linewidth}{@{}*{4}{L}@{}}
\optlabel{A}\,$\dfrac{6 - 3\sqrt3}{16}$ & \optlabel{B}\,$\dfrac13$ & \optlabel{C}\,$\dfrac{4}{\pi^2}$ & \correctlabel{D}\,$\dfrac{6 + 3\sqrt3}{16}$
\end{tabularx}
\end{qitem}
\begin{solbox}{D}
\textbf{\color{brand}What you need to know.} $P(y) = \frac{\sin^2(\pi2^t\delta)}{2^{2t}\sin^2(\pi\delta)}$, $\delta = \varphi - y/2^t$.

\textbf{\color{brand}Step-by-step solution.} $P(y) = \frac{1}{16}\left|\sum_{k=0}^{3}e^{2\pi i k\delta}\right|^2 = \frac{1}{16}\,\frac{\sin^2(4\pi\delta)}{\sin^2(\pi\delta)}$ with
$\delta = \varphi - \frac{y}{4}$. For $y = 1$: $\delta = \frac13 - \frac14 = \frac{1}{12}$, so
\[ P(1) = \frac{1}{16}\,\frac{\sin^2(\pi/3)}{\sin^2(\pi/12)} = \frac{1}{16}\cdot\frac{3/4}{(1 - \cos\frac\pi6)/2} = \frac{1}{16}\cdot\frac{3}{2 - \sqrt3}
   = \frac{3(2 + \sqrt3)}{16} = \frac{6 + 3\sqrt3}{16} \approx 0.700 . \]
For completeness $P(0) = \frac{1}{16}$, $P(2) = \frac{3}{16}$, $P(3) \approx 0.050$; the four add to $1$. It exceeds the guaranteed
$\frac{4}{\pi^2} \approx 0.405$ because $\frac13$ is not at the worst position (midway between grid points).

\textbf{\color{brand}Why the other options are wrong.}
\begin{itemize}[leftmargin=2.4em, itemsep=1pt, topsep=2pt]
\item[\textbf{(A)}] $\frac{6-3\sqrt3}{16}$ rationalises with the wrong sign ($\frac{3}{2-\sqrt3} = 3(2+\sqrt3)$).
\item[\textbf{(B)}] $\frac13$ confuses the phase with a probability.
\item[\textbf{(C)}] $\frac{4}{\pi^2}$ is only a lower bound.
\end{itemize}
\textbf{\color{brand}Remember.} Use $\sin^2 a = \frac{1-\cos 2a}{2}$ for exact surds.
\end{solbox}
\qsep

\needspace{13\baselineskip}
\begin{qitem}{68}
\qmeta{Single-correct MCQ}{Qubit budget of Shor's algorithm}
For Shor's algorithm with a work register of $n = \lceil\log_2 N\rceil$ qubits and a counting register of $2n$ qubits,
the total number of qubits needed to factor $N = 21$ is
\par\smallskip\noindent\begin{tabularx}{\linewidth}{@{}*{4}{L}@{}}
\correctlabel{A}\,$15$ & \optlabel{B}\,$12$ & \optlabel{C}\,$21$ & \optlabel{D}\,$10$
\end{tabularx}
\end{qitem}
\begin{solbox}{A}
\textbf{\color{brand}What you need to know.} $n = \lceil\log_2N\rceil$, total $2n + n = 3n$.

\textbf{\color{brand}Step-by-step solution.} $\log_2 21 \approx 4.39$, so $n = 5$ (5 bits hold the residues $0, \dots, 20$). Counting register $2n = 10$ qubits, work
register $5$ qubits: total $15$. The $2n$ counting qubits make $M = 2^{2n} \ge N^2$, which guarantees that each peak
$y/M$ is within $\frac{1}{2r^2}$ of some $k/r$, so continued fractions can recover $r$. (The passage uses exactly these
$10 + 5$ qubits.)

\textbf{\color{brand}Why the other options are wrong.}
\begin{itemize}[leftmargin=2.4em, itemsep=1pt, topsep=2pt]
\item[\textbf{(B)}] $12$ uses $n = 4$ (too small: $2^4 = 16 < 21$).
\item[\textbf{(C)}] $21$ uses $N$ itself.
\item[\textbf{(D)}] $10$ is the counting register only.
\end{itemize}
\textbf{\color{brand}Remember.} $M = 2^{2n} \ge N^2$ guarantees continued fractions work.
\end{solbox}
\qsep

\needspace{13\baselineskip}
\begin{qitem}{69}
\qmeta{Single-correct MCQ}{Order finding, $a = 7$, $N = 15$}
Order finding for $a = 7$, $N = 15$ uses a counting register of $t = 8$ qubits ($M = 256$). The set of possible
measurement outcomes $y$ is
\par\smallskip\noindent\begin{tabularx}{\linewidth}{@{}*{4}{L}@{}}
\optlabel{A}\,$\{0, 128\}$ & \correctlabel{B}\,$\{0, 64, 128, 192\}$ & \optlabel{C}\,$\{0, 4, 8, 12\}$ & \optlabel{D}\,$\{0, 85, 171\}$
\end{tabularx}
\end{qitem}
\begin{solbox}{B}
\textbf{\color{brand}What you need to know.} $r \mid M$ $\Rightarrow$ outcomes exactly $kM/r$.

\textbf{\color{brand}Step-by-step solution.} Powers of 7 mod 15: $7, 49 \equiv 4, 28 \equiv 13, 91 \equiv 1$, so $r = 4$. After measuring (or ignoring) the work register,
the counting register is periodic with period 4: $\frac{1}{8}\sum_{j=0}^{63}\ket{x_0 + 4j}$. Since $4 \mid 256$, by Q.61 the
outcomes are exactly the multiples of $M/r = 64$: $y \in \{0, 64, 128, 192\}$, each with probability $\frac14$, with no
leakage. ($\{0,85,171\}$ would correspond to $r = 3$.)

\textbf{\color{brand}Why the other options are wrong.}
\begin{itemize}[leftmargin=2.4em, itemsep=1pt, topsep=2pt]
\item[\textbf{(A)}] $\{0, 128\}$ is for $r = 2$.
\item[\textbf{(C)}] $\{0,4,8,12\}$ uses $M = 16$.
\item[\textbf{(D)}] $\{0, 85, 171\}$ is for $r = 3$.
\end{itemize}
\textbf{\color{brand}Remember.} Find $r$ classically to predict the outcomes.
\end{solbox}
\qsep

\needspace{13\baselineskip}
\begin{qitem}{70}
\qmeta{Single-correct MCQ}{Failure mode $\gcd(k, r) > 1$}
In the setting of the previous question (order finding for $a = 7$, $N = 15$, $M = 256$), the continued-fraction step
turns $y/256$ into $k/r'$. Which outcome leads to the \emph{wrong} candidate $r' = 2$?
\par\smallskip\noindent\begin{tabularx}{\linewidth}{@{}*{4}{L}@{}}
\optlabel{A}\,$y = 64$ & \correctlabel{B}\,$y = 128$ & \optlabel{C}\,$y = 0$ & \optlabel{D}\,$y = 192$
\end{tabularx}
\end{qitem}
\begin{solbox}{B}
\textbf{\color{brand}What you need to know.} $\gcd(k, r) > 1$ reduces $k/r$, giving only a divisor of $r$.

\textbf{\color{brand}Step-by-step solution.} $\frac{64}{256} = \frac14 \Rightarrow r' = 4$ \checkmark; $\frac{192}{256} = \frac34 \Rightarrow r' = 4$ \checkmark;
$\frac{128}{256} = \frac12 \Rightarrow r' = 2$; $\frac{0}{256} = \frac01 \Rightarrow r' = 1$ (no information).
For $y = 128$, $k = 2$ shares the factor 2 with $r = 4$, so the fraction reduces and only a \emph{divisor} of $r$ appears.
The check $7^2 = 49 \equiv 4 \not\equiv 1 \pmod{15}$ rejects it; one repeats the run or takes the lcm of several candidates.
With the correct $r = 4$: $\gcd(7^2 - 1, 15) = \gcd(48, 15) = 3$, $\gcd(7^2 + 1, 15) = \gcd(50, 15) = 5$.

\textbf{\color{brand}Why the other options are wrong.}
\begin{itemize}[leftmargin=2.4em, itemsep=1pt, topsep=2pt]
\item[\textbf{(A)}] $64/256 = \frac14$ gives $r = 4$ correctly.
\item[\textbf{(C)}] $0$ gives $\frac01$, no information ($r' = 1$), not $2$.
\item[\textbf{(D)}] $192/256 = \frac34$ gives $r = 4$ correctly.
\end{itemize}
\textbf{\color{brand}Remember.} Check $a^{r'} \equiv 1$; take the lcm of candidates.
\end{solbox}
\qsep

\needspace{13\baselineskip}
\begin{qitem}{71}
\qmeta{Single-correct MCQ}{Approximate QFT gate savings}
An approximate QFT on $n = 20$ qubits keeps only the controlled-$R_k$ gates with $k \le m = 6$ and drops the rest.
The number of controlled rotations dropped is
\par\smallskip\noindent\begin{tabularx}{\linewidth}{@{}*{4}{L}@{}}
\optlabel{A}\,$114$ & \optlabel{B}\,$85$ & \optlabel{C}\,$95$ & \correctlabel{D}\,$105$
\end{tabularx}
\end{qitem}
\begin{solbox}{D}
\textbf{\color{brand}What you need to know.} $R_k$ appears $n - k + 1$ times; total $\binom n2$.

\textbf{\color{brand}Step-by-step solution.} In the textbook circuit the gate $R_k$ (for $k = 2, \dots, n$) appears $n - k + 1$ times (once for each target qubit $j$
with $j + k - 1 \le n$). Total $= \binom{20}{2} = 190$. Kept: $k = 2, \dots, 6$: $19 + 18 + 17 + 16 + 15 = 85$. Dropped:
$190 - 85 = 105$. The approximate QFT therefore needs only $O(nm)$ gates instead of $O(n^2)$, at the price of a small error
(Q.77).

\textbf{\color{brand}Why the other options are wrong.}
\begin{itemize}[leftmargin=2.4em, itemsep=1pt, topsep=2pt]
\item[\textbf{(A)}] $114$ is not a count from the circuit.
\item[\textbf{(B)}] $85$ is the number \emph{kept}.
\item[\textbf{(C)}] $95$ is not a count from the circuit.
\end{itemize}
\textbf{\color{brand}Remember.} Approximate QFT needs $O(nm)$ gates.
\end{solbox}
\qsep

\needspace{13\baselineskip}
\begin{qitem}{72}
\qmeta{Single-correct MCQ}{Periodic state with $r\nmid N$}
On $n = 6$ qubits ($N = 64$) the uniform superposition of all $x \equiv 0 \pmod 6$, $0 \le x \le 63$, is prepared and the
QFT applied. The probability of the outcome $y = 32$ is
\par\smallskip\noindent\begin{tabularx}{\linewidth}{@{}*{4}{L}@{}}
\correctlabel{A}\,$\dfrac{11}{64}$ & \optlabel{B}\,$\dfrac{1}{64}$ & \optlabel{C}\,$\dfrac16$ & \optlabel{D}\,$\dfrac{1}{11}$
\end{tabularx}
\end{qitem}
\begin{solbox}{A}
\textbf{\color{brand}What you need to know.} At $y$ with $ry/N$ an integer all phases are 1: $P = K^2/(KN) = K/N$.

\textbf{\color{brand}Step-by-step solution.} The allowed $x$ are $0, 6, 12, \dots, 60$: $K = 11$ terms ($6\times10 = 60 \le 63 < 66$). Then
\[ \beta_{32} = \frac{1}{\sqrt{11}}\cdot\frac{1}{\sqrt{64}}\sum_{j=0}^{10}e^{2\pi i\,6j\cdot32/64} = \frac{1}{\sqrt{704}}\sum_{j=0}^{10}e^{2\pi i\cdot3j}
   = \frac{11}{\sqrt{704}}, \]
so $P(32) = \frac{121}{704} = \frac{11}{64} \approx 0.172$. The period $6$ does not divide $64$, so the ideal answer $\frac16$ is
only approximate, and most peaks ($y \approx 10.67k$) are smeared over neighbouring integers; $y = 32 = 3\cdot\frac{64}{6}$
happens to fall exactly on a grid point.

\textbf{\color{brand}Why the other options are wrong.}
\begin{itemize}[leftmargin=2.4em, itemsep=1pt, topsep=2pt]
\item[\textbf{(B)}] $\frac1{64}$ assumes a flat distribution.
\item[\textbf{(C)}] $\frac16$ is the idealised value when $r \mid N$.
\item[\textbf{(D)}] $\frac1{11}$ is the probability per input term, not of the outcome.
\end{itemize}
\textbf{\color{brand}Remember.} Count terms exactly when $r \nmid N$.
\end{solbox}
\qsep

\needspace{13\baselineskip}
\begin{qitem}{73}
\qmeta{Single-correct MCQ}{Eigenvalues of the QFT}
Every eigenvalue $\lambda$ of the $n$-qubit QFT matrix ($n \ge 1$) satisfies
\begin{itemize}[leftmargin=2.6em, labelsep=0.4em, itemsep=2pt, topsep=3pt]
\item[\optlabel{A}] $\lambda^{N} = 1$, and all $N$ of the $N$-th roots of unity occur
\item[\optlabel{B}] $|\lambda| < 1$ for at least one eigenvalue
\item[\optlabel{C}] $\lambda \in \{1, -1\}$
\item[\correctlabel{D}] $\lambda \in \{1, -1, i, -i\}$
\end{itemize}
\end{qitem}
\begin{solbox}{D}
\textbf{\color{brand}What you need to know.} $\QFTop^4 = I$ $\Rightarrow$ $\lambda^4 = 1$.

\textbf{\color{brand}Step-by-step solution.} From Q.48, $\mathrm{QFT}^2\ket x = \ket{-x}$, hence $\mathrm{QFT}^4 = I$. If $F\mathbf v = \lambda\mathbf v$ then
$\mathbf v = F^4\mathbf v = \lambda^4\mathbf v$, so $\lambda^4 = 1$: $\lambda \in \{1, -1, i, -i\}$. The QFT is unitary, so $|\lambda| = 1$
always. Imaginary eigenvalues do occur (e.g.\ $i$ for $N = 4$, where $\operatorname{tr}F = 1 + i$), so $\{1, -1\}$ is too small;
for $N \ge 8$ all four values occur, but never all $N$ roots of unity once $N > 4$.

\textbf{\color{brand}Why the other options are wrong.}
\begin{itemize}[leftmargin=2.4em, itemsep=1pt, topsep=2pt]
\item[\textbf{(A)}] Not all $N$-th roots occur; only the four 4th roots.
\item[\textbf{(B)}] $|\lambda| = 1$ for every eigenvalue of a unitary.
\item[\textbf{(C)}] $\{1,-1\}$ is too small: $i$ occurs already for $N = 4$.
\end{itemize}
\textbf{\color{brand}Remember.} Spectrum of the Fourier transform: $\{\pm1, \pm i\}$.
\end{solbox}
\qsep

\needspace{13\baselineskip}
\begin{qitem}{74}
\qmeta{Single-correct MCQ}{Semiclassical QFT}
In phase estimation the inverse QFT is immediately followed by measurement of every counting qubit. Using the
semiclassical (Griffiths--Niu) QFT, the number of \emph{two-qubit} gates needed for this inverse QFT on $n$ qubits is
\par\smallskip\noindent\begin{tabularx}{\linewidth}{@{}*{4}{L}@{}}
\optlabel{A}\,$n - 1$ & \optlabel{B}\,$\dfrac{n(n-1)}{2}$ & \correctlabel{C}\,$0$ & \optlabel{D}\,$n$
\end{tabularx}
\end{qitem}
\begin{solbox}{C}
\textbf{\color{brand}What you need to know.} Measure-and-feed-forward replaces each controlled rotation by a classically controlled single-qubit rotation.

\textbf{\color{brand}Step-by-step solution.} Controlled-phase gates are symmetric (Q.27), so in each $\mathrm{C}R_k^\dagger$ the qubit that is measured \emph{first} can be
taken as the control. Once a qubit has been measured, its value $y_j$ is a classical bit, and ``controlled on a measured
qubit'' becomes ``apply $R_k^\dagger$ if $y_j = 1$'': a classically controlled single-qubit gate. Measuring the qubits one
after the other and feeding each result forward removes all $\frac{n(n-1)}{2}$ two-qubit gates. This needs only
mid-circuit measurement plus feed-forward, and even allows one counting qubit to be recycled $n$ times.

\textbf{\color{brand}Why the other options are wrong.}
\begin{itemize}[leftmargin=2.4em, itemsep=1pt, topsep=2pt]
\item[\textbf{(A)}] $n-1$ would keep one gate per qubit; none are needed.
\item[\textbf{(B)}] $\binom n2$ is the fully quantum circuit.
\item[\textbf{(D)}] $n$ confuses gates with measurements.
\end{itemize}
\textbf{\color{brand}Remember.} Symmetric controlled phases + early measurement = no two-qubit gates.
\end{solbox}
\qsep

\needspace{13\baselineskip}
\begin{qitem}{75}
\qmeta{Matching lists}{Applications of the QFT}
Match each algorithm in List-I with the role played by the QFT, or the related step, in List-II.
\par\smallskip\noindent{\renewcommand{\arraystretch}{1.35}%
\begin{tabularx}{\linewidth}{@{}c>{\raggedright\arraybackslash}p{0.33\linewidth}@{\hspace{1.2em}}cL@{}}
\toprule
\multicolumn{2}{@{}l}{\textbf{List-I}} & \multicolumn{2}{l@{}}{\textbf{List-II}} \\ \midrule
\textbf{(P)} & Quantum phase estimation & \textbf{(1)} & adds a number by applying controlled phase rotations to a register held in the Fourier basis \\
\textbf{(Q)} & Shor's factoring algorithm & \textbf{(2)} & turns the phases $e^{2\pi i\,2^j\varphi}$ in the counting register into the binary digits of $\varphi$ \\
\textbf{(R)} & Draper's quantum adder & \textbf{(3)} & after a QPE step, applies an ancilla rotation proportional to $1/\lambda$ for each eigenvalue $\lambda$ \\
\textbf{(S)} & HHL linear-system solver & \textbf{(4)} & extracts an order $r$ from peaks near $kM/r$ using continued fractions \\
\bottomrule
\end{tabularx}}
\par\smallskip\noindent\begin{tabularx}{\linewidth}{@{}*{2}{L}@{}}
\correctlabel{A}\,P$\to$2, Q$\to$4, R$\to$1, S$\to$3 & \optlabel{B}\,P$\to$4, Q$\to$2, R$\to$1, S$\to$3 \\[3pt]
\optlabel{C}\,P$\to$2, Q$\to$1, R$\to$4, S$\to$3 & \optlabel{D}\,P$\to$2, Q$\to$4, R$\to$3, S$\to$1
\end{tabularx}
\end{qitem}
\begin{solbox}{A}
\textbf{\color{brand}What you need to know.} Know the role of the QFT in each algorithm.

\textbf{\color{brand}Step-by-step solution.} QPE: controlled-$U^{2^j}$ writes $\frac{1}{\sqrt{2^t}}\sum_k e^{2\pi i\varphi k}\ket k \approx \mathrm{QFT}\ket{2^t\varphi}$; the inverse QFT reads
the digits (2). Shor: QPE applied to modular multiplication, whose eigenphases are $k/r$; continued fractions give $r$ (4).
Draper: in the Fourier basis, adding $a$ to $\ket{x}$ is just the phase ramp $\omega^{ay}$ (Q.63), done with single-qubit or
controlled rotations, then an inverse QFT (1). HHL: QPE puts eigenvalues $\lambda_j$ in a register, a controlled rotation
of an ancilla by an angle with $\sin\theta \propto 1/\lambda_j$ performs $A^{-1}$, then QPE is uncomputed (3).
Hence P$\to$2, Q$\to$4, R$\to$1, S$\to$3.

\textbf{\color{brand}Why the other options are wrong.}
\begin{itemize}[leftmargin=2.4em, itemsep=1pt, topsep=2pt]
\item[\textbf{(B)}] P$\to$4, Q$\to$2 swaps QPE and Shor.
\item[\textbf{(C)}] Q$\to$1, R$\to$4 swaps Shor and Draper.
\item[\textbf{(D)}] R$\to$3, S$\to$1 swaps Draper and HHL.
\end{itemize}
\textbf{\color{brand}Remember.} QPE reads phases; Shor = QPE + continued fractions; Draper adds in the Fourier basis; HHL inverts eigenvalues.
\end{solbox}
\qsep

\needspace{13\baselineskip}
\begin{qitem}{76}
\qmeta{Matching lists}{Period $\leftrightarrow$ peak spacing}
On $n = 4$ qubits ($N = 16$) the QFT is applied to each input in List-I and the result measured. Match with List-II.
\par\smallskip\noindent{\renewcommand{\arraystretch}{1.35}%
\begin{tabularx}{\linewidth}{@{}c>{\raggedright\arraybackslash}p{0.33\linewidth}@{\hspace{1.2em}}cL@{}}
\toprule
\multicolumn{2}{@{}l}{\textbf{List-I}} & \multicolumn{2}{l@{}}{\textbf{List-II}} \\ \midrule
\textbf{(P)} & $\ket{0}$ & \textbf{(1)} & only $y = 0$ and $y = 8$, each with probability $\frac12$ \\
\textbf{(Q)} & $\frac12\big(\ket0 + \ket4 + \ket8 + \ket{12}\big)$ & \textbf{(2)} & all 16 outcomes, each with probability $\frac{1}{16}$ \\
\textbf{(R)} & $\frac{1}{\sqrt8}\sum_{x\ \mathrm{even}}\ket{x}$ & \textbf{(3)} & only $y = 0, 4, 8, 12$, each with probability $\frac14$ \\
\textbf{(S)} & $\frac14\sum_{x=0}^{15}\ket{x}$ & \textbf{(4)} & $y = 0$ with certainty \\
\bottomrule
\end{tabularx}}
\par\smallskip\noindent\begin{tabularx}{\linewidth}{@{}*{2}{L}@{}}
\optlabel{A}\,P$\to$4, Q$\to$3, R$\to$1, S$\to$2 & \optlabel{B}\,P$\to$2, Q$\to$3, R$\to$4, S$\to$1 \\[3pt]
\correctlabel{C}\,P$\to$2, Q$\to$3, R$\to$1, S$\to$4 & \optlabel{D}\,P$\to$2, Q$\to$1, R$\to$3, S$\to$4
\end{tabularx}
\end{qitem}
\begin{solbox}{C}
\textbf{\color{brand}What you need to know.} Period $r \mid 16$ $\Rightarrow$ peaks every $16/r$.

\textbf{\color{brand}Step-by-step solution.} P: a basis state transforms to equal-magnitude amplitudes (Q.42): uniform over 16 (2).
Q: period $r = 4$ dividing 16: peaks at multiples of $16/4 = 4$, i.e.\ $0, 4, 8, 12$, each $\frac14$ (3).
R: the even numbers have period $r = 2$: peaks at multiples of $16/2 = 8$, i.e.\ $0$ and $8$, each $\frac12$ (1).
S: the uniform superposition has period 1 and goes to $\ket0$ (Q.49) (4).
Pattern: period $r$ in $x$ $\Leftrightarrow$ $r$ peaks spaced $N/r$ in $y$ --- the ``uncertainty principle'' of the Fourier
transform. Hence P$\to$2, Q$\to$3, R$\to$1, S$\to$4.

\textbf{\color{brand}Why the other options are wrong.}
\begin{itemize}[leftmargin=2.4em, itemsep=1pt, topsep=2pt]
\item[\textbf{(A)}] P$\to$4 says $\QFTop\ket0 = \ket0$; it is uniform.
\item[\textbf{(B)}] R$\to$4, S$\to$1 swaps even-$x$ and uniform states.
\item[\textbf{(D)}] Q$\to$1, R$\to$3 swaps period 4 and period 2.
\end{itemize}
\textbf{\color{brand}Remember.} Larger period $\Rightarrow$ more, closer peaks.
\end{solbox}
\qsep

\needspace{13\baselineskip}
\begin{qitem}{77}
\qmeta{Assertion--Reason}{Approximate QFT error bound}
\textbf{Assertion (A):} For large $n$, dropping every controlled-$R_k$ with $k > m$, where $m$ is only slightly larger than
$\log_2 n$, changes the output of the QFT circuit only slightly.\\[1mm]
\textbf{Reason (R):} Each dropped gate differs from the identity only by a phase $2\pi/2^{k}$, so its error is at most about
$2\pi/2^{k}$; at most $n$ gates are dropped for each $k$, giving a total error of order $n\cdot 2\pi/2^{m}$, which is small once
$2^m \gg n$.
\begin{itemize}[leftmargin=2.6em, labelsep=0.4em, itemsep=2pt, topsep=3pt]
\item[\correctlabel{A}] Both Assertion (A) and Reason (R) are true, and R is the correct explanation of A
\item[\optlabel{B}] Both A and R are true, but R is \emph{not} the correct explanation of A
\item[\optlabel{C}] A is true, but R is false
\item[\optlabel{D}] A is false, but R is true
\end{itemize}
\end{qitem}
\begin{solbox}{A}
\textbf{\color{brand}What you need to know.} $\|CR_k - I\| \le 2\pi/2^k$; errors of a product add.

\textbf{\color{brand}Step-by-step solution.} \textbf{R is true:} $\|\mathrm{C}R_k - I\| = |e^{2\pi i/2^k} - 1| = 2\sin(\pi/2^k) \le 2\pi/2^k$. Errors of a product of unitaries add
at most linearly (triangle inequality), and $R_k$ appears $n - k + 1 \le n$ times, so the total error is at most
$\sum_{k > m} n\cdot\frac{2\pi}{2^k} < \frac{2\pi n}{2^{m}}$. \textbf{A is true} \emph{because} of R: choosing $m = \log_2 n + c$ makes the error
at most $2\pi/2^{c}$, independent of $n$. (Coppersmith's approximate QFT; in practice gate noise is larger than the
error from these tiny rotations, so dropping them even \emph{improves} real results.) Answer: (A).

\textbf{\color{brand}Why the other options are wrong.}
\begin{itemize}[leftmargin=2.4em, itemsep=1pt, topsep=2pt]
\item[\textbf{(B)}] needs A and R both true but R not explaining A; R is exactly the error estimate behind A.
\item[\textbf{(C)}] needs A true and R false; R is true.
\item[\textbf{(D)}] needs A false and R true; A is true.
\end{itemize}
\textbf{\color{brand}Remember.} Tiny rotations can be dropped with bounded error.
\end{solbox}
\qsep

\needspace{13\baselineskip}
\begin{qitem}{78}
\qmeta{Statement I / II}{Non-terminating binary phases}
\textbf{Statement I:} For $\varphi = \frac13$, phase estimation with any finite number $t$ of counting qubits never returns an
outcome $y$ with $y/2^t = \varphi$ exactly, with certainty or otherwise.\\[1mm]
\textbf{Statement II:} The number $\frac13$ has no terminating binary expansion ($\frac13 = 0.010101\ldots_2$).
\begin{itemize}[leftmargin=2.6em, labelsep=0.4em, itemsep=2pt, topsep=3pt]
\item[\correctlabel{A}] Both Statement I and Statement II are true
\item[\optlabel{B}] Both Statement I and Statement II are false
\item[\optlabel{C}] Statement I is true, but Statement II is false
\item[\optlabel{D}] Statement I is false, but Statement II is true
\end{itemize}
\end{qitem}
\begin{solbox}{A}
\textbf{\color{brand}What you need to know.} Outcomes are integers $y$; $y/2^t = \frac13$ would need $3 \mid 2^t$.

\textbf{\color{brand}Step-by-step solution.} \textbf{II is true:} $\frac13 = \frac14 + \frac1{16} + \frac1{64} + \dots = 0.\overline{01}_2$; a terminating expansion would mean
$\frac13 = \frac{y}{2^t}$, i.e.\ $2^t = 3y$, impossible since $3 \nmid 2^t$. \textbf{I is true:} the outcome is always an integer
$y$, so the estimate is $y/2^t \ne \frac13$ for every $t$ by the same argument. The output distribution is peaked near
$2^t/3$ (see Q.67: $t = 2$ gives $y = 1$ with probability $\approx 0.70$). Both statements are true. Answer: (A).

\textbf{\color{brand}Why the other options are wrong.}
\begin{itemize}[leftmargin=2.4em, itemsep=1pt, topsep=2pt]
\item[\textbf{(B)}] needs both statements false; both are true.
\item[\textbf{(C)}] needs I true and II false; II is true.
\item[\textbf{(D)}] needs I false and II true; I is true.
\end{itemize}
\textbf{\color{brand}Remember.} Non-dyadic phases are only approximated.
\end{solbox}
\qsep

\needspace{20\baselineskip}
\begin{passagebox}{Paragraph for Questions 79 and 80}
\textbf{Shor's algorithm for $N = 21$ with $a = 2$.} A counting register of $t = 10$ qubits ($M = 2^{10} = 1024$) and a
5-qubit work register initialised to $\ket{1}$ are used. Steps: (i) $H^{\otimes 10}$ on the counting register;
(ii) $\ket{x}\ket{1} \mapsto \ket{x}\ket{2^x \bmod 21}$; (iii) inverse QFT on the counting register; (iv) measure the counting
register, obtaining $y$. Classical post-processing replaces $y/1024$ by the closest fraction $k/r$ with $r < 21$ (continued
fractions) and keeps $r$ if $2^r \equiv 1 \pmod{21}$. The true order of $2$ modulo $21$ is $6$, since $2^6 = 64 = 3\cdot21 + 1$.
\end{passagebox}
\needspace{13\baselineskip}
\begin{qitem}{79}
\qmeta{Paragraph based}{Shor, $N = 21$: continued fractions}
Which of the following measured values $y$ leads the post-processing to the correct order $r = 6$?
\par\smallskip\noindent\begin{tabularx}{\linewidth}{@{}*{4}{L}@{}}
\optlabel{A}\,$683$ & \optlabel{B}\,$512$ & \correctlabel{C}\,$171$ & \optlabel{D}\,$341$
\end{tabularx}
\end{qitem}
\begin{solbox}{C}
\textbf{\color{brand}What you need to know.} Continued fractions of $y/M$; good outcomes have $\gcd(k, r) = 1$.

\textbf{\color{brand}Step-by-step solution.} The peaks lie near $y \approx k\cdot\frac{1024}{6} = 170.67k$: $0, 170.67, 341.33, 512, 682.67, 853.33$.
\begin{itemize}\itemsep0pt
\item $\frac{171}{1024} = 0.16699$: continued fraction $[0;\,5,\,1,\,84,\,2]$, convergents $\frac15, \frac16, \frac{85}{509}, \frac{171}{1024}$; the last with
      denominator $< 21$ is $\frac16$ $\Rightarrow$ $r = 6$, and $2^6 \equiv 1$ \checkmark.
\item $\frac{341}{1024} \approx \frac13$ ($k = 2$), $\frac{512}{1024} = \frac12$ ($k = 3$), $\frac{683}{1024} \approx \frac23$ ($k = 4$): $\gcd(k, 6) > 1$,
      so only the divisors $3, 2, 3$ appear and $2^3 = 8$, $2^2 = 4$ are $\not\equiv 1 \pmod{21}$.
\end{itemize}
Factoring: $2^{r/2} = 8$, $\gcd(8 - 1, 21) = 7$ and $\gcd(8 + 1, 21) = 3$; indeed $21 = 3\times7$. Only $k = 1, 5$ (coprime to 6)
give $r$ directly, i.e.\ about $\frac{\varphi(6)}{6} = \frac13$ of the runs.

\textbf{\color{brand}Why the other options are wrong.}
\begin{itemize}[leftmargin=2.4em, itemsep=1pt, topsep=2pt]
\item[\textbf{(A)}] $683 \to \frac23$ ($k = 4$): divisor 3.
\item[\textbf{(B)}] $512 \to \frac12$ ($k = 3$): divisor 2.
\item[\textbf{(D)}] $341 \to \frac13$ ($k = 2$): divisor 3.
\end{itemize}
\textbf{\color{brand}Remember.} Only $\varphi(r)/r$ of the peaks give $r$ directly.
\end{solbox}
\qsep

\needspace{13\baselineskip}
\begin{qitem}{80}
\qmeta{Paragraph based}{Shor, $N = 21$: peak probability}
The probability that the measured value is exactly $y = 512$ is closest to
\par\smallskip\noindent\begin{tabularx}{\linewidth}{@{}*{4}{L}@{}}
\optlabel{A}\,$0.333$ & \optlabel{B}\,$0.500$ & \correctlabel{C}\,$0.167$ & \optlabel{D}\,$0.001$
\end{tabularx}
\end{qitem}
\begin{solbox}{C}
\textbf{\color{brand}What you need to know.} Group by work-register value; count terms per branch; on an exact grid point all phases are 1.

\textbf{\color{brand}Step-by-step solution.} After step (ii) the state is $\frac{1}{\sqrt{1024}}\sum_{x=0}^{1023}\ket{x}\ket{2^x \bmod 21}$. Group $x = x_0 + 6j$ with $x_0 \in \{0, \dots, 5\}$;
because $1024 = 6\cdot170 + 4$, the offsets $x_0 = 0,1,2,3$ have $K = 171$ terms and $x_0 = 4,5$ have $K = 170$ terms
(check: $4\cdot171 + 2\cdot170 = 1024$). The six work values $2^{x_0} \bmod 21$ are distinct, so the branches do not interfere. After
the inverse QFT, the amplitude of $\ket{y}$ in branch $x_0$ is $\frac{1}{1024}\sum_{j}e^{-2\pi i(x_0 + 6j)y/1024}$. For $y = 512$,
$e^{-2\pi i\cdot6j\cdot512/1024} = e^{-6\pi i j} = 1$, so its modulus is $K_{x_0}/1024$ and
\[ P(512) = \frac{4\cdot171^2 + 2\cdot170^2}{1024^2} = \frac{174\,764}{1\,048\,576} \approx 0.16667 \approx \frac16 . \]
$y = 512 = 3\cdot\frac{1024}{6}$ falls exactly on a grid point, so this peak is ``perfect''; the peaks at $k = 1, 2, 4, 5$ are smeared
over neighbouring $y$ (the nearest integers $171, 341, 683, 853$ get only about $0.11$ each).

\textbf{\color{brand}Why the other options are wrong.}
\begin{itemize}[leftmargin=2.4em, itemsep=1pt, topsep=2pt]
\item[\textbf{(A)}] $0.333$ would double the peak weight.
\item[\textbf{(B)}] $0.500$ treats $r = 2$.
\item[\textbf{(D)}] $0.001 \approx 1/1024$ assumes a flat distribution.
\end{itemize}
\textbf{\color{brand}Remember.} $P(\text{exact peak}) = \sum K_{x_0}^2/M^2 \approx 1/r$.
\end{solbox}
\qsep

\end{document}